\documentclass[11pt]{exam}
%%%%%%%%%%%%%%%%%%%%%%%%%%%%%%%%%%%%%%%%%%%%%%%%%%%%%%%%%%%%%%%%%%%%%%%%%%%%%%
% Standortbestimmung Mathematik – Eintritt Gymnasium (BfKA-Grundlagen)
%
% Version ohne Lösungen:   pdflatex Standortbestimmung_Mathematik.tex
% Version mit Lösungen:    pdflatex -jobname=Standortbestimmung_Mathematik_Loesungen \
%                              "\def\mitloesungen{1}\documentclass[11pt]{exam}
%%%%%%%%%%%%%%%%%%%%%%%%%%%%%%%%%%%%%%%%%%%%%%%%%%%%%%%%%%%%%%%%%%%%%%%%%%%%%%
% Standortbestimmung Mathematik – Eintritt Gymnasium (BfKA-Grundlagen)
%
% Version ohne Lösungen:   pdflatex Standortbestimmung_Mathematik.tex
% Version mit Lösungen:    pdflatex -jobname=Standortbestimmung_Mathematik_Loesungen \
%                              "\def\mitloesungen{1}\documentclass[11pt]{exam}
%%%%%%%%%%%%%%%%%%%%%%%%%%%%%%%%%%%%%%%%%%%%%%%%%%%%%%%%%%%%%%%%%%%%%%%%%%%%%%
% Standortbestimmung Mathematik – Eintritt Gymnasium (BfKA-Grundlagen)
%
% Version ohne Lösungen:   pdflatex Standortbestimmung_Mathematik.tex
% Version mit Lösungen:    pdflatex -jobname=Standortbestimmung_Mathematik_Loesungen \
%                              "\def\mitloesungen{1}\documentclass[11pt]{exam}
%%%%%%%%%%%%%%%%%%%%%%%%%%%%%%%%%%%%%%%%%%%%%%%%%%%%%%%%%%%%%%%%%%%%%%%%%%%%%%
% Standortbestimmung Mathematik – Eintritt Gymnasium (BfKA-Grundlagen)
%
% Version ohne Lösungen:   pdflatex Standortbestimmung_Mathematik.tex
% Version mit Lösungen:    pdflatex -jobname=Standortbestimmung_Mathematik_Loesungen \
%                              "\def\mitloesungen{1}\input{Standortbestimmung_Mathematik}"
% Jeweils zweimal kompilieren (Seitenzahlen).
%
% Aufbau: 9 Stränge × 4 Niveaus
%   A Grundlagen  (Aufgaben x.1, x.2)
%   B Numerisch   (Aufgaben x.3, x.4)
%   C Symbolisch  (Aufgabe  x.5)
%   D Formal      (Aufgabe  x.6)
%%%%%%%%%%%%%%%%%%%%%%%%%%%%%%%%%%%%%%%%%%%%%%%%%%%%%%%%%%%%%%%%%%%%%%%%%%%%%%

\ifdefined\mitloesungen \printanswers \fi

\usepackage[T1]{fontenc}
\usepackage{lmodern}
\usepackage[nswissgerman]{babel}
\usepackage[hidelinks]{hyperref}
\usepackage{amssymb, amsmath, setspace, multicol, array, needspace}
\usepackage{tikz}
\usepackage[margin=2cm]{geometry}

\usepackage{color}
\colorgrids
\definecolor{GridColor}{gray}{0.7}
\setlength{\gridsize}{4mm}

\singlespacing
\parindent 0ex

% ---------------------------------------------------------------------------
% Kopfzeilen und Angaben
% ---------------------------------------------------------------------------
\newcommand{\titel}{Standortbestimmung Mathematik}
\newcommand{\schuljahr}{Schuljahr 2026/27}
\newcommand{\zeit}{90 Minuten (Teile A und B)}
\newcommand{\hilfsmittel}{keine (kein Taschenrechner)}

% ---------------------------------------------------------------------------
% Hilfsbefehle
% ---------------------------------------------------------------------------
% Aufgabe mit fester Nummer (Strang.Nummer) statt fortlaufender Zählung
\newcommand{\aufgabe}[1]{%
  \needspace{8\baselineskip}%
  \renewcommand{\questionlabel}{\textbf{#1}}%
  \question}

% Teil-Überschrift über die volle Breite
\newcommand{\teil}[2]{\uplevel{\vspace{1ex}\textbf{\large Teil #1: #2}\par\vspace{1ex}}}

% Divisionszeichen (Doppelpunkt) mit korrektem Abstand
\newcommand{\geteilt}{\mathbin{:}}
% Punkt mit Koordinaten (x | y)
\newcommand{\pkt}[2]{(#1\,|\,#2)}
% Leeres Kästchen für Lückenaufgaben
\newcommand{\kaestchen}{\fbox{\rule{0pt}{1.1em}\hspace{1.6em}}}
% Platz für die Antwort (Raster), nur in der Version ohne Lösungen
\newcommand{\raum}[1]{\ifprintanswers\else\fillwithgrid{#1}\fi}
% Ankreuzfelder: \wahl = falsche, \wahlr = richtige Option
\newcommand{\wahl}[1]{\mbox{$\Box$\,#1}}
\newcommand{\wahlr}[1]{\mbox{\ifprintanswers$\blacksquare$\,\textbf{#1}\else$\Box$\,#1\fi}}
% Mehr Zeilenabstand in Aufgabenteilen mit Brüchen
\newcommand{\bruchabstand}{\setlength{\itemsep}{2.5ex}}
% Interpretationshinweis in den Lösungen
\newcommand{\diagnose}[1]{\par\smallskip\textit{Interpretation:} #1}

% ---------------------------------------------------------------------------
% Deutsche Bezeichnungen für die exam-Klasse
% ---------------------------------------------------------------------------
\renewcommand{\solutiontitle}{\noindent\textbf{Lösung:}\enspace}
\checkboxchar{$\Box$}
\checkedchar{$\blacksquare$}
\CorrectChoiceEmphasis{\bfseries}
\pointsinmargin
\nopointsinmargin

\begin{document}

\pagestyle{headandfoot}
\firstpageheader{}{}{}
\runningheader{\titel\ifprintanswers\ – Lösungen\fi}{}{\schuljahr}
\runningheadrule
\footer{}{Seite \thepage\ von \numpages}{}

% ===========================================================================
% Titelseite
% ===========================================================================
\begin{center}
  {\LARGE\bfseries \titel\ifprintanswers\ – Lösungen\fi}\\[1ex]
  {\large Eintritt ins Gymnasium – \schuljahr}
\end{center}
\vspace{1em}

\begin{tabular}{@{}p{0.5\textwidth}l@{}}
  \textbf{Name:} \makebox[5.5cm]{\hrulefill} & \textbf{Klasse:} \makebox[3cm]{\hrulefill} \\[2ex]
  \textbf{Zeit:} \zeit                       & \textbf{Hilfsmittel:} \hilfsmittel
\end{tabular}

\vspace{1.5em}
\rule{\textwidth}{.4pt}

\subsection*{Hinweise}

\begin{itemize}
  \item Diese Standortbestimmung ist \textbf{keine Prüfung} und wird nicht benotet. Sie zeigt
        uns, wo Sie in Mathematik stehen, damit wir Sie gezielt unterstützen können.
  \item Die Aufgaben werden von Teil zu Teil schwieriger. Es ist normal, dass
        Sie nicht alle Aufgaben lösen können. Versuchen Sie trotzdem jede Aufgabe.
  \item Lösen Sie die Aufgaben \textbf{ohne Taschenrechner}.
  \item Schreiben Sie Ihre Überlegungen und Zwischenschritte auf. Ein Lösungsweg
        sagt uns oft mehr als das Ergebnis.
  \item Wenn Sie eine Aufgabe nicht lösen können, gehen Sie zur nächsten weiter.
\end{itemize}

\vfill
\newpage

% ===========================================================================
\begin{questions}
% ===========================================================================

\teil{A}{Grundlagen}

% ---------------------------------------------------------------------------
% 1 Zahlen und Stellenwert
% ---------------------------------------------------------------------------
\aufgabe{1.1}
Ordnen Sie die Zahlen der Grösse nach, beginnend mit der kleinsten:
\[
  0.5; \quad 0.45; \quad 0.405; \quad 0.54
\]
\begin{solution}
  $0.405 < 0.45 < 0.5 < 0.54$
  \diagnose{Wird $0.405$ oder $0.45$ grösser als $0.5$ eingeordnet, so gilt
  vermutlich die Vorstellung «mehr Ziffern bedeutet eine grössere Zahl».
  Der Stellenwert der Dezimalstellen ist nicht gesichert.}
\end{solution}
\raum{2.5cm}

\aufgabe{1.2}
\begin{parts}
  \part Berechnen Sie: \quad $3.7 \cdot 100$ \qquad $45 \geteilt 1000$
  \part Welche Zahl liegt genau in der Mitte zwischen $0.3$ und $0.4$?
\end{parts}
\begin{solution}
  a) $370$; \ $0.045$ \qquad b) $0.35$
  \diagnose{Ergebnisse wie $3.700$ oder $0.45$ deuten auf die auswendig gelernte
  Regel «Nullen anhängen bzw.\ wegstreichen» hin, ohne Verständnis des
  Stellenwerts. Bei b) zeigt die Antwort «gibt es nicht», dass die Dichte der
  Dezimalzahlen nicht verstanden ist.}
\end{solution}
\raum{3cm}

% ---------------------------------------------------------------------------
% 2 Operationen und Rechenreihenfolge
% ---------------------------------------------------------------------------
\aufgabe{2.1}
Ein Band von $3.6$\,m Länge wird in Stücke von $0.4$\,m Länge geschnitten.
Welche Rechnung passt zu der Frage, wie viele Stücke entstehen? Kreuzen Sie an.

\begin{oneparcheckboxes}
  \choice $3.6 \cdot 0.4$
  \CorrectChoice $3.6 \geteilt 0.4$
  \choice $0.4 \geteilt 3.6$
  \choice $3.6 - 0.4$
\end{oneparcheckboxes}
\begin{solution}
  $3.6 \geteilt 0.4 = 9$ Stücke.
  \diagnose{Wer die Division meidet, hat oft die Vorstellung, man dürfe nur durch
  ganze Zahlen teilen. Die Wahl von $0.4 \geteilt 3.6$ zeigt die Regel «immer die
  grössere durch die kleinere Zahl» in umgekehrter Form bzw.\ fehlendes Verständnis
  der Division als Aufteilen.}
\end{solution}

\aufgabe{2.2}
Entscheiden Sie \textbf{ohne zu rechnen}, ob das Ergebnis grösser oder kleiner
als $50$ ist.
\begin{parts}
  \part $50 \cdot 0.8$ \hfill \wahl{grösser als 50}\quad \wahlr{kleiner als 50}
  \part $50 \geteilt 0.8$ \hfill \wahlr{grösser als 50}\quad \wahl{kleiner als 50}
  \part $50 \geteilt 1.25$ \hfill \wahl{grösser als 50}\quad \wahlr{kleiner als 50}
\end{parts}
\begin{solution}
  a) kleiner ($40$) \quad b) grösser ($62.5$) \quad c) kleiner ($40$)
  \diagnose{Fehler bei a) und b) zeigen die Fehlvorstellung «Multiplizieren
  vergrössert, Dividieren verkleinert», die aus dem Rechnen mit natürlichen
  Zahlen stammt.}
\end{solution}

% ---------------------------------------------------------------------------
% 3 Brüche
% ---------------------------------------------------------------------------
\aufgabe{3.1}
Markieren Sie die Zahlen $\dfrac{3}{4}$ und $\dfrac{5}{4}$ auf dem Zahlenstrahl.

\begin{center}
\begin{tikzpicture}[x=6cm]
  \draw[->, thick] (-0.05,0) -- (2.15,0);
  \foreach \x in {0,1,2} {
    \draw[thick] (\x,0.15) -- (\x,-0.15) node[below] {$\x$};
  }
  \ifprintanswers
    \foreach \x/\l in {0.75/{\frac{3}{4}}, 1.25/{\frac{5}{4}}} {
      \draw[very thick] (\x,0.25) -- (\x,-0.25);
      \node[above] at (\x,0.25) {$\l$};
    }
  \fi
\end{tikzpicture}
\end{center}
\begin{solution}
  Markierungen bei $0.75$ und $1.25$ (siehe Zahlenstrahl).
  \diagnose{Wer $\frac{5}{4}$ nicht einträgt oder für unmöglich hält, versteht
  Brüche nur als «Teil eines Ganzen», nicht als Zahlen. Falsche Positionen von
  $\frac{3}{4}$ zeigen Schwierigkeiten beim Einteilen der Einheit.}
\end{solution}

\aufgabe{3.2}
\begin{parts}
  \part Welcher Bruch ist jeweils grösser? Begründen Sie kurz.
  \[
    \frac{2}{3} \text{ oder } \frac{3}{4} \qquad\qquad
    \frac{5}{8} \text{ oder } \frac{5}{6} \qquad\qquad
    \frac{7}{8} \text{ oder } \frac{8}{9}
  \]
  \part Ergänzen Sie: \quad $\dfrac{2}{3} = \dfrac{\kaestchen}{12}$
\end{parts}
\begin{solution}
  a) $\frac{3}{4}$; \ $\frac{5}{6}$; \ $\frac{8}{9}$ (z.\,B.\ weil $\frac{8}{9}$
  nur $\frac{1}{9}$ von $1$ entfernt ist, $\frac{7}{8}$ aber $\frac{1}{8}$)
  \qquad b) $8$
  \diagnose{Wird $\frac{5}{8}$ als grösser gewählt («8 ist grösser als 6»), so
  werden Zähler und Nenner getrennt als natürliche Zahlen betrachtet. Das letzte
  Paar unterscheidet auswendig gelernte Verfahren (gleichnamig machen) von
  Zahlverständnis (Abstand zu~1).}
\end{solution}
\raum{3.5cm}

% ---------------------------------------------------------------------------
% 4 Proportionalität und Prozent
% ---------------------------------------------------------------------------
\aufgabe{4.1}
3 Hefte kosten CHF 7.50. Wie viel kosten 12 Hefte? Wie viel kosten 5 Hefte?
\begin{solution}
  12 Hefte: CHF 30.00 \qquad 5 Hefte: CHF 12.50 (1 Heft kostet CHF 2.50)
  \diagnose{Gelingen die 12 Hefte, die 5 Hefte aber nicht, so wird nur
  vervielfacht, der Weg über den Einheitspreis (Dreisatz) fehlt.}
\end{solution}
\raum{3cm}

\aufgabe{4.2}
Für 4 Personen braucht ein Rezept 6 Eier. Wie viele Eier braucht man für
6 Personen?
\begin{solution}
  $9$ Eier
  \diagnose{Die Antwort $8$ («2 Personen mehr, also 2 Eier mehr») ist das
  klassische Zeichen für additives statt multiplikatives Denken.}
\end{solution}
\raum{3cm}

% ---------------------------------------------------------------------------
% 5 Terme
% ---------------------------------------------------------------------------
\aufgabe{5.1}
Aus Hölzchen werden Quadrate in einer Reihe gelegt.

\begin{center}
\begin{tikzpicture}[x=0.8cm, y=0.8cm, line width=1.6pt, line cap=round]
  \newcommand{\reihe}[2]{% #1 = x-Start, #2 = Anzahl Quadrate
    \foreach \i in {1,...,#2} {
      \draw (#1+\i-1,0) -- (#1+\i,0);
      \draw (#1+\i-1,1) -- (#1+\i,1);
      \draw (#1+\i,0) -- (#1+\i,1);
    }
    \draw (#1,0) -- (#1,1);
  }
  \reihe{0}{1}  \node at (0.5,-0.6) {\small 4 Hölzchen};
  \reihe{3}{2}  \node at (4,-0.6)   {\small 7 Hölzchen};
  \reihe{7}{3}  \node at (8.5,-0.6) {\small 10 Hölzchen};
\end{tikzpicture}
\end{center}
\begin{parts}
  \part Wie viele Hölzchen braucht man für 10 Quadrate?
  \part Wie viele Hölzchen braucht man für $n$ Quadrate?
\end{parts}
\begin{solution}
  a) $31$ \qquad b) $3n + 1$ (oder gleichwertig, z.\,B.\ $4 + 3 \cdot (n-1)$)
  \diagnose{Gelingt a), aber nicht b), so kann ein Muster an Beispielen
  fortgesetzt, aber nicht verallgemeinert werden. Antworten wie $4n$ zeigen, dass
  die gemeinsamen Hölzchen nicht beachtet werden.}
\end{solution}
\raum{3.5cm}

\aufgabe{5.2}
Welcher Ausdruck ist grösser: $n + 2$ oder $2n$? Begründen Sie.
\begin{solution}
  Das hängt von $n$ ab: Für $n = 2$ sind beide gleich, für $n < 2$ ist $n + 2$
  grösser, für $n > 2$ ist $2n$ grösser.
  \diagnose{Die Antwort «$2n$, weil verdoppeln mehr ist» oder die Prüfung mit
  nur einem Wert zeigen, dass die Variable als feste Zahl und nicht als
  Platzhalter für viele Zahlen verstanden wird.}
\end{solution}
\raum{3cm}

% ---------------------------------------------------------------------------
% 6 Gleichungen
% ---------------------------------------------------------------------------
\aufgabe{6.1}
Welche Zahl gehört ins Kästchen?
\begin{parts}
  \part $8 + 4 = \kaestchen + 5$
  \part $3 \cdot \kaestchen = 24 \geteilt 2$
\end{parts}
\begin{solution}
  a) $7$ \qquad b) $4$
  \diagnose{Die Antworten $12$ oder $17$ bei a) zeigen, dass das
  Gleichheitszeichen als Aufforderung «jetzt kommt das Ergebnis» gelesen wird
  und nicht als Beziehung zwischen zwei gleichwertigen Seiten.}
\end{solution}

\aufgabe{6.2}
Ist die Aussage wahr oder falsch? Begründen Sie, möglichst ohne die Ergebnisse
auszurechnen.
\begin{parts}
  \part $37 + 25 = 36 + 26$ \hfill \wahlr{wahr}\quad \wahl{falsch}
  \part $47 - 18 = 48 - 17$ \hfill \wahl{wahr}\quad \wahlr{falsch}
  \part $6 \cdot 15 = 3 \cdot 30$ \hfill \wahlr{wahr}\quad \wahl{falsch}
\end{parts}
\begin{solution}
  a) wahr \quad b) falsch ($29 \neq 31$) \quad c) wahr
  \diagnose{Begründungen wie «links eins weniger, rechts eins mehr» zeigen
  relationales Denken. Wer b) als wahr einstuft, überträgt die
  Ausgleichsregel der Addition unreflektiert auf die Subtraktion.}
\end{solution}
\raum{3cm}

% ---------------------------------------------------------------------------
% 7 Funktionen und Graphen
% ---------------------------------------------------------------------------
\aufgabe{7.1}
Geben Sie die Koordinaten der Punkte $A$, $B$ und $C$ an.

\begin{center}
\begin{tikzpicture}[scale=0.55]
  \draw[gray!40, thin] (-5,-5) grid (5,5);
  \draw[->, thick] (-5.4,0) -- (5.6,0) node[right] {$x$};
  \draw[->, thick] (0,-5.4) -- (0,5.6) node[above] {$y$};
  \foreach \t in {-4,-2,2,4} {
    \node[below, font=\scriptsize] at (\t,0) {$\t$};
    \node[left,  font=\scriptsize] at (0,\t) {$\t$};
  }
  \foreach \p/\x/\y/\pos in {A/3/-2/right, B/-2/3/left, C/0/4/right} {
    \fill (\x,\y) circle (3.5pt);
    \node[\pos, font=\small] at (\x,\y) {$\p$};
  }
\end{tikzpicture}
\end{center}
$A = $ \fillin[$\pkt{3}{-2}$][2.5cm] \qquad
$B = $ \fillin[$\pkt{-2}{3}$][2.5cm] \qquad
$C = $ \fillin[$\pkt{0}{4}$][2.5cm]
\begin{solution}
  $A\pkt{3}{-2}$, \ $B\pkt{-2}{3}$, \ $C\pkt{0}{4}$
  \diagnose{Vertauschte Koordinaten bei $A$ und $B$ zeigen, dass die
  Reihenfolge $(x\,|\,y)$ nicht gesichert ist. Fehler bei $C$ (z.\,B.\ $(4\,|\,0)$)
  deuten auf Unsicherheit bei Punkten auf den Achsen hin.}
\end{solution}

\aufgabe{7.2}
Die $y$-Werte in der Tabelle werden alle nach derselben Regel aus den
$x$-Werten berechnet. Ergänzen Sie die Tabelle sinngemäss.
\begin{center}
  \renewcommand{\arraystretch}{1.4}
  \begin{tabular}{|c|*{4}{>{\centering\arraybackslash}p{1.2cm}|}}
    \hline
    $x$ & $0$ & $1$ & $3$  & $6$ \\ \hline
    $y$ & $2$ & $5$ & $11$ & \ifprintanswers $\mathbf{20}$\fi \\ \hline
  \end{tabular}
\end{center}
\begin{solution}
  $y = 20$ (Zuordnung $y = 3x + 2$)
  \diagnose{Da die $x$-Werte nicht in gleichen Schritten wachsen, scheitert ein
  rein schrittweises Weiterzählen (z.\,B.\ «immer $+6$» ergibt $17$). Das zeigt,
  ob der Zusammenhang zwischen $x$ und $y$ erkannt wird.}
\end{solution}
\raum{2.5cm}

% ---------------------------------------------------------------------------
% 8 Geometrie und Grössen
% ---------------------------------------------------------------------------
\aufgabe{8.1}
Ergänzen Sie.
\begin{parts}
  \part $3.5\text{ m} = $ \fillin[$350$][2cm] cm
  \part $250\text{ g} = $ \fillin[$0.25$][2cm] kg
  \part $1\text{ m}^2 = $ \fillin[$10\,000$][2cm] cm$^2$
  \part $2\text{ h }15\text{ min} = $ \fillin[$2.25$][2cm] h
\end{parts}
\begin{solution}
  a) $350$ \quad b) $0.25$ \quad c) $10\,000$ \quad d) $2.25$
  \diagnose{Die Antwort $100$ bei c) zeigt, dass die Umrechnungszahl für
  Flächeneinheiten nicht verstanden ist. Die Antwort $2.15$ bei d) zeigt, dass
  Zeitangaben wie Dezimalzahlen behandelt werden.}
\end{solution}

\aufgabe{8.2}
Ein Rechteck ist $6$\,cm lang und $4$\,cm breit.
\begin{parts}
  \part Berechnen Sie den Umfang und den Flächeninhalt.
  \part Hat jedes Rechteck mit demselben Umfang auch denselben Flächeninhalt?
        Begründen Sie.
\end{parts}
\begin{solution}
  a) Umfang $20$\,cm, Flächeninhalt $24$\,cm$^2$
  \qquad b) Nein, z.\,B.\ hat ein Rechteck mit $7$\,cm und $3$\,cm ebenfalls den
  Umfang $20$\,cm, aber den Flächeninhalt $21$\,cm$^2$.
  \diagnose{Verwechslungen von Umfang und Fläche (auch bei den Einheiten) sowie
  die Antwort «ja» bei b) zeigen, dass die beiden Grössen nicht klar
  unterschieden werden.}
\end{solution}
\raum{3.5cm}

% ---------------------------------------------------------------------------
% 9 Sprache ↔ Mathematik
% ---------------------------------------------------------------------------
\aufgabe{9.1}
Schreiben Sie die Rechnung auf und berechnen Sie das Ergebnis.
\begin{parts}
  \part die Summe von $12$ und $8$
  \part die Differenz von $20$ und $5$
  \part das Produkt von $6$ und $7$
  \part der Quotient von $36$ und $4$
\end{parts}
\begin{solution}
  a) $12 + 8 = 20$ \quad b) $20 - 5 = 15$ \quad c) $6 \cdot 7 = 42$ \quad
  d) $36 \geteilt 4 = 9$
  \diagnose{Unbekannte Fachbegriffe zeigen sich hier direkt. Bei b) und d)
  auf die Reihenfolge in der aufgeschriebenen Rechnung achten ($5 - 20$,
  $4 \geteilt 36$).}
\end{solution}
\raum{3cm}

\aufgabe{9.2}
\begin{parts}
  \part Mia hat dreimal so viele Karten wie Leo. Leo hat 12 Karten. Wie viele
        Karten hat Mia?
  \part Nora hat 12 Karten. Das sind dreimal so viele wie Tim. Wie viele Karten
        hat Tim?
\end{parts}
\begin{solution}
  a) $36$ \qquad b) $4$
  \diagnose{Die Antwort $36$ bei b) zeigt eine Schlüsselwort-Strategie
  («mal» bedeutet multiplizieren), ohne die Beziehung zwischen den Grössen zu
  erfassen.}
\end{solution}
\raum{3cm}

% ===========================================================================
\newpage
\teil{B}{Rechnen mit Zahlen}
% ===========================================================================

% ---------------------------------------------------------------------------
% 1 Zahlen und Stellenwert
% ---------------------------------------------------------------------------
\aufgabe{1.3}
\begin{parts}
  \part Runden Sie.
  \begin{subparts}
    \subpart $3.4567$ auf zwei Dezimalstellen
    \subpart $0.0496$ auf zwei Dezimalstellen
    \subpart $28\,349$ auf Tausender
  \end{subparts}
  \part Welches Ergebnis ist plausibel? Kreuzen Sie an, ohne genau zu rechnen.
        \[
          0.48 \cdot 205 \approx \quad
          \text{\wahl{$10$}} \quad \text{\wahlr{$100$}} \quad
          \text{\wahl{$1000$}} \quad \text{\wahl{$10\,000$}}
        \]
\end{parts}
\begin{solution}
  a) $3.46$; \ $0.05$; \ $28\,000$ \qquad b) $100$ (genau: $98.4$)
  \diagnose{$3.45$ zeigt Abschneiden statt Runden. Antworten wie $0.5$ oder
  $0.04$ bei $0.0496$ zeigen, dass der Stellenwert beim Runden verloren geht.
  Fehler bei b) deuten auf fehlende Überschlagsrechnung hin.}
\end{solution}
\raum{3cm}

\aufgabe{1.4}
\begin{parts}
  \part Ordnen Sie der Grösse nach, beginnend mit der kleinsten Zahl:
        \quad $-3.5; \quad -3; \quad 0.5; \quad -4; \quad -0.5$
  \part Welche Zahl ist um $5$ grösser als $-3$?
\end{parts}
\begin{solution}
  a) $-4 < -3.5 < -3 < -0.5 < 0.5$ \qquad b) $2$
  \diagnose{Wird $-3.5$ grösser als $-3$ eingeordnet, so werden negative Zahlen
  nach ihrem Betrag geordnet. Antworten wie $-8$ oder $8$ bei b) zeigen
  Unsicherheit beim Bewegen auf dem Zahlenstrahl.}
\end{solution}
\raum{2.5cm}

% ---------------------------------------------------------------------------
% 2 Operationen und Rechenreihenfolge
% ---------------------------------------------------------------------------
\aufgabe{2.3}
Berechnen Sie.
\begin{multicols}{3}
\begin{parts}
  \part $-7 + 12$
  \part $-4 - 9$
  \part $(-6) \cdot (-3)$
  \part $(-15) \geteilt 3$
  \part $5 - (-8)$
\end{parts}
\end{multicols}
\begin{solution}
  a) $5$ \quad b) $-13$ \quad c) $18$ \quad d) $-5$ \quad e) $13$
  \diagnose{Typische Fehler sind $-5$ oder $13$ bei b) sowie $-3$ bei e). Sie
  zeigen, dass Vorzeichen und Rechenzeichen nicht sicher unterschieden werden.}
\end{solution}
\raum{2.5cm}

\aufgabe{2.4}
Berechnen Sie.
\begin{multicols}{3}
\begin{parts}
  \part $12 - 3 \cdot 4$
  \part $(12 - 3) \cdot 4$
  \part $20 \geteilt 4 \cdot 5$
  \part $2 \cdot 3^2$
  \part $18 - 6 \geteilt 2 + 1$
\end{parts}
\end{multicols}
\begin{solution}
  a) $0$ \quad b) $36$ \quad c) $25$ \quad d) $18$ \quad e) $16$
  \diagnose{Typische Fehler: c) $1$ (Multiplikation vor Division), d) $36$
  (zuerst multipliziert, dann quadriert), e) $7$ (von links nach rechts
  gerechnet) oder $14$ (Addition vor Subtraktion). Sie zeigen, welche Regel der
  Rechenreihenfolge fehlt.}
\end{solution}
\raum{3cm}

% ---------------------------------------------------------------------------
% 3 Brüche
% ---------------------------------------------------------------------------
\aufgabe{3.3}
Berechnen Sie und kürzen Sie das Ergebnis so weit wie möglich.
\begin{multicols}{3}
\begin{parts}\bruchabstand
  \part $\dfrac{2}{3} + \dfrac{1}{4}$
  \part $\dfrac{5}{6} - \dfrac{1}{2}$
  \part $\dfrac{3}{4} \cdot \dfrac{2}{9}$
  \part $\dfrac{2}{5} \geteilt \dfrac{4}{15}$
  \part $3 \cdot \dfrac{2}{5}$
\end{parts}
\end{multicols}
\begin{solution}
  a) $\frac{11}{12}$ \quad b) $\frac{1}{3}$ \quad c) $\frac{1}{6}$ \quad
  d) $\frac{3}{2}$ \quad e) $\frac{6}{5}$
  \diagnose{$\frac{3}{7}$ bei a) zeigt, dass Zähler und Nenner getrennt addiert
  werden. $\frac{6}{15}$ bei e) zeigt, dass Zähler und Nenner mit $3$
  multipliziert werden (Verwechslung mit dem Erweitern). Fehler nur bei d)
  deuten auf eine nicht verstandene Divisionsregel hin.}
\end{solution}
\raum{4cm}

\aufgabe{3.4}
Ergänzen Sie die Tabelle.
\begin{center}
  \newcommand{\lsg}[1]{\ifprintanswers$\mathbf{#1}$\fi}
  \newcommand{\zh}{\rule[-2.4ex]{0pt}{6.6ex}}% Zeilenhöhe für Brüche
  \begin{tabular}{|c|c|c|}
    \hline
    \zh\makebox[2.6cm]{\textbf{Bruch}} & \makebox[2.6cm]{\textbf{Dezimalzahl}}
      & \makebox[2.6cm]{\textbf{Prozent}} \\ \hline
    \zh$\dfrac{3}{4}$          & \lsg{0.75}           & \lsg{75\,\%} \\ \hline
    \zh\lsg{\dfrac{1}{20}}     & $0.05$               & \lsg{5\,\%} \\ \hline
    \zh\lsg{\dfrac{5}{4}}      & \lsg{1.25}           & $125\,\%$ \\ \hline
    \zh$\dfrac{1}{3}$          & \lsg{0.333\ldots}    & \lsg{33.3\,\%} \\ \hline
  \end{tabular}
\end{center}
\begin{solution}
  Siehe Tabelle. Für $\frac{1}{3}$ sind auch $0.\overline{3}$ und
  $33\frac{1}{3}\,\%$ korrekt.
  \diagnose{$0.05 = 50\,\%$ zeigt Unsicherheit beim Stellenwert. Wer bei
  $125\,\%$ keinen Bruch findet, hält Prozentangaben über $100\,\%$ oder Brüche
  grösser als $1$ für unmöglich.}
\end{solution}
\raum{2.5cm}

% ---------------------------------------------------------------------------
% 4 Proportionalität und Prozent
% ---------------------------------------------------------------------------
\aufgabe{4.3}
Ein Pullover kostet CHF 80.
\begin{parts}
  \part Berechnen Sie $15\,\%$ von CHF 80.
  \part Der Preis des Pullovers wird um $25\,\%$ reduziert. Wie viel kostet er
        jetzt?
\end{parts}
\begin{solution}
  a) CHF 12 \qquad b) CHF 60
  \diagnose{Wer bei b) nur den Rabatt (CHF 20) angibt, unterscheidet nicht
  zwischen Prozentwert und neuem Wert. Fehler bei a) zeigen, dass die
  Grundidee von «Prozent als Hundertstel» fehlt.}
\end{solution}
\raum{3cm}

\aufgabe{4.4}
Ein Preis steigt von CHF 40 auf CHF 50.
\begin{parts}
  \part Um wie viel Prozent ist der Preis gestiegen?
  \part Danach sinkt der Preis wieder von CHF 50 auf CHF 40. Um wie viel Prozent
        ist er gesunken?
\end{parts}
\begin{solution}
  a) $25\,\%$ \qquad b) $20\,\%$
  \diagnose{Die Antwort $25\,\%$ bei b) (oder $20\,\%$ bei a)) zeigt, dass
  nicht erkannt wird, dass sich der Prozentsatz auf den jeweiligen Grundwert
  bezieht. Dieser Fehler ist sehr häufig.}
\end{solution}
\raum{3cm}

% ---------------------------------------------------------------------------
% 5 Terme
% ---------------------------------------------------------------------------
\aufgabe{5.3}
Vereinfachen Sie so weit wie möglich.
\begin{multicols}{3}
\begin{parts}
  \part $3a + 5 + 2a - 1$
  \part $4x + 3y - x + 2y$
  \part $2a + 3b$
  \part $3x \cdot 4x$
  \part $a^2 + a + a^2$
\end{parts}
\end{multicols}
\begin{solution}
  a) $5a + 4$ \quad b) $3x + 5y$ \quad c) nicht vereinfachbar \quad
  d) $12x^2$ \quad e) $2a^2 + a$
  \diagnose{$5ab$ bei c) sowie $a^4$, $3a^2$ oder $2a^3$ bei e) zeigen, dass
  ungleichartige Terme zusammengefasst werden. Die Aussage «$2a + 3b$ ist keine
  Lösung» weist auf die Vorstellung hin, dass ein Term ein einzelnes Ergebnis
  haben muss.}
\end{solution}
\raum{3cm}

\aufgabe{5.4}
Ist die Aussage richtig oder falsch? Kreuzen Sie an.
\begin{parts}
  \part $2a + 3b = 5ab$     \hfill \wahl{richtig}\quad \wahlr{falsch}
  \part $x + x = x^2$       \hfill \wahl{richtig}\quad \wahlr{falsch}
  \part $3x \cdot 2 = 6x$   \hfill \wahlr{richtig}\quad \wahl{falsch}
  \part $a + a^2 = a^3$     \hfill \wahl{richtig}\quad \wahlr{falsch}
  \part $4y - y = 4$        \hfill \wahl{richtig}\quad \wahlr{falsch}
\end{parts}
\begin{solution}
  a) falsch \quad b) falsch \quad c) richtig \quad d) falsch \quad e) falsch
  \diagnose{Wiedererkennungsversion von 5.3. Wer hier richtig entscheidet, bei
  5.3 aber Fehler macht, kennt die Regeln, wendet sie aber nicht sicher an.
  Bei e) zeigt «richtig», dass $y$ als Einheit oder Abkürzung gelesen wird
  («4 Äpfel minus 1 Apfel»).}
\end{solution}

% ---------------------------------------------------------------------------
% 6 Gleichungen
% ---------------------------------------------------------------------------
\aufgabe{6.3}
Lösen Sie die Gleichungen.
\begin{multicols}{3}
\begin{parts}\bruchabstand
  \part $x + 7 = 3$
  \part $3x = 21$
  \part $4x - 5 = 11$
  \part $12 = 5x + 2$
  \part $\dfrac{x}{3} = 6$
\end{parts}
\end{multicols}
\begin{solution}
  a) $x = -4$ \quad b) $x = 7$ \quad c) $x = 4$ \quad d) $x = 2$ \quad
  e) $x = 18$
  \diagnose{Bei d) steht die Variable rechts, was manche Schülerinnen und
  Schüler verunsichert. $x = 2$ bei e) zeigt, dass Umkehroperationen
  verwechselt werden.}
\end{solution}
\raum{3.5cm}

\aufgabe{6.4}
Lösen Sie die Gleichungen.
\begin{multicols}{3}
\begin{parts}
  \part $5x + 3 = 2x + 15$
  \part $7 - 2x = 1$
  \part $3x = 0$
\end{parts}
\end{multicols}
\begin{solution}
  a) $x = 4$ \quad b) $x = 3$ \quad c) $x = 0$
  \diagnose{$x = -3$ bei b) zeigt Vorzeichenfehler beim Umformen. Bei c)
  antworten manche «keine Lösung» oder «$x$ ist nicht definiert», weil $0$
  nicht als Lösung akzeptiert wird.}
\end{solution}
\raum{3.5cm}

% ---------------------------------------------------------------------------
% 7 Funktionen und Graphen
% ---------------------------------------------------------------------------
\aufgabe{7.3}
Der Graph zeigt, wie weit eine Person von zu Hause entfernt ist.

\begin{center}
\begin{tikzpicture}[x=0.22cm, y=1.1cm]
  \draw[gray!35, thin, xstep=5, ystep=0.5] (0,0) grid (45,3.5);
  \draw[->, thick] (0,0) -- (47,0) node[right] {\small Zeit (min)};
  \draw[->, thick] (0,0) -- (0,3.9) node[above] {\small Entfernung (km)};
  \foreach \t in {0,10,20,30,40} \node[below, font=\scriptsize] at (\t,0) {$\t$};
  \foreach \s in {1,2,3} \node[left, font=\scriptsize] at (0,\s) {$\s$};
  \draw[very thick] (0,0) -- (20,3) -- (30,3) -- (40,0);
\end{tikzpicture}
\end{center}
\begin{parts}
  \part Wie weit ist die Person nach 20 Minuten von zu Hause entfernt?
  \part Was geschieht zwischen der 20.\ und der 30.\ Minute?
  \part In welchem Abschnitt bewegt sich die Person am schnellsten?
        Begründen Sie.
  \part Berechnen Sie die Geschwindigkeit im ersten Abschnitt in km/h.
\end{parts}
\begin{solution}
  a) $3$\,km \quad b) Die Person bleibt stehen (macht eine Pause). \quad
  c) Auf dem Rückweg (30.--40.\ Minute), weil der Graph dort am steilsten ist.
  \quad d) $3$\,km in $20$\,min, also $9$\,km/h.
  \diagnose{Wird bei c) der «höchste» Abschnitt gewählt, so wird die Höhe des
  Graphen statt seiner Steigung gelesen. Bei b) zeigt «sie geht geradeaus
  weiter», dass der Graph als Bild des Weges gedeutet wird.}
\end{solution}
\raum{3.5cm}

\aufgabe{7.4}
Eine Velofahrerin fährt über einen Hügel: zuerst bergauf, dann bergab.
Welcher Graph zeigt ihre Geschwindigkeit $v$ in Abhängigkeit von der Zeit $t$
am besten? Kreuzen Sie an.

\begin{center}
\newcommand{\achsen}{%
  \draw[->] (0,0) -- (2.4,0) node[below] {\scriptsize $t$};
  \draw[->] (0,0) -- (0,1.7) node[left]  {\scriptsize $v$};}
\begin{tikzpicture}[thick]
  \begin{scope}[xshift=0cm]
    \achsen \draw[very thick] (0,0.3) .. controls (0.8,1.6) and (1.4,1.6) .. (2.2,0.3);
    \node at (1.1,-0.5) {\wahl{Graph 1}};
  \end{scope}
  \begin{scope}[xshift=3.3cm]
    \achsen \draw[very thick] (0,1.0) .. controls (0.6,0.3) and (0.9,0.3) .. (1.1,0.35)
                                     .. controls (1.4,0.5) and (1.8,1.3) .. (2.2,1.5);
    \node at (1.1,-0.5) {\wahlr{Graph 2}};
  \end{scope}
  \begin{scope}[xshift=6.6cm]
    \achsen \draw[very thick] (0,1.4) -- (2.2,0.3);
    \node at (1.1,-0.5) {\wahl{Graph 3}};
  \end{scope}
  \begin{scope}[xshift=9.9cm]
    \achsen \draw[very thick] (0,0.3) -- (1.1,1.3) -- (2.2,1.3);
    \node at (1.1,-0.5) {\wahl{Graph 4}};
  \end{scope}
\end{tikzpicture}
\end{center}
\begin{solution}
  Graph 2: Bergauf nimmt die Geschwindigkeit ab, bergab nimmt sie zu.
  \diagnose{Die Wahl von Graph 1 zeigt die klassische Fehlvorstellung «Graph
  als Bild der Situation»: Der Graph sieht aus wie der Hügel, beschreibt aber
  nicht die Geschwindigkeit.}
\end{solution}

% ---------------------------------------------------------------------------
% 8 Geometrie und Grössen
% ---------------------------------------------------------------------------
\aufgabe{8.3}
Aus einem Rechteck wird an einer Ecke ein kleineres Rechteck herausgeschnitten.
Berechnen Sie den Flächeninhalt und den Umfang der Figur.

\begin{center}
\begin{tikzpicture}[scale=0.5]
  \draw[very thick, fill=gray!10] (0,0) -- (8,0) -- (8,4) -- (5,4) -- (5,6) -- (0,6) -- cycle;
  \draw[dashed, gray] (5,4) -- (5,6) (5,6) -- (8,6) -- (8,4);
  \node[below] at (4,0) {\small $8$\,cm};
  \node[left]  at (0,3) {\small $6$\,cm};
  \node[above] at (6.5,6) {\small $3$\,cm};
  \node[right] at (8,5) {\small $2$\,cm};
\end{tikzpicture}
\end{center}
\begin{solution}
  Flächeninhalt: $8 \cdot 6 - 3 \cdot 2 = 42$\,cm$^2$ \qquad
  Umfang: $28$\,cm (gleich wie beim ganzen Rechteck)
  \diagnose{Wer beim Umfang etwas abzieht (z.\,B.\ $28 - 10 = 18$\,cm),
  überträgt die Überlegung zur Fläche auf den Umfang. Das zeigt, dass Umfang
  und Fläche nicht als verschiedene Grössen verstanden werden.}
\end{solution}
\raum{3.5cm}

\aufgabe{8.4}
\begin{parts}
  \part In einem Dreieck sind zwei Winkel $47^\circ$ und $68^\circ$ gross. Wie
        gross ist der dritte Winkel?
  \part Ein Aquarium ist $50$\,cm lang, $40$\,cm breit und $30$\,cm hoch. Wie
        viele Liter Wasser fasst es?
\end{parts}
\begin{solution}
  a) $65^\circ$ \qquad b) $60\,000$\,cm$^3 = 60$\,dm$^3 = 60$\,Liter
  \diagnose{Ergebnisse wie $600$\,L oder $6$\,L bei b) zeigen, dass der
  Zusammenhang $1\,\text{dm}^3 = 1$\,L bzw.\ die Umrechnung von
  Volumeneinheiten nicht gesichert ist.}
\end{solution}
\raum{3cm}

% ---------------------------------------------------------------------------
% 9 Sprache ↔ Mathematik
% ---------------------------------------------------------------------------
\aufgabe{9.3}
Schreiben Sie als \textbf{einen} Rechenausdruck und berechnen Sie das Ergebnis.
\begin{parts}
  \part das Dreifache der Differenz von $12$ und $5$
  \part die Differenz aus dem Dreifachen von $12$ und $5$
  \part Subtrahieren Sie $4$ von $10$.
  \part die Hälfte der Summe von $9$ und $15$
\end{parts}
\begin{solution}
  a) $3 \cdot (12 - 5) = 21$ \quad b) $3 \cdot 12 - 5 = 31$ \quad
  c) $10 - 4 = 6$ \quad d) $(9 + 15) \geteilt 2 = 12$
  \diagnose{a) und b) enthalten dieselben Wörter in anderer Reihenfolge und
  ergeben verschiedene Strukturen. Fehlende Klammern bei a) oder d) zeigen,
  dass die Struktur des Satzes nicht in die Rechnung übertragen wird. $4 - 10$
  bei c) ist die typische Umkehrung der Reihenfolge.}
\end{solution}
\raum{3.5cm}

\aufgabe{9.4}
Eine Klasse mit 24 Schülerinnen und Schülern macht einen Ausflug. Die Bahnfahrt
kostet pro Person CHF 18.50, der Museumseintritt pro Person CHF 12. Die Klasse
fährt um 8:15 Uhr ab und ist um 16:40 Uhr zurück. Die Schule übernimmt
CHF 300 der Gesamtkosten. Wie viel muss jede Person selbst bezahlen?
\begin{solution}
  Kosten pro Person: $18.50 + 12 = 30.50$; Gesamtkosten:
  $24 \cdot 30.50 = 732$; nach Abzug: $732 - 300 = 432$; pro Person:
  $432 \geteilt 24 = 18$. Jede Person bezahlt CHF 18.
  \diagnose{Die Zeitangaben sind überflüssig. Wer sie verwendet, kann die
  relevanten Informationen nicht auswählen. Wird CHF 300 von den Kosten einer
  Person abgezogen, so fehlt die Unterscheidung zwischen Gesamt- und
  Einzelgrösse.}
\end{solution}
\raum{4.5cm}

% ===========================================================================
\newpage
\teil{C}{Rechnen mit Variablen}
% ===========================================================================

% ---------------------------------------------------------------------------
% 1 Zahlen und Stellenwert
% ---------------------------------------------------------------------------
\aufgabe{1.5}
\begin{parts}
  \part Schreiben Sie in wissenschaftlicher Schreibweise:
        \quad $45\,000\,000$ \qquad $0.00072$
  \part Schreiben Sie als Dezimalzahl:
        \quad $3.1 \cdot 10^5$ \qquad $6 \cdot 10^{-3}$
  \part Berechnen Sie die beiden Ausdrücke und geben Sie die Ergebnisse in
        wissenschaftlicher Schreibweise bzw.\ als Zehnerpotenz an.
  \begin{multicols}{2}
  \begin{subparts}
    \subpart $(2 \cdot 10^3) \cdot (4 \cdot 10^5)$
    \subpart $10^6 \geteilt 10^2$
  \end{subparts}
  \end{multicols}
\end{parts}
\begin{solution}
  a) $4.5 \cdot 10^7$; \ $7.2 \cdot 10^{-4}$ \quad
  b) $310\,000$; \ $0.006$ \quad
  c) $8 \cdot 10^8$; \ $10^4$
  \diagnose{Falsch gezählte Stellen zeigen fehlendes Stellenwertverständnis.
  Werden negative Exponenten als negative Zahlen gedeutet ($6 \cdot 10^{-3}
  = -6000$), fehlt die Bedeutung der Schreibweise. $10^{15}$ bzw.\ $10^3$ bei
  c) zeigen, dass Exponenten multipliziert bzw.\ dividiert statt addiert bzw.\
  subtrahiert werden.}
\end{solution}
\raum{3.5cm}

% ---------------------------------------------------------------------------
% 2 Operationen und Rechenreihenfolge
% ---------------------------------------------------------------------------
\aufgabe{2.5}
Berechnen Sie den Wert der Terme für $a = -2$ und $b = 3$.
\begin{multicols}{3}
\begin{parts}
  \part $3a + b$
  \part $a^2$
  \part $-a^2$
  \part $2b - a$
  \part $a \cdot b^2$
\end{parts}
\end{multicols}
\begin{solution}
  a) $-3$ \quad b) $4$ \quad c) $-4$ \quad d) $8$ \quad e) $-18$
  \diagnose{Eine der aussagekräftigsten Aufgaben. $-4$ bei b) zeigt, dass ohne
  Klammern eingesetzt wird ($-2^2$ statt $(-2)^2$). $4$ bei c) zeigt
  Unsicherheit mit dem Minuszeichen vor der Potenz. $36$ bei e) zeigt, dass
  $(ab)^2$ gerechnet wird.}
\end{solution}
\raum{3cm}

% ---------------------------------------------------------------------------
% 3 Brüche
% ---------------------------------------------------------------------------
\aufgabe{3.5}
Fassen Sie zusammen und vereinfachen Sie so weit wie möglich.
\begin{multicols}{2}
\begin{parts}\bruchabstand
  \part $\dfrac{a}{4} + \dfrac{b}{6}$
  \part $\dfrac{x}{3} - \dfrac{x}{5}$
  \part $\dfrac{3}{x} + \dfrac{2}{x}$
  \part $\dfrac{2}{a} \cdot \dfrac{a}{6}$
\end{parts}
\end{multicols}
\begin{solution}
  a) $\frac{3a + 2b}{12}$ \quad b) $\frac{2x}{15}$ \quad c) $\frac{5}{x}$ \quad
  d) $\frac{1}{3}$
  \diagnose{Diese Aufgabe entspricht 3.3 mit Variablen. Gelingt 3.3, aber nicht
  3.5, so liegt die Schwierigkeit beim Übergang zur Algebra, nicht bei den
  Bruchregeln. $\frac{5}{2x}$ bei c) oder $\frac{a+b}{10}$ bei a) zeigen
  getrenntes Addieren von Zählern und Nennern.}
\end{solution}
\raum{4cm}

% ---------------------------------------------------------------------------
% 4 Proportionalität und Prozent
% ---------------------------------------------------------------------------
\aufgabe{4.5}
\begin{parts}
  \part Entscheiden Sie bei jeder Tabelle, ob die Zuordnung proportional,
        umgekehrt proportional oder keines von beidem ist. Geben Sie bei
        proportionalen Zuordnungen die Gleichung $y = k \cdot x$ an.
  \begin{center}
    \renewcommand{\arraystretch}{1.3}
    \newcommand{\tab}[5]{%
      \begin{tabular}[t]{@{}c@{}}
        \small Tabelle #1 \\[0.5ex]
        \begin{tabular}{|c|c|c|c|c|}
          \hline $x$ & 1 & 2 & 4 & 5 \\ \hline
                 $y$ & #2 & #3 & #4 & #5 \\ \hline
        \end{tabular}
      \end{tabular}}
    \tab{1}{3}{6}{12}{15} \qquad \tab{2}{5}{7}{11}{13} \qquad \tab{3}{20}{10}{5}{4}
  \end{center}
  \part $4$ Personen brauchen für eine Arbeit $6$ Tage. Wie lange brauchen
        $3$ Personen bei gleichem Arbeitstempo?
\end{parts}
\begin{solution}
  a) Tabelle 1: proportional, $y = 3x$; \ Tabelle 2: keines von beidem
  ($y = 2x + 3$); \ Tabelle 3: umgekehrt proportional ($x \cdot y = 20$)
  \qquad b) $8$ Tage
  \diagnose{Tabelle 2 ist die zentrale Falle: «gleichmässig wachsend» wird mit
  proportional verwechselt. $4.5$ Tage bei b) zeigen, dass ohne Prüfung direkt
  proportional gerechnet wird.}
\end{solution}
\raum{3.5cm}

% ---------------------------------------------------------------------------
% 5 Terme
% ---------------------------------------------------------------------------
\aufgabe{5.5}
\begin{parts}
  \part Multiplizieren Sie aus und vereinfachen Sie.
  \begin{multicols}{4}
    $3 \cdot (x + 4)$ \columnbreak

    $(-2) \cdot (a - 5)$ \columnbreak

    $(x + 3)(x + 2)$ \columnbreak

    $(a + b)^2$
  \end{multicols}
  \part Klammern Sie so weit wie möglich aus.
        \qquad $6x + 9$ \qquad\qquad $x^2 - 5x$
\end{parts}
\begin{solution}
  a) $3x + 12$; \ $-2a + 10$; \ $x^2 + 5x + 6$; \ $a^2 + 2ab + b^2$ \quad
  b) $3 \cdot (2x + 3)$; \ $x(x - 5)$
  \diagnose{$3x + 4$ zeigt, dass nur der erste Summand multipliziert wird.
  $-2a - 10$ zeigt Vorzeichenfehler. $(a + b)^2 = a^2 + b^2$ ist der häufigste
  Fehler überhaupt und zeigt, dass Potenzieren über die Addition «verteilt»
  wird.}
\end{solution}
\raum{4cm}

% ---------------------------------------------------------------------------
% 6 Gleichungen
% ---------------------------------------------------------------------------
\aufgabe{6.5}
\begin{parts}
  \part Lösen Sie die Gleichungen.
  \begin{multicols}{3}
    $3 \cdot (x - 2) = 2x + 1$ \columnbreak

    $\dfrac{x}{2} + \dfrac{x}{3} = 10$ \columnbreak

    $\dfrac{x + 1}{4} = \dfrac{x - 2}{3}$
  \end{multicols}
  \part Prüfen Sie, \textbf{ohne die Gleichung zu lösen}, ob $x = 2$ eine Lösung
        der folgenden Gleichung ist.
        \[ 5x - 3 = 2 \cdot (x + 1) + 1 \]
\end{parts}
\begin{solution}
  a) $x = 7$; \ $x = 12$; \ $x = 11$ \qquad
  b) Ja: Für $x = 2$ ergeben beide Seiten $7$.
  \diagnose{Teil b) prüft, ob verstanden wird, was eine Lösung ist, unabhängig
  vom Lösungsverfahren. Wer b) nur durch Lösen der Gleichung beantworten kann,
  verwendet ein Verfahren ohne Begriffsverständnis.}
\end{solution}
\raum{4.5cm}

% ---------------------------------------------------------------------------
% 7 Funktionen und Graphen
% ---------------------------------------------------------------------------
\aufgabe{7.5}
\begin{parts}
  \part Geben Sie für die Gerade im linken Koordinatensystem die Steigung, den
        $y$-Achsenabschnitt und die Funktionsgleichung an.
  \part Zeichnen Sie den Graphen von $y = -\frac{1}{2}x + 3$ ins rechte
        Koordinatensystem.
\end{parts}

\begin{center}
\newcommand{\koord}[4]{% xmin, xmax, ymin, ymax
  \draw[gray!35, thin] (#1,#3) grid (#2,#4);
  \draw[->, thick] (#1-0.4,0) -- (#2+0.5,0) node[right] {\small $x$};
  \draw[->, thick] (0,#3-0.4) -- (0,#4+0.5) node[above] {\small $y$};
  \foreach \t in {#1,...,#2} {\ifnum\t=0\else
    \draw (\t,0.1) -- (\t,-0.1) node[below, font=\tiny] {$\t$};\fi}
  \foreach \t in {#3,...,#4} {\ifnum\t=0\else
    \draw (0.1,\t) -- (-0.1,\t) node[left, font=\tiny] {$\t$};\fi}}
\begin{tikzpicture}[scale=0.5]
  \koord{-3}{4}{-3}{7}
  \draw[very thick, domain=-2:3] plot (\x, {2*\x + 1});
  \fill (0,1) circle (4pt) (2,5) circle (4pt);
\end{tikzpicture}
\hspace{1.5cm}
\begin{tikzpicture}[scale=0.5]
  \koord{-2}{8}{-2}{6}
  \ifprintanswers
    \draw[very thick, domain=-2:8] plot (\x, {-0.5*\x + 3});
  \fi
\end{tikzpicture}
\end{center}
\begin{parts}
  \setcounter{partno}{2}
  \part Liegt der Punkt $\pkt{4}{7}$ auf der Geraden $y = 2x - 1$? Begründen Sie.
  \part Die Werte in der Tabelle gehören zu einer linearen Funktion.
        Bestimmen Sie die Funktionsgleichung.
  \quad
  \begin{tabular}{|c|c|c|c|}
    \hline $x$ & $0$ & $1$ & $2$  \\ \hline
           $y$ & $4$ & $1$ & $-2$ \\ \hline
  \end{tabular}
\end{parts}
\begin{solution}
  a) Steigung $2$, $y$-Achsenabschnitt $1$, $y = 2x + 1$ \quad
  b) Gerade durch $\pkt{0}{3}$ und $\pkt{6}{0}$ (siehe Grafik) \quad
  c) Ja, denn $2 \cdot 4 - 1 = 7$. \quad
  d) $y = -3x + 4$
  \diagnose{Die Aufgabe prüft alle Wechsel zwischen Graph, Gleichung und
  Tabelle. Typische Fehler: Steigung als $\frac{\Delta x}{\Delta y}$ berechnet
  ($\frac{1}{2}$), negative Steigung steigend gezeichnet, $y$-Achsenabschnitt
  und Nullstelle verwechselt.}
\end{solution}
\raum{3.5cm}

% ---------------------------------------------------------------------------
% 8 Geometrie und Grössen
% ---------------------------------------------------------------------------
\aufgabe{8.5}
\begin{parts}
  \part Ein rechtwinkliges Dreieck hat die Katheten $6$\,cm und $8$\,cm.
        Berechnen Sie die Länge der Hypotenuse.
  \part In einem rechtwinkligen Dreieck ist die Hypotenuse $13$\,cm und eine
        Kathete $5$\,cm lang. Berechnen Sie die Länge der anderen Kathete.
  \part Eine $10$\,m lange Leiter lehnt an einer senkrechten Wand. Ihr Fuss
        steht $6$\,m von der Wand entfernt. In welcher Höhe berührt die Leiter
        die Wand?
  \part Ist ein Dreieck mit den Seitenlängen $7$\,cm, $24$\,cm und $25$\,cm
        rechtwinklig? Begründen Sie.
\end{parts}
\begin{solution}
  a) $10$\,cm \quad b) $12$\,cm \quad c) $8$\,m \quad
  d) Ja, denn $7^2 + 24^2 = 49 + 576 = 625 = 25^2$.
  \diagnose{$\sqrt{194}$ bei b) zeigt, dass $a^2 + b^2 = c^2$ mechanisch
  angewendet wird, ohne die Hypotenuse zu identifizieren. c) verlangt, das
  rechtwinklige Dreieck in einer Situation zu erkennen. d) prüft die Umkehrung
  des Satzes von Pythagoras.}
\end{solution}
\raum{5cm}

% ---------------------------------------------------------------------------
% 9 Sprache ↔ Mathematik
% ---------------------------------------------------------------------------
\aufgabe{9.5}
\begin{parts}
  \part Schreiben Sie als Term: \quad «$5$ weniger als das Doppelte von $x$»
        \quad und \quad «um $3$ grösser als die Hälfte von $y$»
  \part An einer Schule gibt es sechsmal so viele Schülerinnen und Schüler wie
        Lehrpersonen. Schreiben Sie diese Aussage als Gleichung. Verwenden Sie
        $S$ für die Anzahl Schülerinnen und Schüler und $L$ für die Anzahl
        Lehrpersonen.
  \part Beschreiben Sie die Terme in Worten: \quad $3 \cdot (x + 4)$ \quad und \quad
        $3x + 4$
  \part Verdoppelt man eine Zahl und addiert $7$, so erhält man $31$. Stellen
        Sie eine Gleichung auf und bestimmen Sie die Zahl.
\end{parts}
\begin{solution}
  a) $2x - 5$; \ $\frac{y}{2} + 3$ \quad
  b) $S = 6L$ \quad
  c) «das Dreifache der Summe von $x$ und $4$»; «die Summe aus dem Dreifachen
  von $x$ und $4$» \quad
  d) $2x + 7 = 31$, $x = 12$
  \diagnose{b) ist die klassische «Studierende-und-Professoren»-Aufgabe: $6S = L$
  entsteht durch Wort-für-Wort-Übersetzung, ohne die Beziehung zu erfassen.
  $5 - 2x$ bei a) zeigt die Umkehrung bei «weniger als». c) ist die
  Gegenrichtung zu 9.3.}
\end{solution}
\raum{4.5cm}

% ===========================================================================
\newpage
\teil{D}{Vertiefung}
% ===========================================================================

% ---------------------------------------------------------------------------
% 1 Zahlen und Stellenwert
% ---------------------------------------------------------------------------
\aufgabe{1.6}
\begin{parts}
  \part Geben Sie für jede Zahl alle Zahlenmengen ($\mathbb{N}$, $\mathbb{Z}$,
        $\mathbb{Q}$, $\mathbb{R}$) an, zu denen sie gehört.
        \[
          -3 \qquad 0.75 \qquad \sqrt{2} \qquad \frac{12}{4} \qquad 0.333\ldots
        \]
  \part Kann das Ergebnis stimmen? Entscheiden Sie \textbf{ohne nachzurechnen}
        und begründen Sie.
        \begin{itemize}
          \item Lena rechnet: $0.6 \geteilt 0.03 = 2$.
          \item Ein Schüler berechnet das Volumen eines Schulzimmers und erhält
                $18$\,m$^3$.
        \end{itemize}
\end{parts}
\begin{solution}
  a) $-3$: $\mathbb{Z}, \mathbb{Q}, \mathbb{R}$; \
  $0.75$: $\mathbb{Q}, \mathbb{R}$; \
  $\sqrt{2}$: $\mathbb{R}$; \
  $\frac{12}{4}$: $\mathbb{N}, \mathbb{Z}, \mathbb{Q}, \mathbb{R}$; \
  $0.333\ldots$: $\mathbb{Q}, \mathbb{R}$ \\
  b) Beides ist nicht plausibel: Division durch eine Zahl kleiner als $1$
  vergrössert (richtig: $20$). Ein Schulzimmer hat etwa $150$ bis
  $200$\,m$^3$.
  \diagnose{Wird $\frac{12}{4}$ nur $\mathbb{Q}$ zugeordnet, so wird die Zahl
  nach ihrer Schreibweise statt nach ihrem Wert beurteilt. Die Zahl $0$ wurde
  bewusst vermieden, da Lehrmittel uneinheitlich sind, ob $0 \in \mathbb{N}$.}
\end{solution}
\raum{4cm}

% ---------------------------------------------------------------------------
% 2 Operationen und Rechenreihenfolge
% ---------------------------------------------------------------------------
\aufgabe{2.6}
\begin{parts}
  \part Welche Terme sind für alle Werte von $x$ gleichwertig mit $2 \cdot (x - 3)$?
        Kreuzen Sie an.

        \wahl{$2x - 3$} \quad \wahlr{$2x - 6$} \quad \wahlr{$(-2) \cdot (3 - x)$}
        \quad \wahlr{$(x - 3) \cdot 2$} \quad \wahl{$2 \cdot (3 - x)$}
  \part Gilt $(a - b)^2 = (b - a)^2$ für alle Zahlen $a$ und $b$? Begründen Sie.
\end{parts}
\begin{solution}
  a) $2x - 6$, \ $(-2) \cdot (3 - x)$, \ $(x - 3) \cdot 2$ \qquad
  b) Ja, denn $b - a = -(a - b)$, und beim Quadrieren fällt das Vorzeichen weg.
  \diagnose{Die Aufgabe prüft, ob Gleichwertigkeit an der Struktur erkannt wird.
  Wer nur einzelne Werte einsetzt, kann b) nicht allgemein begründen.}
\end{solution}
\raum{3cm}

% ---------------------------------------------------------------------------
% 3 Brüche
% ---------------------------------------------------------------------------
\aufgabe{3.6}
\begin{parts}
  \part Fassen Sie zusammen: \quad $\dfrac{1}{x} + \dfrac{1}{x + 1}$
  \part Kürzen Sie, falls möglich:
        \quad $\dfrac{3x + 6}{x + 2}$ \qquad $\dfrac{x^2 - 9}{x - 3}$
        \qquad $\dfrac{x + 5}{5}$
\end{parts}
\begin{solution}
  a) $\frac{2x + 1}{x(x + 1)}$ \quad
  b) $3$; \ $x + 3$; \ nicht kürzbar
  \diagnose{$\frac{x + 5}{5} = x + 1$ oder $= x$ zeigt «Kürzen aus Summen». Wer
  bei den ersten beiden Brüchen nicht faktorisiert, erkennt die Struktur nicht.
  Die Angabe der Definitionslücken ($x \neq -2$, $x \neq 3$) wird auf diesem
  Niveau nicht verlangt.}
\end{solution}
\raum{3.5cm}

% ---------------------------------------------------------------------------
% 4 Proportionalität und Prozent
% ---------------------------------------------------------------------------
\aufgabe{4.6}
\begin{parts}
  \part Ein Preis wird zuerst um $20\,\%$ erhöht und danach um $20\,\%$
        gesenkt. Ist er danach gleich hoch wie zu Beginn? Begründen Sie.
  \part Nach einer Erhöhung um $25\,\%$ kostet ein Artikel CHF 150. Wie viel
        kostete er vorher?
  \part Eine Bevölkerung von $P$ Personen wächst jährlich um $3\,\%$. Geben Sie
        einen Term für die Bevölkerung nach $n$ Jahren an.
\end{parts}
\begin{solution}
  a) Nein: $1.2 \cdot 0.8 = 0.96$, der Preis ist $4\,\%$ tiefer. \quad
  b) $150 \geteilt 1.25 = 120$, also CHF 120 \quad
  c) $P \cdot 1.03^n$
  \diagnose{CHF 112.50 bei b) zeigt, dass $25\,\%$ vom neuen statt vom alten
  Wert abgezogen werden. $P + 0.03n$ oder $P \cdot 0.03^n$ bei c) zeigen, dass
  der Wachstumsfaktor nicht verstanden ist.}
\end{solution}
\raum{4cm}

% ---------------------------------------------------------------------------
% 5 Terme
% ---------------------------------------------------------------------------
\aufgabe{5.6}
\begin{parts}
  \part Faktorisieren Sie so weit wie möglich.
  \begin{multicols}{4}
    $x^2 - 9$ \columnbreak

    $x^2 + 6x + 9$ \columnbreak

    $4a^2b - 6ab^2$ \columnbreak

    $ax + ay - bx - by$
  \end{multicols}
  \part Ergänzen Sie: \quad $x^2 + \kaestchen\,x + 25 = (x + \kaestchen)^2$
\end{parts}
\begin{solution}
  a) $(x - 3)(x + 3)$; \ $(x + 3)^2$; \ $2ab(2a - 3b)$; \ $(a - b)(x + y)$
  \quad b) $10$ und $5$
  \diagnose{Faktorisieren verlangt, die binomischen Formeln «rückwärts» zu
  lesen. Wer 5.5 lösen kann, hier aber scheitert, kennt die Formeln nur in eine
  Richtung.}
\end{solution}
\raum{3.5cm}

% ---------------------------------------------------------------------------
% 6 Gleichungen
% ---------------------------------------------------------------------------
\aufgabe{6.6}
\begin{parts}
  \part Lösen Sie nach der angegebenen Variable auf.
  \begin{multicols}{3}
    $A = \frac{1}{2} \cdot g \cdot h$ \ nach $h$ \columnbreak

    $v = \dfrac{s}{t}$ \ nach $t$ \columnbreak

    $F = \frac{9}{5} \cdot C + 32$ \ nach $C$
  \end{multicols}
  \part Lösen Sie nach $x$ auf ($a \neq 3$): \quad $ax = 3x + 5$
  \part Bestimmen Sie die Lösungsmenge.
        \quad $2 \cdot (x + 1) = 2x + 2$ \qquad $2 \cdot (x + 1) = 2x + 3$
\end{parts}
\begin{solution}
  a) $h = \frac{2A}{g}$; \ $t = \frac{s}{v}$; \ $C = \frac{5}{9} \cdot (F - 32)$
  \quad b) $x = \frac{5}{a - 3}$ \quad
  c) $\mathbb{L} = \mathbb{R}$; \ $\mathbb{L} = \{\,\}$
  \diagnose{Bei b) muss $x$ ausgeklammert werden ($x(a - 3) = 5$). Wer nur
  Gleichungen mit Zahlen lösen kann, kommt hier nicht weiter. Bei c) zeigt
  die Reaktion «$0 = 1$, also habe ich mich verrechnet», dass Lösungsmengen
  ohne genau eine Lösung unbekannt sind.}
\end{solution}
\raum{4.5cm}

% ---------------------------------------------------------------------------
% 7 Funktionen und Graphen
% ---------------------------------------------------------------------------
\aufgabe{7.6}
\begin{parts}
  \part Was haben alle Geraden mit der Gleichung $y = mx + 2$ gemeinsam? Was
        haben alle Geraden mit der Gleichung $y = 3x + q$ gemeinsam?
  \part Berechnen Sie den Schnittpunkt der Geraden $y = 2x + 1$ und
        $y = -x + 7$.
  \part Für welchen Wert von $m$ verläuft die Gerade $y = mx + 2$ durch den
        Punkt $\pkt{4}{0}$?
\end{parts}
\begin{solution}
  a) Alle verlaufen durch $\pkt{0}{2}$; alle sind parallel (Steigung $3$).
  \quad b) $2x + 1 = -x + 7 \Rightarrow x = 2$, also $S\pkt{2}{5}$
  \quad c) $0 = 4m + 2 \Rightarrow m = -\frac{1}{2}$
  \diagnose{Die Aufgabe prüft, ob ein Buchstabe als Parameter (fest, aber
  beliebig) und nicht als Variable verstanden wird, und ob ein Schnittpunkt als
  Lösung einer Gleichung erkannt wird.}
\end{solution}
\raum{4cm}

% ---------------------------------------------------------------------------
% 8 Geometrie und Grössen
% ---------------------------------------------------------------------------
\aufgabe{8.6}
\begin{parts}
  \part Die Geraden $AA'$ und $BB'$ sind parallel. Berechnen Sie die Längen
        $\overline{SB'}$ und $\overline{BB'}$. (Skizze nicht massstäblich)

  \begin{center}
  \begin{tikzpicture}[scale=0.8]
    \coordinate (S) at (0,0);
    \coordinate (A) at (4,0);   \coordinate (B) at (6,0);
    \coordinate (A2) at (3.2,2.4); \coordinate (B2) at (4.8,3.6);
    \draw[thick] (S) -- (7,0);
    \draw[thick] (S) -- (5.6,4.2);
    \draw[thick] (A) -- (A2);
    \draw[thick] (B) -- (B2);
    \foreach \p/\pos in {S/left, A/below, B/below, A2/above left, B2/above left}
      \fill (\p) circle (2pt);
    \node[left]        at (S)  {$S$};
    \node[below]       at (A)  {$A$};
    \node[below]       at (B)  {$B$};
    \node[above left]  at (A2) {$A'$};
    \node[above left]  at (B2) {$B'$};
    \node[below, font=\small] at (2,0)       {$4$\,cm};
    \node[below, font=\small] at (5,0)       {$2$\,cm};
    \node[above left, font=\small] at (1.6,1.2) {$5$\,cm};
    \node[right, font=\small] at (3.6,1.2)   {$3$\,cm};
  \end{tikzpicture}
  \end{center}
  \part Eine $1.80$\,m grosse Person wirft einen $1.20$\,m langen Schatten. Ein
        Baum wirft gleichzeitig einen $8$\,m langen Schatten. Wie hoch ist der
        Baum?
  \part Ein Quadrat mit der Seitenlänge $a$ wird mit dem Faktor $k$ vergrössert.
        Geben Sie den Umfang und den Flächeninhalt des neuen Quadrats an.
\end{parts}
\begin{solution}
  a) $\overline{SB} = 6$\,cm, also $\overline{SB'} = 5 \cdot \frac{6}{4} =
  7.5$\,cm und $\overline{BB'} = 3 \cdot \frac{6}{4} = 4.5$\,cm \quad
  b) $12$\,m \quad
  c) Umfang $4ka$, Flächeninhalt $k^2a^2$
  \diagnose{Bei a) zeigen Ergebnisse wie $\overline{BB'} = 1.5$\,cm oder
  $\overline{SB'} = 2.5$\,cm, dass mit der Teilstrecke $\overline{AB}$ statt
  mit $\overline{SB}$ gerechnet wird. Additive Ansätze ($3 + 2 = 5$) zeigen
  fehlendes Verständnis der Ähnlichkeit. $ka^2$ bei c) zeigt, dass die Fläche
  wie die Länge skaliert wird.}
\end{solution}
\raum{4cm}

% ---------------------------------------------------------------------------
% 9 Sprache ↔ Mathematik
% ---------------------------------------------------------------------------
\aufgabe{9.6}
\begin{parts}
  \part Handyabo A kostet CHF 20 pro Monat plus CHF 0.10 pro Gesprächsminute.
        Abo B kostet pauschal CHF 35 pro Monat. Stellen Sie für beide Abos
        einen Term für die Monatskosten bei $x$ Minuten auf. Ab wie vielen
        Minuten ist Abo B günstiger?
  \part Für einen Schulausflug fahren $200$ Personen mit Bussen, die je
        $45$ Plätze haben. Wie viele Busse braucht es?
\end{parts}
\begin{solution}
  a) A: $20 + 0.1x$; B: $35$. Aus $20 + 0.1x = 35$ folgt $x = 150$. Abo B ist
  ab mehr als $150$ Minuten günstiger. \quad
  b) $200 \geteilt 45 \approx 4.4$, also braucht es $5$ Busse.
  \diagnose{Die Antwort «$150$» ohne Aussage, was darunter und darüber gilt,
  zeigt, dass die Lösung nicht im Kontext gedeutet wird. Bei b) sind $4.44$
  oder $4$ Busse die klassischen Zeichen dafür, dass ein Rechenergebnis nicht
  auf die Situation zurückbezogen wird.}
\end{solution}
\raum{4.5cm}


\end{questions}

% ===========================================================================
% Auswertungsraster (für die Lehrperson)
% ===========================================================================
\newpage
\section*{Auswertung \normalsize(wird von der Lehrperson ausgefüllt)}

Pro Aufgabe: \quad \textbf{2} = sicher gelöst \quad \textbf{1} = teilweise
gelöst \quad \textbf{0} = nicht gelöst. \\
Die \emph{Bruchstelle} ist das tiefste Niveau eines Strangs, auf dem die
Aufgaben nicht mehr sicher gelöst werden.

\begin{center}
\small
\renewcommand{\arraystretch}{2.2}
\newcommand{\feld}[1]{#1\,\framebox[0.6cm]{\rule{0pt}{0.9em}}}
\begin{tabular}{|>{\raggedright\arraybackslash}p{3.5cm}|c|c|c|c|p{1.9cm}|}
  \hline
  \textbf{Strang}
    & \textbf{A}\,\footnotesize Grundlagen & \textbf{B}\,\footnotesize Numerisch
    & \textbf{C}\,\footnotesize Symbol. & \textbf{D}\,\footnotesize Formal
    & \footnotesize\textbf{Bruchstelle} \\ \hline
  1 Zahlen und Stellenwert          & \feld{1.1}\ \feld{1.2} & \feld{1.3}\ \feld{1.4} & \feld{1.5} & \feld{1.6} & \\ \hline
  2 Operationen und Rechenreihenfolge & \feld{2.1}\ \feld{2.2} & \feld{2.3}\ \feld{2.4} & \feld{2.5} & \feld{2.6} & \\ \hline
  3 Brüche                          & \feld{3.1}\ \feld{3.2} & \feld{3.3}\ \feld{3.4} & \feld{3.5} & \feld{3.6} & \\ \hline
  4 Proportionalität und Prozent    & \feld{4.1}\ \feld{4.2} & \feld{4.3}\ \feld{4.4} & \feld{4.5} & \feld{4.6} & \\ \hline
  5 Terme                           & \feld{5.1}\ \feld{5.2} & \feld{5.3}\ \feld{5.4} & \feld{5.5} & \feld{5.6} & \\ \hline
  6 Gleichungen                     & \feld{6.1}\ \feld{6.2} & \feld{6.3}\ \feld{6.4} & \feld{6.5} & \feld{6.6} & \\ \hline
  7 Funktionen und Graphen          & \feld{7.1}\ \feld{7.2} & \feld{7.3}\ \feld{7.4} & \feld{7.5} & \feld{7.6} & \\ \hline
  8 Geometrie und Grössen           & \feld{8.1}\ \feld{8.2} & \feld{8.3}\ \feld{8.4} & \feld{8.5} & \feld{8.6} & \\ \hline
  9 Sprache $\leftrightarrow$ Mathematik & \feld{9.1}\ \feld{9.2} & \feld{9.3}\ \feld{9.4} & \feld{9.5} & \feld{9.6} & \\ \hline
\end{tabular}
\end{center}

\textbf{Bemerkungen:}

\vspace{4cm}

\end{document}
"
% Jeweils zweimal kompilieren (Seitenzahlen).
%
% Aufbau: 9 Stränge × 4 Niveaus
%   A Grundlagen  (Aufgaben x.1, x.2)
%   B Numerisch   (Aufgaben x.3, x.4)
%   C Symbolisch  (Aufgabe  x.5)
%   D Formal      (Aufgabe  x.6)
%%%%%%%%%%%%%%%%%%%%%%%%%%%%%%%%%%%%%%%%%%%%%%%%%%%%%%%%%%%%%%%%%%%%%%%%%%%%%%

\ifdefined\mitloesungen \printanswers \fi

\usepackage[T1]{fontenc}
\usepackage{lmodern}
\usepackage[nswissgerman]{babel}
\usepackage[hidelinks]{hyperref}
\usepackage{amssymb, amsmath, setspace, multicol, array, needspace}
\usepackage{tikz}
\usepackage[margin=2cm]{geometry}

\usepackage{color}
\colorgrids
\definecolor{GridColor}{gray}{0.7}
\setlength{\gridsize}{4mm}

\singlespacing
\parindent 0ex

% ---------------------------------------------------------------------------
% Kopfzeilen und Angaben
% ---------------------------------------------------------------------------
\newcommand{\titel}{Standortbestimmung Mathematik}
\newcommand{\schuljahr}{Schuljahr 2026/27}
\newcommand{\zeit}{90 Minuten (Teile A und B)}
\newcommand{\hilfsmittel}{keine (kein Taschenrechner)}

% ---------------------------------------------------------------------------
% Hilfsbefehle
% ---------------------------------------------------------------------------
% Aufgabe mit fester Nummer (Strang.Nummer) statt fortlaufender Zählung
\newcommand{\aufgabe}[1]{%
  \needspace{8\baselineskip}%
  \renewcommand{\questionlabel}{\textbf{#1}}%
  \question}

% Teil-Überschrift über die volle Breite
\newcommand{\teil}[2]{\uplevel{\vspace{1ex}\textbf{\large Teil #1: #2}\par\vspace{1ex}}}

% Divisionszeichen (Doppelpunkt) mit korrektem Abstand
\newcommand{\geteilt}{\mathbin{:}}
% Punkt mit Koordinaten (x | y)
\newcommand{\pkt}[2]{(#1\,|\,#2)}
% Leeres Kästchen für Lückenaufgaben
\newcommand{\kaestchen}{\fbox{\rule{0pt}{1.1em}\hspace{1.6em}}}
% Platz für die Antwort (Raster), nur in der Version ohne Lösungen
\newcommand{\raum}[1]{\ifprintanswers\else\fillwithgrid{#1}\fi}
% Ankreuzfelder: \wahl = falsche, \wahlr = richtige Option
\newcommand{\wahl}[1]{\mbox{$\Box$\,#1}}
\newcommand{\wahlr}[1]{\mbox{\ifprintanswers$\blacksquare$\,\textbf{#1}\else$\Box$\,#1\fi}}
% Mehr Zeilenabstand in Aufgabenteilen mit Brüchen
\newcommand{\bruchabstand}{\setlength{\itemsep}{2.5ex}}
% Interpretationshinweis in den Lösungen
\newcommand{\diagnose}[1]{\par\smallskip\textit{Interpretation:} #1}

% ---------------------------------------------------------------------------
% Deutsche Bezeichnungen für die exam-Klasse
% ---------------------------------------------------------------------------
\renewcommand{\solutiontitle}{\noindent\textbf{Lösung:}\enspace}
\checkboxchar{$\Box$}
\checkedchar{$\blacksquare$}
\CorrectChoiceEmphasis{\bfseries}
\pointsinmargin
\nopointsinmargin

\begin{document}

\pagestyle{headandfoot}
\firstpageheader{}{}{}
\runningheader{\titel\ifprintanswers\ – Lösungen\fi}{}{\schuljahr}
\runningheadrule
\footer{}{Seite \thepage\ von \numpages}{}

% ===========================================================================
% Titelseite
% ===========================================================================
\begin{center}
  {\LARGE\bfseries \titel\ifprintanswers\ – Lösungen\fi}\\[1ex]
  {\large Eintritt ins Gymnasium – \schuljahr}
\end{center}
\vspace{1em}

\begin{tabular}{@{}p{0.5\textwidth}l@{}}
  \textbf{Name:} \makebox[5.5cm]{\hrulefill} & \textbf{Klasse:} \makebox[3cm]{\hrulefill} \\[2ex]
  \textbf{Zeit:} \zeit                       & \textbf{Hilfsmittel:} \hilfsmittel
\end{tabular}

\vspace{1.5em}
\rule{\textwidth}{.4pt}

\subsection*{Hinweise}

\begin{itemize}
  \item Diese Standortbestimmung ist \textbf{keine Prüfung} und wird nicht benotet. Sie zeigt
        uns, wo Sie in Mathematik stehen, damit wir Sie gezielt unterstützen können.
  \item Die Aufgaben werden von Teil zu Teil schwieriger. Es ist normal, dass
        Sie nicht alle Aufgaben lösen können. Versuchen Sie trotzdem jede Aufgabe.
  \item Lösen Sie die Aufgaben \textbf{ohne Taschenrechner}.
  \item Schreiben Sie Ihre Überlegungen und Zwischenschritte auf. Ein Lösungsweg
        sagt uns oft mehr als das Ergebnis.
  \item Wenn Sie eine Aufgabe nicht lösen können, gehen Sie zur nächsten weiter.
\end{itemize}

\vfill
\newpage

% ===========================================================================
\begin{questions}
% ===========================================================================

\teil{A}{Grundlagen}

% ---------------------------------------------------------------------------
% 1 Zahlen und Stellenwert
% ---------------------------------------------------------------------------
\aufgabe{1.1}
Ordnen Sie die Zahlen der Grösse nach, beginnend mit der kleinsten:
\[
  0.5; \quad 0.45; \quad 0.405; \quad 0.54
\]
\begin{solution}
  $0.405 < 0.45 < 0.5 < 0.54$
  \diagnose{Wird $0.405$ oder $0.45$ grösser als $0.5$ eingeordnet, so gilt
  vermutlich die Vorstellung «mehr Ziffern bedeutet eine grössere Zahl».
  Der Stellenwert der Dezimalstellen ist nicht gesichert.}
\end{solution}
\raum{2.5cm}

\aufgabe{1.2}
\begin{parts}
  \part Berechnen Sie: \quad $3.7 \cdot 100$ \qquad $45 \geteilt 1000$
  \part Welche Zahl liegt genau in der Mitte zwischen $0.3$ und $0.4$?
\end{parts}
\begin{solution}
  a) $370$; \ $0.045$ \qquad b) $0.35$
  \diagnose{Ergebnisse wie $3.700$ oder $0.45$ deuten auf die auswendig gelernte
  Regel «Nullen anhängen bzw.\ wegstreichen» hin, ohne Verständnis des
  Stellenwerts. Bei b) zeigt die Antwort «gibt es nicht», dass die Dichte der
  Dezimalzahlen nicht verstanden ist.}
\end{solution}
\raum{3cm}

% ---------------------------------------------------------------------------
% 2 Operationen und Rechenreihenfolge
% ---------------------------------------------------------------------------
\aufgabe{2.1}
Ein Band von $3.6$\,m Länge wird in Stücke von $0.4$\,m Länge geschnitten.
Welche Rechnung passt zu der Frage, wie viele Stücke entstehen? Kreuzen Sie an.

\begin{oneparcheckboxes}
  \choice $3.6 \cdot 0.4$
  \CorrectChoice $3.6 \geteilt 0.4$
  \choice $0.4 \geteilt 3.6$
  \choice $3.6 - 0.4$
\end{oneparcheckboxes}
\begin{solution}
  $3.6 \geteilt 0.4 = 9$ Stücke.
  \diagnose{Wer die Division meidet, hat oft die Vorstellung, man dürfe nur durch
  ganze Zahlen teilen. Die Wahl von $0.4 \geteilt 3.6$ zeigt die Regel «immer die
  grössere durch die kleinere Zahl» in umgekehrter Form bzw.\ fehlendes Verständnis
  der Division als Aufteilen.}
\end{solution}

\aufgabe{2.2}
Entscheiden Sie \textbf{ohne zu rechnen}, ob das Ergebnis grösser oder kleiner
als $50$ ist.
\begin{parts}
  \part $50 \cdot 0.8$ \hfill \wahl{grösser als 50}\quad \wahlr{kleiner als 50}
  \part $50 \geteilt 0.8$ \hfill \wahlr{grösser als 50}\quad \wahl{kleiner als 50}
  \part $50 \geteilt 1.25$ \hfill \wahl{grösser als 50}\quad \wahlr{kleiner als 50}
\end{parts}
\begin{solution}
  a) kleiner ($40$) \quad b) grösser ($62.5$) \quad c) kleiner ($40$)
  \diagnose{Fehler bei a) und b) zeigen die Fehlvorstellung «Multiplizieren
  vergrössert, Dividieren verkleinert», die aus dem Rechnen mit natürlichen
  Zahlen stammt.}
\end{solution}

% ---------------------------------------------------------------------------
% 3 Brüche
% ---------------------------------------------------------------------------
\aufgabe{3.1}
Markieren Sie die Zahlen $\dfrac{3}{4}$ und $\dfrac{5}{4}$ auf dem Zahlenstrahl.

\begin{center}
\begin{tikzpicture}[x=6cm]
  \draw[->, thick] (-0.05,0) -- (2.15,0);
  \foreach \x in {0,1,2} {
    \draw[thick] (\x,0.15) -- (\x,-0.15) node[below] {$\x$};
  }
  \ifprintanswers
    \foreach \x/\l in {0.75/{\frac{3}{4}}, 1.25/{\frac{5}{4}}} {
      \draw[very thick] (\x,0.25) -- (\x,-0.25);
      \node[above] at (\x,0.25) {$\l$};
    }
  \fi
\end{tikzpicture}
\end{center}
\begin{solution}
  Markierungen bei $0.75$ und $1.25$ (siehe Zahlenstrahl).
  \diagnose{Wer $\frac{5}{4}$ nicht einträgt oder für unmöglich hält, versteht
  Brüche nur als «Teil eines Ganzen», nicht als Zahlen. Falsche Positionen von
  $\frac{3}{4}$ zeigen Schwierigkeiten beim Einteilen der Einheit.}
\end{solution}

\aufgabe{3.2}
\begin{parts}
  \part Welcher Bruch ist jeweils grösser? Begründen Sie kurz.
  \[
    \frac{2}{3} \text{ oder } \frac{3}{4} \qquad\qquad
    \frac{5}{8} \text{ oder } \frac{5}{6} \qquad\qquad
    \frac{7}{8} \text{ oder } \frac{8}{9}
  \]
  \part Ergänzen Sie: \quad $\dfrac{2}{3} = \dfrac{\kaestchen}{12}$
\end{parts}
\begin{solution}
  a) $\frac{3}{4}$; \ $\frac{5}{6}$; \ $\frac{8}{9}$ (z.\,B.\ weil $\frac{8}{9}$
  nur $\frac{1}{9}$ von $1$ entfernt ist, $\frac{7}{8}$ aber $\frac{1}{8}$)
  \qquad b) $8$
  \diagnose{Wird $\frac{5}{8}$ als grösser gewählt («8 ist grösser als 6»), so
  werden Zähler und Nenner getrennt als natürliche Zahlen betrachtet. Das letzte
  Paar unterscheidet auswendig gelernte Verfahren (gleichnamig machen) von
  Zahlverständnis (Abstand zu~1).}
\end{solution}
\raum{3.5cm}

% ---------------------------------------------------------------------------
% 4 Proportionalität und Prozent
% ---------------------------------------------------------------------------
\aufgabe{4.1}
3 Hefte kosten CHF 7.50. Wie viel kosten 12 Hefte? Wie viel kosten 5 Hefte?
\begin{solution}
  12 Hefte: CHF 30.00 \qquad 5 Hefte: CHF 12.50 (1 Heft kostet CHF 2.50)
  \diagnose{Gelingen die 12 Hefte, die 5 Hefte aber nicht, so wird nur
  vervielfacht, der Weg über den Einheitspreis (Dreisatz) fehlt.}
\end{solution}
\raum{3cm}

\aufgabe{4.2}
Für 4 Personen braucht ein Rezept 6 Eier. Wie viele Eier braucht man für
6 Personen?
\begin{solution}
  $9$ Eier
  \diagnose{Die Antwort $8$ («2 Personen mehr, also 2 Eier mehr») ist das
  klassische Zeichen für additives statt multiplikatives Denken.}
\end{solution}
\raum{3cm}

% ---------------------------------------------------------------------------
% 5 Terme
% ---------------------------------------------------------------------------
\aufgabe{5.1}
Aus Hölzchen werden Quadrate in einer Reihe gelegt.

\begin{center}
\begin{tikzpicture}[x=0.8cm, y=0.8cm, line width=1.6pt, line cap=round]
  \newcommand{\reihe}[2]{% #1 = x-Start, #2 = Anzahl Quadrate
    \foreach \i in {1,...,#2} {
      \draw (#1+\i-1,0) -- (#1+\i,0);
      \draw (#1+\i-1,1) -- (#1+\i,1);
      \draw (#1+\i,0) -- (#1+\i,1);
    }
    \draw (#1,0) -- (#1,1);
  }
  \reihe{0}{1}  \node at (0.5,-0.6) {\small 4 Hölzchen};
  \reihe{3}{2}  \node at (4,-0.6)   {\small 7 Hölzchen};
  \reihe{7}{3}  \node at (8.5,-0.6) {\small 10 Hölzchen};
\end{tikzpicture}
\end{center}
\begin{parts}
  \part Wie viele Hölzchen braucht man für 10 Quadrate?
  \part Wie viele Hölzchen braucht man für $n$ Quadrate?
\end{parts}
\begin{solution}
  a) $31$ \qquad b) $3n + 1$ (oder gleichwertig, z.\,B.\ $4 + 3 \cdot (n-1)$)
  \diagnose{Gelingt a), aber nicht b), so kann ein Muster an Beispielen
  fortgesetzt, aber nicht verallgemeinert werden. Antworten wie $4n$ zeigen, dass
  die gemeinsamen Hölzchen nicht beachtet werden.}
\end{solution}
\raum{3.5cm}

\aufgabe{5.2}
Welcher Ausdruck ist grösser: $n + 2$ oder $2n$? Begründen Sie.
\begin{solution}
  Das hängt von $n$ ab: Für $n = 2$ sind beide gleich, für $n < 2$ ist $n + 2$
  grösser, für $n > 2$ ist $2n$ grösser.
  \diagnose{Die Antwort «$2n$, weil verdoppeln mehr ist» oder die Prüfung mit
  nur einem Wert zeigen, dass die Variable als feste Zahl und nicht als
  Platzhalter für viele Zahlen verstanden wird.}
\end{solution}
\raum{3cm}

% ---------------------------------------------------------------------------
% 6 Gleichungen
% ---------------------------------------------------------------------------
\aufgabe{6.1}
Welche Zahl gehört ins Kästchen?
\begin{parts}
  \part $8 + 4 = \kaestchen + 5$
  \part $3 \cdot \kaestchen = 24 \geteilt 2$
\end{parts}
\begin{solution}
  a) $7$ \qquad b) $4$
  \diagnose{Die Antworten $12$ oder $17$ bei a) zeigen, dass das
  Gleichheitszeichen als Aufforderung «jetzt kommt das Ergebnis» gelesen wird
  und nicht als Beziehung zwischen zwei gleichwertigen Seiten.}
\end{solution}

\aufgabe{6.2}
Ist die Aussage wahr oder falsch? Begründen Sie, möglichst ohne die Ergebnisse
auszurechnen.
\begin{parts}
  \part $37 + 25 = 36 + 26$ \hfill \wahlr{wahr}\quad \wahl{falsch}
  \part $47 - 18 = 48 - 17$ \hfill \wahl{wahr}\quad \wahlr{falsch}
  \part $6 \cdot 15 = 3 \cdot 30$ \hfill \wahlr{wahr}\quad \wahl{falsch}
\end{parts}
\begin{solution}
  a) wahr \quad b) falsch ($29 \neq 31$) \quad c) wahr
  \diagnose{Begründungen wie «links eins weniger, rechts eins mehr» zeigen
  relationales Denken. Wer b) als wahr einstuft, überträgt die
  Ausgleichsregel der Addition unreflektiert auf die Subtraktion.}
\end{solution}
\raum{3cm}

% ---------------------------------------------------------------------------
% 7 Funktionen und Graphen
% ---------------------------------------------------------------------------
\aufgabe{7.1}
Geben Sie die Koordinaten der Punkte $A$, $B$ und $C$ an.

\begin{center}
\begin{tikzpicture}[scale=0.55]
  \draw[gray!40, thin] (-5,-5) grid (5,5);
  \draw[->, thick] (-5.4,0) -- (5.6,0) node[right] {$x$};
  \draw[->, thick] (0,-5.4) -- (0,5.6) node[above] {$y$};
  \foreach \t in {-4,-2,2,4} {
    \node[below, font=\scriptsize] at (\t,0) {$\t$};
    \node[left,  font=\scriptsize] at (0,\t) {$\t$};
  }
  \foreach \p/\x/\y/\pos in {A/3/-2/right, B/-2/3/left, C/0/4/right} {
    \fill (\x,\y) circle (3.5pt);
    \node[\pos, font=\small] at (\x,\y) {$\p$};
  }
\end{tikzpicture}
\end{center}
$A = $ \fillin[$\pkt{3}{-2}$][2.5cm] \qquad
$B = $ \fillin[$\pkt{-2}{3}$][2.5cm] \qquad
$C = $ \fillin[$\pkt{0}{4}$][2.5cm]
\begin{solution}
  $A\pkt{3}{-2}$, \ $B\pkt{-2}{3}$, \ $C\pkt{0}{4}$
  \diagnose{Vertauschte Koordinaten bei $A$ und $B$ zeigen, dass die
  Reihenfolge $(x\,|\,y)$ nicht gesichert ist. Fehler bei $C$ (z.\,B.\ $(4\,|\,0)$)
  deuten auf Unsicherheit bei Punkten auf den Achsen hin.}
\end{solution}

\aufgabe{7.2}
Die $y$-Werte in der Tabelle werden alle nach derselben Regel aus den
$x$-Werten berechnet. Ergänzen Sie die Tabelle sinngemäss.
\begin{center}
  \renewcommand{\arraystretch}{1.4}
  \begin{tabular}{|c|*{4}{>{\centering\arraybackslash}p{1.2cm}|}}
    \hline
    $x$ & $0$ & $1$ & $3$  & $6$ \\ \hline
    $y$ & $2$ & $5$ & $11$ & \ifprintanswers $\mathbf{20}$\fi \\ \hline
  \end{tabular}
\end{center}
\begin{solution}
  $y = 20$ (Zuordnung $y = 3x + 2$)
  \diagnose{Da die $x$-Werte nicht in gleichen Schritten wachsen, scheitert ein
  rein schrittweises Weiterzählen (z.\,B.\ «immer $+6$» ergibt $17$). Das zeigt,
  ob der Zusammenhang zwischen $x$ und $y$ erkannt wird.}
\end{solution}
\raum{2.5cm}

% ---------------------------------------------------------------------------
% 8 Geometrie und Grössen
% ---------------------------------------------------------------------------
\aufgabe{8.1}
Ergänzen Sie.
\begin{parts}
  \part $3.5\text{ m} = $ \fillin[$350$][2cm] cm
  \part $250\text{ g} = $ \fillin[$0.25$][2cm] kg
  \part $1\text{ m}^2 = $ \fillin[$10\,000$][2cm] cm$^2$
  \part $2\text{ h }15\text{ min} = $ \fillin[$2.25$][2cm] h
\end{parts}
\begin{solution}
  a) $350$ \quad b) $0.25$ \quad c) $10\,000$ \quad d) $2.25$
  \diagnose{Die Antwort $100$ bei c) zeigt, dass die Umrechnungszahl für
  Flächeneinheiten nicht verstanden ist. Die Antwort $2.15$ bei d) zeigt, dass
  Zeitangaben wie Dezimalzahlen behandelt werden.}
\end{solution}

\aufgabe{8.2}
Ein Rechteck ist $6$\,cm lang und $4$\,cm breit.
\begin{parts}
  \part Berechnen Sie den Umfang und den Flächeninhalt.
  \part Hat jedes Rechteck mit demselben Umfang auch denselben Flächeninhalt?
        Begründen Sie.
\end{parts}
\begin{solution}
  a) Umfang $20$\,cm, Flächeninhalt $24$\,cm$^2$
  \qquad b) Nein, z.\,B.\ hat ein Rechteck mit $7$\,cm und $3$\,cm ebenfalls den
  Umfang $20$\,cm, aber den Flächeninhalt $21$\,cm$^2$.
  \diagnose{Verwechslungen von Umfang und Fläche (auch bei den Einheiten) sowie
  die Antwort «ja» bei b) zeigen, dass die beiden Grössen nicht klar
  unterschieden werden.}
\end{solution}
\raum{3.5cm}

% ---------------------------------------------------------------------------
% 9 Sprache ↔ Mathematik
% ---------------------------------------------------------------------------
\aufgabe{9.1}
Schreiben Sie die Rechnung auf und berechnen Sie das Ergebnis.
\begin{parts}
  \part die Summe von $12$ und $8$
  \part die Differenz von $20$ und $5$
  \part das Produkt von $6$ und $7$
  \part der Quotient von $36$ und $4$
\end{parts}
\begin{solution}
  a) $12 + 8 = 20$ \quad b) $20 - 5 = 15$ \quad c) $6 \cdot 7 = 42$ \quad
  d) $36 \geteilt 4 = 9$
  \diagnose{Unbekannte Fachbegriffe zeigen sich hier direkt. Bei b) und d)
  auf die Reihenfolge in der aufgeschriebenen Rechnung achten ($5 - 20$,
  $4 \geteilt 36$).}
\end{solution}
\raum{3cm}

\aufgabe{9.2}
\begin{parts}
  \part Mia hat dreimal so viele Karten wie Leo. Leo hat 12 Karten. Wie viele
        Karten hat Mia?
  \part Nora hat 12 Karten. Das sind dreimal so viele wie Tim. Wie viele Karten
        hat Tim?
\end{parts}
\begin{solution}
  a) $36$ \qquad b) $4$
  \diagnose{Die Antwort $36$ bei b) zeigt eine Schlüsselwort-Strategie
  («mal» bedeutet multiplizieren), ohne die Beziehung zwischen den Grössen zu
  erfassen.}
\end{solution}
\raum{3cm}

% ===========================================================================
\newpage
\teil{B}{Rechnen mit Zahlen}
% ===========================================================================

% ---------------------------------------------------------------------------
% 1 Zahlen und Stellenwert
% ---------------------------------------------------------------------------
\aufgabe{1.3}
\begin{parts}
  \part Runden Sie.
  \begin{subparts}
    \subpart $3.4567$ auf zwei Dezimalstellen
    \subpart $0.0496$ auf zwei Dezimalstellen
    \subpart $28\,349$ auf Tausender
  \end{subparts}
  \part Welches Ergebnis ist plausibel? Kreuzen Sie an, ohne genau zu rechnen.
        \[
          0.48 \cdot 205 \approx \quad
          \text{\wahl{$10$}} \quad \text{\wahlr{$100$}} \quad
          \text{\wahl{$1000$}} \quad \text{\wahl{$10\,000$}}
        \]
\end{parts}
\begin{solution}
  a) $3.46$; \ $0.05$; \ $28\,000$ \qquad b) $100$ (genau: $98.4$)
  \diagnose{$3.45$ zeigt Abschneiden statt Runden. Antworten wie $0.5$ oder
  $0.04$ bei $0.0496$ zeigen, dass der Stellenwert beim Runden verloren geht.
  Fehler bei b) deuten auf fehlende Überschlagsrechnung hin.}
\end{solution}
\raum{3cm}

\aufgabe{1.4}
\begin{parts}
  \part Ordnen Sie der Grösse nach, beginnend mit der kleinsten Zahl:
        \quad $-3.5; \quad -3; \quad 0.5; \quad -4; \quad -0.5$
  \part Welche Zahl ist um $5$ grösser als $-3$?
\end{parts}
\begin{solution}
  a) $-4 < -3.5 < -3 < -0.5 < 0.5$ \qquad b) $2$
  \diagnose{Wird $-3.5$ grösser als $-3$ eingeordnet, so werden negative Zahlen
  nach ihrem Betrag geordnet. Antworten wie $-8$ oder $8$ bei b) zeigen
  Unsicherheit beim Bewegen auf dem Zahlenstrahl.}
\end{solution}
\raum{2.5cm}

% ---------------------------------------------------------------------------
% 2 Operationen und Rechenreihenfolge
% ---------------------------------------------------------------------------
\aufgabe{2.3}
Berechnen Sie.
\begin{multicols}{3}
\begin{parts}
  \part $-7 + 12$
  \part $-4 - 9$
  \part $(-6) \cdot (-3)$
  \part $(-15) \geteilt 3$
  \part $5 - (-8)$
\end{parts}
\end{multicols}
\begin{solution}
  a) $5$ \quad b) $-13$ \quad c) $18$ \quad d) $-5$ \quad e) $13$
  \diagnose{Typische Fehler sind $-5$ oder $13$ bei b) sowie $-3$ bei e). Sie
  zeigen, dass Vorzeichen und Rechenzeichen nicht sicher unterschieden werden.}
\end{solution}
\raum{2.5cm}

\aufgabe{2.4}
Berechnen Sie.
\begin{multicols}{3}
\begin{parts}
  \part $12 - 3 \cdot 4$
  \part $(12 - 3) \cdot 4$
  \part $20 \geteilt 4 \cdot 5$
  \part $2 \cdot 3^2$
  \part $18 - 6 \geteilt 2 + 1$
\end{parts}
\end{multicols}
\begin{solution}
  a) $0$ \quad b) $36$ \quad c) $25$ \quad d) $18$ \quad e) $16$
  \diagnose{Typische Fehler: c) $1$ (Multiplikation vor Division), d) $36$
  (zuerst multipliziert, dann quadriert), e) $7$ (von links nach rechts
  gerechnet) oder $14$ (Addition vor Subtraktion). Sie zeigen, welche Regel der
  Rechenreihenfolge fehlt.}
\end{solution}
\raum{3cm}

% ---------------------------------------------------------------------------
% 3 Brüche
% ---------------------------------------------------------------------------
\aufgabe{3.3}
Berechnen Sie und kürzen Sie das Ergebnis so weit wie möglich.
\begin{multicols}{3}
\begin{parts}\bruchabstand
  \part $\dfrac{2}{3} + \dfrac{1}{4}$
  \part $\dfrac{5}{6} - \dfrac{1}{2}$
  \part $\dfrac{3}{4} \cdot \dfrac{2}{9}$
  \part $\dfrac{2}{5} \geteilt \dfrac{4}{15}$
  \part $3 \cdot \dfrac{2}{5}$
\end{parts}
\end{multicols}
\begin{solution}
  a) $\frac{11}{12}$ \quad b) $\frac{1}{3}$ \quad c) $\frac{1}{6}$ \quad
  d) $\frac{3}{2}$ \quad e) $\frac{6}{5}$
  \diagnose{$\frac{3}{7}$ bei a) zeigt, dass Zähler und Nenner getrennt addiert
  werden. $\frac{6}{15}$ bei e) zeigt, dass Zähler und Nenner mit $3$
  multipliziert werden (Verwechslung mit dem Erweitern). Fehler nur bei d)
  deuten auf eine nicht verstandene Divisionsregel hin.}
\end{solution}
\raum{4cm}

\aufgabe{3.4}
Ergänzen Sie die Tabelle.
\begin{center}
  \newcommand{\lsg}[1]{\ifprintanswers$\mathbf{#1}$\fi}
  \newcommand{\zh}{\rule[-2.4ex]{0pt}{6.6ex}}% Zeilenhöhe für Brüche
  \begin{tabular}{|c|c|c|}
    \hline
    \zh\makebox[2.6cm]{\textbf{Bruch}} & \makebox[2.6cm]{\textbf{Dezimalzahl}}
      & \makebox[2.6cm]{\textbf{Prozent}} \\ \hline
    \zh$\dfrac{3}{4}$          & \lsg{0.75}           & \lsg{75\,\%} \\ \hline
    \zh\lsg{\dfrac{1}{20}}     & $0.05$               & \lsg{5\,\%} \\ \hline
    \zh\lsg{\dfrac{5}{4}}      & \lsg{1.25}           & $125\,\%$ \\ \hline
    \zh$\dfrac{1}{3}$          & \lsg{0.333\ldots}    & \lsg{33.3\,\%} \\ \hline
  \end{tabular}
\end{center}
\begin{solution}
  Siehe Tabelle. Für $\frac{1}{3}$ sind auch $0.\overline{3}$ und
  $33\frac{1}{3}\,\%$ korrekt.
  \diagnose{$0.05 = 50\,\%$ zeigt Unsicherheit beim Stellenwert. Wer bei
  $125\,\%$ keinen Bruch findet, hält Prozentangaben über $100\,\%$ oder Brüche
  grösser als $1$ für unmöglich.}
\end{solution}
\raum{2.5cm}

% ---------------------------------------------------------------------------
% 4 Proportionalität und Prozent
% ---------------------------------------------------------------------------
\aufgabe{4.3}
Ein Pullover kostet CHF 80.
\begin{parts}
  \part Berechnen Sie $15\,\%$ von CHF 80.
  \part Der Preis des Pullovers wird um $25\,\%$ reduziert. Wie viel kostet er
        jetzt?
\end{parts}
\begin{solution}
  a) CHF 12 \qquad b) CHF 60
  \diagnose{Wer bei b) nur den Rabatt (CHF 20) angibt, unterscheidet nicht
  zwischen Prozentwert und neuem Wert. Fehler bei a) zeigen, dass die
  Grundidee von «Prozent als Hundertstel» fehlt.}
\end{solution}
\raum{3cm}

\aufgabe{4.4}
Ein Preis steigt von CHF 40 auf CHF 50.
\begin{parts}
  \part Um wie viel Prozent ist der Preis gestiegen?
  \part Danach sinkt der Preis wieder von CHF 50 auf CHF 40. Um wie viel Prozent
        ist er gesunken?
\end{parts}
\begin{solution}
  a) $25\,\%$ \qquad b) $20\,\%$
  \diagnose{Die Antwort $25\,\%$ bei b) (oder $20\,\%$ bei a)) zeigt, dass
  nicht erkannt wird, dass sich der Prozentsatz auf den jeweiligen Grundwert
  bezieht. Dieser Fehler ist sehr häufig.}
\end{solution}
\raum{3cm}

% ---------------------------------------------------------------------------
% 5 Terme
% ---------------------------------------------------------------------------
\aufgabe{5.3}
Vereinfachen Sie so weit wie möglich.
\begin{multicols}{3}
\begin{parts}
  \part $3a + 5 + 2a - 1$
  \part $4x + 3y - x + 2y$
  \part $2a + 3b$
  \part $3x \cdot 4x$
  \part $a^2 + a + a^2$
\end{parts}
\end{multicols}
\begin{solution}
  a) $5a + 4$ \quad b) $3x + 5y$ \quad c) nicht vereinfachbar \quad
  d) $12x^2$ \quad e) $2a^2 + a$
  \diagnose{$5ab$ bei c) sowie $a^4$, $3a^2$ oder $2a^3$ bei e) zeigen, dass
  ungleichartige Terme zusammengefasst werden. Die Aussage «$2a + 3b$ ist keine
  Lösung» weist auf die Vorstellung hin, dass ein Term ein einzelnes Ergebnis
  haben muss.}
\end{solution}
\raum{3cm}

\aufgabe{5.4}
Ist die Aussage richtig oder falsch? Kreuzen Sie an.
\begin{parts}
  \part $2a + 3b = 5ab$     \hfill \wahl{richtig}\quad \wahlr{falsch}
  \part $x + x = x^2$       \hfill \wahl{richtig}\quad \wahlr{falsch}
  \part $3x \cdot 2 = 6x$   \hfill \wahlr{richtig}\quad \wahl{falsch}
  \part $a + a^2 = a^3$     \hfill \wahl{richtig}\quad \wahlr{falsch}
  \part $4y - y = 4$        \hfill \wahl{richtig}\quad \wahlr{falsch}
\end{parts}
\begin{solution}
  a) falsch \quad b) falsch \quad c) richtig \quad d) falsch \quad e) falsch
  \diagnose{Wiedererkennungsversion von 5.3. Wer hier richtig entscheidet, bei
  5.3 aber Fehler macht, kennt die Regeln, wendet sie aber nicht sicher an.
  Bei e) zeigt «richtig», dass $y$ als Einheit oder Abkürzung gelesen wird
  («4 Äpfel minus 1 Apfel»).}
\end{solution}

% ---------------------------------------------------------------------------
% 6 Gleichungen
% ---------------------------------------------------------------------------
\aufgabe{6.3}
Lösen Sie die Gleichungen.
\begin{multicols}{3}
\begin{parts}\bruchabstand
  \part $x + 7 = 3$
  \part $3x = 21$
  \part $4x - 5 = 11$
  \part $12 = 5x + 2$
  \part $\dfrac{x}{3} = 6$
\end{parts}
\end{multicols}
\begin{solution}
  a) $x = -4$ \quad b) $x = 7$ \quad c) $x = 4$ \quad d) $x = 2$ \quad
  e) $x = 18$
  \diagnose{Bei d) steht die Variable rechts, was manche Schülerinnen und
  Schüler verunsichert. $x = 2$ bei e) zeigt, dass Umkehroperationen
  verwechselt werden.}
\end{solution}
\raum{3.5cm}

\aufgabe{6.4}
Lösen Sie die Gleichungen.
\begin{multicols}{3}
\begin{parts}
  \part $5x + 3 = 2x + 15$
  \part $7 - 2x = 1$
  \part $3x = 0$
\end{parts}
\end{multicols}
\begin{solution}
  a) $x = 4$ \quad b) $x = 3$ \quad c) $x = 0$
  \diagnose{$x = -3$ bei b) zeigt Vorzeichenfehler beim Umformen. Bei c)
  antworten manche «keine Lösung» oder «$x$ ist nicht definiert», weil $0$
  nicht als Lösung akzeptiert wird.}
\end{solution}
\raum{3.5cm}

% ---------------------------------------------------------------------------
% 7 Funktionen und Graphen
% ---------------------------------------------------------------------------
\aufgabe{7.3}
Der Graph zeigt, wie weit eine Person von zu Hause entfernt ist.

\begin{center}
\begin{tikzpicture}[x=0.22cm, y=1.1cm]
  \draw[gray!35, thin, xstep=5, ystep=0.5] (0,0) grid (45,3.5);
  \draw[->, thick] (0,0) -- (47,0) node[right] {\small Zeit (min)};
  \draw[->, thick] (0,0) -- (0,3.9) node[above] {\small Entfernung (km)};
  \foreach \t in {0,10,20,30,40} \node[below, font=\scriptsize] at (\t,0) {$\t$};
  \foreach \s in {1,2,3} \node[left, font=\scriptsize] at (0,\s) {$\s$};
  \draw[very thick] (0,0) -- (20,3) -- (30,3) -- (40,0);
\end{tikzpicture}
\end{center}
\begin{parts}
  \part Wie weit ist die Person nach 20 Minuten von zu Hause entfernt?
  \part Was geschieht zwischen der 20.\ und der 30.\ Minute?
  \part In welchem Abschnitt bewegt sich die Person am schnellsten?
        Begründen Sie.
  \part Berechnen Sie die Geschwindigkeit im ersten Abschnitt in km/h.
\end{parts}
\begin{solution}
  a) $3$\,km \quad b) Die Person bleibt stehen (macht eine Pause). \quad
  c) Auf dem Rückweg (30.--40.\ Minute), weil der Graph dort am steilsten ist.
  \quad d) $3$\,km in $20$\,min, also $9$\,km/h.
  \diagnose{Wird bei c) der «höchste» Abschnitt gewählt, so wird die Höhe des
  Graphen statt seiner Steigung gelesen. Bei b) zeigt «sie geht geradeaus
  weiter», dass der Graph als Bild des Weges gedeutet wird.}
\end{solution}
\raum{3.5cm}

\aufgabe{7.4}
Eine Velofahrerin fährt über einen Hügel: zuerst bergauf, dann bergab.
Welcher Graph zeigt ihre Geschwindigkeit $v$ in Abhängigkeit von der Zeit $t$
am besten? Kreuzen Sie an.

\begin{center}
\newcommand{\achsen}{%
  \draw[->] (0,0) -- (2.4,0) node[below] {\scriptsize $t$};
  \draw[->] (0,0) -- (0,1.7) node[left]  {\scriptsize $v$};}
\begin{tikzpicture}[thick]
  \begin{scope}[xshift=0cm]
    \achsen \draw[very thick] (0,0.3) .. controls (0.8,1.6) and (1.4,1.6) .. (2.2,0.3);
    \node at (1.1,-0.5) {\wahl{Graph 1}};
  \end{scope}
  \begin{scope}[xshift=3.3cm]
    \achsen \draw[very thick] (0,1.0) .. controls (0.6,0.3) and (0.9,0.3) .. (1.1,0.35)
                                     .. controls (1.4,0.5) and (1.8,1.3) .. (2.2,1.5);
    \node at (1.1,-0.5) {\wahlr{Graph 2}};
  \end{scope}
  \begin{scope}[xshift=6.6cm]
    \achsen \draw[very thick] (0,1.4) -- (2.2,0.3);
    \node at (1.1,-0.5) {\wahl{Graph 3}};
  \end{scope}
  \begin{scope}[xshift=9.9cm]
    \achsen \draw[very thick] (0,0.3) -- (1.1,1.3) -- (2.2,1.3);
    \node at (1.1,-0.5) {\wahl{Graph 4}};
  \end{scope}
\end{tikzpicture}
\end{center}
\begin{solution}
  Graph 2: Bergauf nimmt die Geschwindigkeit ab, bergab nimmt sie zu.
  \diagnose{Die Wahl von Graph 1 zeigt die klassische Fehlvorstellung «Graph
  als Bild der Situation»: Der Graph sieht aus wie der Hügel, beschreibt aber
  nicht die Geschwindigkeit.}
\end{solution}

% ---------------------------------------------------------------------------
% 8 Geometrie und Grössen
% ---------------------------------------------------------------------------
\aufgabe{8.3}
Aus einem Rechteck wird an einer Ecke ein kleineres Rechteck herausgeschnitten.
Berechnen Sie den Flächeninhalt und den Umfang der Figur.

\begin{center}
\begin{tikzpicture}[scale=0.5]
  \draw[very thick, fill=gray!10] (0,0) -- (8,0) -- (8,4) -- (5,4) -- (5,6) -- (0,6) -- cycle;
  \draw[dashed, gray] (5,4) -- (5,6) (5,6) -- (8,6) -- (8,4);
  \node[below] at (4,0) {\small $8$\,cm};
  \node[left]  at (0,3) {\small $6$\,cm};
  \node[above] at (6.5,6) {\small $3$\,cm};
  \node[right] at (8,5) {\small $2$\,cm};
\end{tikzpicture}
\end{center}
\begin{solution}
  Flächeninhalt: $8 \cdot 6 - 3 \cdot 2 = 42$\,cm$^2$ \qquad
  Umfang: $28$\,cm (gleich wie beim ganzen Rechteck)
  \diagnose{Wer beim Umfang etwas abzieht (z.\,B.\ $28 - 10 = 18$\,cm),
  überträgt die Überlegung zur Fläche auf den Umfang. Das zeigt, dass Umfang
  und Fläche nicht als verschiedene Grössen verstanden werden.}
\end{solution}
\raum{3.5cm}

\aufgabe{8.4}
\begin{parts}
  \part In einem Dreieck sind zwei Winkel $47^\circ$ und $68^\circ$ gross. Wie
        gross ist der dritte Winkel?
  \part Ein Aquarium ist $50$\,cm lang, $40$\,cm breit und $30$\,cm hoch. Wie
        viele Liter Wasser fasst es?
\end{parts}
\begin{solution}
  a) $65^\circ$ \qquad b) $60\,000$\,cm$^3 = 60$\,dm$^3 = 60$\,Liter
  \diagnose{Ergebnisse wie $600$\,L oder $6$\,L bei b) zeigen, dass der
  Zusammenhang $1\,\text{dm}^3 = 1$\,L bzw.\ die Umrechnung von
  Volumeneinheiten nicht gesichert ist.}
\end{solution}
\raum{3cm}

% ---------------------------------------------------------------------------
% 9 Sprache ↔ Mathematik
% ---------------------------------------------------------------------------
\aufgabe{9.3}
Schreiben Sie als \textbf{einen} Rechenausdruck und berechnen Sie das Ergebnis.
\begin{parts}
  \part das Dreifache der Differenz von $12$ und $5$
  \part die Differenz aus dem Dreifachen von $12$ und $5$
  \part Subtrahieren Sie $4$ von $10$.
  \part die Hälfte der Summe von $9$ und $15$
\end{parts}
\begin{solution}
  a) $3 \cdot (12 - 5) = 21$ \quad b) $3 \cdot 12 - 5 = 31$ \quad
  c) $10 - 4 = 6$ \quad d) $(9 + 15) \geteilt 2 = 12$
  \diagnose{a) und b) enthalten dieselben Wörter in anderer Reihenfolge und
  ergeben verschiedene Strukturen. Fehlende Klammern bei a) oder d) zeigen,
  dass die Struktur des Satzes nicht in die Rechnung übertragen wird. $4 - 10$
  bei c) ist die typische Umkehrung der Reihenfolge.}
\end{solution}
\raum{3.5cm}

\aufgabe{9.4}
Eine Klasse mit 24 Schülerinnen und Schülern macht einen Ausflug. Die Bahnfahrt
kostet pro Person CHF 18.50, der Museumseintritt pro Person CHF 12. Die Klasse
fährt um 8:15 Uhr ab und ist um 16:40 Uhr zurück. Die Schule übernimmt
CHF 300 der Gesamtkosten. Wie viel muss jede Person selbst bezahlen?
\begin{solution}
  Kosten pro Person: $18.50 + 12 = 30.50$; Gesamtkosten:
  $24 \cdot 30.50 = 732$; nach Abzug: $732 - 300 = 432$; pro Person:
  $432 \geteilt 24 = 18$. Jede Person bezahlt CHF 18.
  \diagnose{Die Zeitangaben sind überflüssig. Wer sie verwendet, kann die
  relevanten Informationen nicht auswählen. Wird CHF 300 von den Kosten einer
  Person abgezogen, so fehlt die Unterscheidung zwischen Gesamt- und
  Einzelgrösse.}
\end{solution}
\raum{4.5cm}

% ===========================================================================
\newpage
\teil{C}{Rechnen mit Variablen}
% ===========================================================================

% ---------------------------------------------------------------------------
% 1 Zahlen und Stellenwert
% ---------------------------------------------------------------------------
\aufgabe{1.5}
\begin{parts}
  \part Schreiben Sie in wissenschaftlicher Schreibweise:
        \quad $45\,000\,000$ \qquad $0.00072$
  \part Schreiben Sie als Dezimalzahl:
        \quad $3.1 \cdot 10^5$ \qquad $6 \cdot 10^{-3}$
  \part Berechnen Sie die beiden Ausdrücke und geben Sie die Ergebnisse in
        wissenschaftlicher Schreibweise bzw.\ als Zehnerpotenz an.
  \begin{multicols}{2}
  \begin{subparts}
    \subpart $(2 \cdot 10^3) \cdot (4 \cdot 10^5)$
    \subpart $10^6 \geteilt 10^2$
  \end{subparts}
  \end{multicols}
\end{parts}
\begin{solution}
  a) $4.5 \cdot 10^7$; \ $7.2 \cdot 10^{-4}$ \quad
  b) $310\,000$; \ $0.006$ \quad
  c) $8 \cdot 10^8$; \ $10^4$
  \diagnose{Falsch gezählte Stellen zeigen fehlendes Stellenwertverständnis.
  Werden negative Exponenten als negative Zahlen gedeutet ($6 \cdot 10^{-3}
  = -6000$), fehlt die Bedeutung der Schreibweise. $10^{15}$ bzw.\ $10^3$ bei
  c) zeigen, dass Exponenten multipliziert bzw.\ dividiert statt addiert bzw.\
  subtrahiert werden.}
\end{solution}
\raum{3.5cm}

% ---------------------------------------------------------------------------
% 2 Operationen und Rechenreihenfolge
% ---------------------------------------------------------------------------
\aufgabe{2.5}
Berechnen Sie den Wert der Terme für $a = -2$ und $b = 3$.
\begin{multicols}{3}
\begin{parts}
  \part $3a + b$
  \part $a^2$
  \part $-a^2$
  \part $2b - a$
  \part $a \cdot b^2$
\end{parts}
\end{multicols}
\begin{solution}
  a) $-3$ \quad b) $4$ \quad c) $-4$ \quad d) $8$ \quad e) $-18$
  \diagnose{Eine der aussagekräftigsten Aufgaben. $-4$ bei b) zeigt, dass ohne
  Klammern eingesetzt wird ($-2^2$ statt $(-2)^2$). $4$ bei c) zeigt
  Unsicherheit mit dem Minuszeichen vor der Potenz. $36$ bei e) zeigt, dass
  $(ab)^2$ gerechnet wird.}
\end{solution}
\raum{3cm}

% ---------------------------------------------------------------------------
% 3 Brüche
% ---------------------------------------------------------------------------
\aufgabe{3.5}
Fassen Sie zusammen und vereinfachen Sie so weit wie möglich.
\begin{multicols}{2}
\begin{parts}\bruchabstand
  \part $\dfrac{a}{4} + \dfrac{b}{6}$
  \part $\dfrac{x}{3} - \dfrac{x}{5}$
  \part $\dfrac{3}{x} + \dfrac{2}{x}$
  \part $\dfrac{2}{a} \cdot \dfrac{a}{6}$
\end{parts}
\end{multicols}
\begin{solution}
  a) $\frac{3a + 2b}{12}$ \quad b) $\frac{2x}{15}$ \quad c) $\frac{5}{x}$ \quad
  d) $\frac{1}{3}$
  \diagnose{Diese Aufgabe entspricht 3.3 mit Variablen. Gelingt 3.3, aber nicht
  3.5, so liegt die Schwierigkeit beim Übergang zur Algebra, nicht bei den
  Bruchregeln. $\frac{5}{2x}$ bei c) oder $\frac{a+b}{10}$ bei a) zeigen
  getrenntes Addieren von Zählern und Nennern.}
\end{solution}
\raum{4cm}

% ---------------------------------------------------------------------------
% 4 Proportionalität und Prozent
% ---------------------------------------------------------------------------
\aufgabe{4.5}
\begin{parts}
  \part Entscheiden Sie bei jeder Tabelle, ob die Zuordnung proportional,
        umgekehrt proportional oder keines von beidem ist. Geben Sie bei
        proportionalen Zuordnungen die Gleichung $y = k \cdot x$ an.
  \begin{center}
    \renewcommand{\arraystretch}{1.3}
    \newcommand{\tab}[5]{%
      \begin{tabular}[t]{@{}c@{}}
        \small Tabelle #1 \\[0.5ex]
        \begin{tabular}{|c|c|c|c|c|}
          \hline $x$ & 1 & 2 & 4 & 5 \\ \hline
                 $y$ & #2 & #3 & #4 & #5 \\ \hline
        \end{tabular}
      \end{tabular}}
    \tab{1}{3}{6}{12}{15} \qquad \tab{2}{5}{7}{11}{13} \qquad \tab{3}{20}{10}{5}{4}
  \end{center}
  \part $4$ Personen brauchen für eine Arbeit $6$ Tage. Wie lange brauchen
        $3$ Personen bei gleichem Arbeitstempo?
\end{parts}
\begin{solution}
  a) Tabelle 1: proportional, $y = 3x$; \ Tabelle 2: keines von beidem
  ($y = 2x + 3$); \ Tabelle 3: umgekehrt proportional ($x \cdot y = 20$)
  \qquad b) $8$ Tage
  \diagnose{Tabelle 2 ist die zentrale Falle: «gleichmässig wachsend» wird mit
  proportional verwechselt. $4.5$ Tage bei b) zeigen, dass ohne Prüfung direkt
  proportional gerechnet wird.}
\end{solution}
\raum{3.5cm}

% ---------------------------------------------------------------------------
% 5 Terme
% ---------------------------------------------------------------------------
\aufgabe{5.5}
\begin{parts}
  \part Multiplizieren Sie aus und vereinfachen Sie.
  \begin{multicols}{4}
    $3 \cdot (x + 4)$ \columnbreak

    $(-2) \cdot (a - 5)$ \columnbreak

    $(x + 3)(x + 2)$ \columnbreak

    $(a + b)^2$
  \end{multicols}
  \part Klammern Sie so weit wie möglich aus.
        \qquad $6x + 9$ \qquad\qquad $x^2 - 5x$
\end{parts}
\begin{solution}
  a) $3x + 12$; \ $-2a + 10$; \ $x^2 + 5x + 6$; \ $a^2 + 2ab + b^2$ \quad
  b) $3 \cdot (2x + 3)$; \ $x(x - 5)$
  \diagnose{$3x + 4$ zeigt, dass nur der erste Summand multipliziert wird.
  $-2a - 10$ zeigt Vorzeichenfehler. $(a + b)^2 = a^2 + b^2$ ist der häufigste
  Fehler überhaupt und zeigt, dass Potenzieren über die Addition «verteilt»
  wird.}
\end{solution}
\raum{4cm}

% ---------------------------------------------------------------------------
% 6 Gleichungen
% ---------------------------------------------------------------------------
\aufgabe{6.5}
\begin{parts}
  \part Lösen Sie die Gleichungen.
  \begin{multicols}{3}
    $3 \cdot (x - 2) = 2x + 1$ \columnbreak

    $\dfrac{x}{2} + \dfrac{x}{3} = 10$ \columnbreak

    $\dfrac{x + 1}{4} = \dfrac{x - 2}{3}$
  \end{multicols}
  \part Prüfen Sie, \textbf{ohne die Gleichung zu lösen}, ob $x = 2$ eine Lösung
        der folgenden Gleichung ist.
        \[ 5x - 3 = 2 \cdot (x + 1) + 1 \]
\end{parts}
\begin{solution}
  a) $x = 7$; \ $x = 12$; \ $x = 11$ \qquad
  b) Ja: Für $x = 2$ ergeben beide Seiten $7$.
  \diagnose{Teil b) prüft, ob verstanden wird, was eine Lösung ist, unabhängig
  vom Lösungsverfahren. Wer b) nur durch Lösen der Gleichung beantworten kann,
  verwendet ein Verfahren ohne Begriffsverständnis.}
\end{solution}
\raum{4.5cm}

% ---------------------------------------------------------------------------
% 7 Funktionen und Graphen
% ---------------------------------------------------------------------------
\aufgabe{7.5}
\begin{parts}
  \part Geben Sie für die Gerade im linken Koordinatensystem die Steigung, den
        $y$-Achsenabschnitt und die Funktionsgleichung an.
  \part Zeichnen Sie den Graphen von $y = -\frac{1}{2}x + 3$ ins rechte
        Koordinatensystem.
\end{parts}

\begin{center}
\newcommand{\koord}[4]{% xmin, xmax, ymin, ymax
  \draw[gray!35, thin] (#1,#3) grid (#2,#4);
  \draw[->, thick] (#1-0.4,0) -- (#2+0.5,0) node[right] {\small $x$};
  \draw[->, thick] (0,#3-0.4) -- (0,#4+0.5) node[above] {\small $y$};
  \foreach \t in {#1,...,#2} {\ifnum\t=0\else
    \draw (\t,0.1) -- (\t,-0.1) node[below, font=\tiny] {$\t$};\fi}
  \foreach \t in {#3,...,#4} {\ifnum\t=0\else
    \draw (0.1,\t) -- (-0.1,\t) node[left, font=\tiny] {$\t$};\fi}}
\begin{tikzpicture}[scale=0.5]
  \koord{-3}{4}{-3}{7}
  \draw[very thick, domain=-2:3] plot (\x, {2*\x + 1});
  \fill (0,1) circle (4pt) (2,5) circle (4pt);
\end{tikzpicture}
\hspace{1.5cm}
\begin{tikzpicture}[scale=0.5]
  \koord{-2}{8}{-2}{6}
  \ifprintanswers
    \draw[very thick, domain=-2:8] plot (\x, {-0.5*\x + 3});
  \fi
\end{tikzpicture}
\end{center}
\begin{parts}
  \setcounter{partno}{2}
  \part Liegt der Punkt $\pkt{4}{7}$ auf der Geraden $y = 2x - 1$? Begründen Sie.
  \part Die Werte in der Tabelle gehören zu einer linearen Funktion.
        Bestimmen Sie die Funktionsgleichung.
  \quad
  \begin{tabular}{|c|c|c|c|}
    \hline $x$ & $0$ & $1$ & $2$  \\ \hline
           $y$ & $4$ & $1$ & $-2$ \\ \hline
  \end{tabular}
\end{parts}
\begin{solution}
  a) Steigung $2$, $y$-Achsenabschnitt $1$, $y = 2x + 1$ \quad
  b) Gerade durch $\pkt{0}{3}$ und $\pkt{6}{0}$ (siehe Grafik) \quad
  c) Ja, denn $2 \cdot 4 - 1 = 7$. \quad
  d) $y = -3x + 4$
  \diagnose{Die Aufgabe prüft alle Wechsel zwischen Graph, Gleichung und
  Tabelle. Typische Fehler: Steigung als $\frac{\Delta x}{\Delta y}$ berechnet
  ($\frac{1}{2}$), negative Steigung steigend gezeichnet, $y$-Achsenabschnitt
  und Nullstelle verwechselt.}
\end{solution}
\raum{3.5cm}

% ---------------------------------------------------------------------------
% 8 Geometrie und Grössen
% ---------------------------------------------------------------------------
\aufgabe{8.5}
\begin{parts}
  \part Ein rechtwinkliges Dreieck hat die Katheten $6$\,cm und $8$\,cm.
        Berechnen Sie die Länge der Hypotenuse.
  \part In einem rechtwinkligen Dreieck ist die Hypotenuse $13$\,cm und eine
        Kathete $5$\,cm lang. Berechnen Sie die Länge der anderen Kathete.
  \part Eine $10$\,m lange Leiter lehnt an einer senkrechten Wand. Ihr Fuss
        steht $6$\,m von der Wand entfernt. In welcher Höhe berührt die Leiter
        die Wand?
  \part Ist ein Dreieck mit den Seitenlängen $7$\,cm, $24$\,cm und $25$\,cm
        rechtwinklig? Begründen Sie.
\end{parts}
\begin{solution}
  a) $10$\,cm \quad b) $12$\,cm \quad c) $8$\,m \quad
  d) Ja, denn $7^2 + 24^2 = 49 + 576 = 625 = 25^2$.
  \diagnose{$\sqrt{194}$ bei b) zeigt, dass $a^2 + b^2 = c^2$ mechanisch
  angewendet wird, ohne die Hypotenuse zu identifizieren. c) verlangt, das
  rechtwinklige Dreieck in einer Situation zu erkennen. d) prüft die Umkehrung
  des Satzes von Pythagoras.}
\end{solution}
\raum{5cm}

% ---------------------------------------------------------------------------
% 9 Sprache ↔ Mathematik
% ---------------------------------------------------------------------------
\aufgabe{9.5}
\begin{parts}
  \part Schreiben Sie als Term: \quad «$5$ weniger als das Doppelte von $x$»
        \quad und \quad «um $3$ grösser als die Hälfte von $y$»
  \part An einer Schule gibt es sechsmal so viele Schülerinnen und Schüler wie
        Lehrpersonen. Schreiben Sie diese Aussage als Gleichung. Verwenden Sie
        $S$ für die Anzahl Schülerinnen und Schüler und $L$ für die Anzahl
        Lehrpersonen.
  \part Beschreiben Sie die Terme in Worten: \quad $3 \cdot (x + 4)$ \quad und \quad
        $3x + 4$
  \part Verdoppelt man eine Zahl und addiert $7$, so erhält man $31$. Stellen
        Sie eine Gleichung auf und bestimmen Sie die Zahl.
\end{parts}
\begin{solution}
  a) $2x - 5$; \ $\frac{y}{2} + 3$ \quad
  b) $S = 6L$ \quad
  c) «das Dreifache der Summe von $x$ und $4$»; «die Summe aus dem Dreifachen
  von $x$ und $4$» \quad
  d) $2x + 7 = 31$, $x = 12$
  \diagnose{b) ist die klassische «Studierende-und-Professoren»-Aufgabe: $6S = L$
  entsteht durch Wort-für-Wort-Übersetzung, ohne die Beziehung zu erfassen.
  $5 - 2x$ bei a) zeigt die Umkehrung bei «weniger als». c) ist die
  Gegenrichtung zu 9.3.}
\end{solution}
\raum{4.5cm}

% ===========================================================================
\newpage
\teil{D}{Vertiefung}
% ===========================================================================

% ---------------------------------------------------------------------------
% 1 Zahlen und Stellenwert
% ---------------------------------------------------------------------------
\aufgabe{1.6}
\begin{parts}
  \part Geben Sie für jede Zahl alle Zahlenmengen ($\mathbb{N}$, $\mathbb{Z}$,
        $\mathbb{Q}$, $\mathbb{R}$) an, zu denen sie gehört.
        \[
          -3 \qquad 0.75 \qquad \sqrt{2} \qquad \frac{12}{4} \qquad 0.333\ldots
        \]
  \part Kann das Ergebnis stimmen? Entscheiden Sie \textbf{ohne nachzurechnen}
        und begründen Sie.
        \begin{itemize}
          \item Lena rechnet: $0.6 \geteilt 0.03 = 2$.
          \item Ein Schüler berechnet das Volumen eines Schulzimmers und erhält
                $18$\,m$^3$.
        \end{itemize}
\end{parts}
\begin{solution}
  a) $-3$: $\mathbb{Z}, \mathbb{Q}, \mathbb{R}$; \
  $0.75$: $\mathbb{Q}, \mathbb{R}$; \
  $\sqrt{2}$: $\mathbb{R}$; \
  $\frac{12}{4}$: $\mathbb{N}, \mathbb{Z}, \mathbb{Q}, \mathbb{R}$; \
  $0.333\ldots$: $\mathbb{Q}, \mathbb{R}$ \\
  b) Beides ist nicht plausibel: Division durch eine Zahl kleiner als $1$
  vergrössert (richtig: $20$). Ein Schulzimmer hat etwa $150$ bis
  $200$\,m$^3$.
  \diagnose{Wird $\frac{12}{4}$ nur $\mathbb{Q}$ zugeordnet, so wird die Zahl
  nach ihrer Schreibweise statt nach ihrem Wert beurteilt. Die Zahl $0$ wurde
  bewusst vermieden, da Lehrmittel uneinheitlich sind, ob $0 \in \mathbb{N}$.}
\end{solution}
\raum{4cm}

% ---------------------------------------------------------------------------
% 2 Operationen und Rechenreihenfolge
% ---------------------------------------------------------------------------
\aufgabe{2.6}
\begin{parts}
  \part Welche Terme sind für alle Werte von $x$ gleichwertig mit $2 \cdot (x - 3)$?
        Kreuzen Sie an.

        \wahl{$2x - 3$} \quad \wahlr{$2x - 6$} \quad \wahlr{$(-2) \cdot (3 - x)$}
        \quad \wahlr{$(x - 3) \cdot 2$} \quad \wahl{$2 \cdot (3 - x)$}
  \part Gilt $(a - b)^2 = (b - a)^2$ für alle Zahlen $a$ und $b$? Begründen Sie.
\end{parts}
\begin{solution}
  a) $2x - 6$, \ $(-2) \cdot (3 - x)$, \ $(x - 3) \cdot 2$ \qquad
  b) Ja, denn $b - a = -(a - b)$, und beim Quadrieren fällt das Vorzeichen weg.
  \diagnose{Die Aufgabe prüft, ob Gleichwertigkeit an der Struktur erkannt wird.
  Wer nur einzelne Werte einsetzt, kann b) nicht allgemein begründen.}
\end{solution}
\raum{3cm}

% ---------------------------------------------------------------------------
% 3 Brüche
% ---------------------------------------------------------------------------
\aufgabe{3.6}
\begin{parts}
  \part Fassen Sie zusammen: \quad $\dfrac{1}{x} + \dfrac{1}{x + 1}$
  \part Kürzen Sie, falls möglich:
        \quad $\dfrac{3x + 6}{x + 2}$ \qquad $\dfrac{x^2 - 9}{x - 3}$
        \qquad $\dfrac{x + 5}{5}$
\end{parts}
\begin{solution}
  a) $\frac{2x + 1}{x(x + 1)}$ \quad
  b) $3$; \ $x + 3$; \ nicht kürzbar
  \diagnose{$\frac{x + 5}{5} = x + 1$ oder $= x$ zeigt «Kürzen aus Summen». Wer
  bei den ersten beiden Brüchen nicht faktorisiert, erkennt die Struktur nicht.
  Die Angabe der Definitionslücken ($x \neq -2$, $x \neq 3$) wird auf diesem
  Niveau nicht verlangt.}
\end{solution}
\raum{3.5cm}

% ---------------------------------------------------------------------------
% 4 Proportionalität und Prozent
% ---------------------------------------------------------------------------
\aufgabe{4.6}
\begin{parts}
  \part Ein Preis wird zuerst um $20\,\%$ erhöht und danach um $20\,\%$
        gesenkt. Ist er danach gleich hoch wie zu Beginn? Begründen Sie.
  \part Nach einer Erhöhung um $25\,\%$ kostet ein Artikel CHF 150. Wie viel
        kostete er vorher?
  \part Eine Bevölkerung von $P$ Personen wächst jährlich um $3\,\%$. Geben Sie
        einen Term für die Bevölkerung nach $n$ Jahren an.
\end{parts}
\begin{solution}
  a) Nein: $1.2 \cdot 0.8 = 0.96$, der Preis ist $4\,\%$ tiefer. \quad
  b) $150 \geteilt 1.25 = 120$, also CHF 120 \quad
  c) $P \cdot 1.03^n$
  \diagnose{CHF 112.50 bei b) zeigt, dass $25\,\%$ vom neuen statt vom alten
  Wert abgezogen werden. $P + 0.03n$ oder $P \cdot 0.03^n$ bei c) zeigen, dass
  der Wachstumsfaktor nicht verstanden ist.}
\end{solution}
\raum{4cm}

% ---------------------------------------------------------------------------
% 5 Terme
% ---------------------------------------------------------------------------
\aufgabe{5.6}
\begin{parts}
  \part Faktorisieren Sie so weit wie möglich.
  \begin{multicols}{4}
    $x^2 - 9$ \columnbreak

    $x^2 + 6x + 9$ \columnbreak

    $4a^2b - 6ab^2$ \columnbreak

    $ax + ay - bx - by$
  \end{multicols}
  \part Ergänzen Sie: \quad $x^2 + \kaestchen\,x + 25 = (x + \kaestchen)^2$
\end{parts}
\begin{solution}
  a) $(x - 3)(x + 3)$; \ $(x + 3)^2$; \ $2ab(2a - 3b)$; \ $(a - b)(x + y)$
  \quad b) $10$ und $5$
  \diagnose{Faktorisieren verlangt, die binomischen Formeln «rückwärts» zu
  lesen. Wer 5.5 lösen kann, hier aber scheitert, kennt die Formeln nur in eine
  Richtung.}
\end{solution}
\raum{3.5cm}

% ---------------------------------------------------------------------------
% 6 Gleichungen
% ---------------------------------------------------------------------------
\aufgabe{6.6}
\begin{parts}
  \part Lösen Sie nach der angegebenen Variable auf.
  \begin{multicols}{3}
    $A = \frac{1}{2} \cdot g \cdot h$ \ nach $h$ \columnbreak

    $v = \dfrac{s}{t}$ \ nach $t$ \columnbreak

    $F = \frac{9}{5} \cdot C + 32$ \ nach $C$
  \end{multicols}
  \part Lösen Sie nach $x$ auf ($a \neq 3$): \quad $ax = 3x + 5$
  \part Bestimmen Sie die Lösungsmenge.
        \quad $2 \cdot (x + 1) = 2x + 2$ \qquad $2 \cdot (x + 1) = 2x + 3$
\end{parts}
\begin{solution}
  a) $h = \frac{2A}{g}$; \ $t = \frac{s}{v}$; \ $C = \frac{5}{9} \cdot (F - 32)$
  \quad b) $x = \frac{5}{a - 3}$ \quad
  c) $\mathbb{L} = \mathbb{R}$; \ $\mathbb{L} = \{\,\}$
  \diagnose{Bei b) muss $x$ ausgeklammert werden ($x(a - 3) = 5$). Wer nur
  Gleichungen mit Zahlen lösen kann, kommt hier nicht weiter. Bei c) zeigt
  die Reaktion «$0 = 1$, also habe ich mich verrechnet», dass Lösungsmengen
  ohne genau eine Lösung unbekannt sind.}
\end{solution}
\raum{4.5cm}

% ---------------------------------------------------------------------------
% 7 Funktionen und Graphen
% ---------------------------------------------------------------------------
\aufgabe{7.6}
\begin{parts}
  \part Was haben alle Geraden mit der Gleichung $y = mx + 2$ gemeinsam? Was
        haben alle Geraden mit der Gleichung $y = 3x + q$ gemeinsam?
  \part Berechnen Sie den Schnittpunkt der Geraden $y = 2x + 1$ und
        $y = -x + 7$.
  \part Für welchen Wert von $m$ verläuft die Gerade $y = mx + 2$ durch den
        Punkt $\pkt{4}{0}$?
\end{parts}
\begin{solution}
  a) Alle verlaufen durch $\pkt{0}{2}$; alle sind parallel (Steigung $3$).
  \quad b) $2x + 1 = -x + 7 \Rightarrow x = 2$, also $S\pkt{2}{5}$
  \quad c) $0 = 4m + 2 \Rightarrow m = -\frac{1}{2}$
  \diagnose{Die Aufgabe prüft, ob ein Buchstabe als Parameter (fest, aber
  beliebig) und nicht als Variable verstanden wird, und ob ein Schnittpunkt als
  Lösung einer Gleichung erkannt wird.}
\end{solution}
\raum{4cm}

% ---------------------------------------------------------------------------
% 8 Geometrie und Grössen
% ---------------------------------------------------------------------------
\aufgabe{8.6}
\begin{parts}
  \part Die Geraden $AA'$ und $BB'$ sind parallel. Berechnen Sie die Längen
        $\overline{SB'}$ und $\overline{BB'}$. (Skizze nicht massstäblich)

  \begin{center}
  \begin{tikzpicture}[scale=0.8]
    \coordinate (S) at (0,0);
    \coordinate (A) at (4,0);   \coordinate (B) at (6,0);
    \coordinate (A2) at (3.2,2.4); \coordinate (B2) at (4.8,3.6);
    \draw[thick] (S) -- (7,0);
    \draw[thick] (S) -- (5.6,4.2);
    \draw[thick] (A) -- (A2);
    \draw[thick] (B) -- (B2);
    \foreach \p/\pos in {S/left, A/below, B/below, A2/above left, B2/above left}
      \fill (\p) circle (2pt);
    \node[left]        at (S)  {$S$};
    \node[below]       at (A)  {$A$};
    \node[below]       at (B)  {$B$};
    \node[above left]  at (A2) {$A'$};
    \node[above left]  at (B2) {$B'$};
    \node[below, font=\small] at (2,0)       {$4$\,cm};
    \node[below, font=\small] at (5,0)       {$2$\,cm};
    \node[above left, font=\small] at (1.6,1.2) {$5$\,cm};
    \node[right, font=\small] at (3.6,1.2)   {$3$\,cm};
  \end{tikzpicture}
  \end{center}
  \part Eine $1.80$\,m grosse Person wirft einen $1.20$\,m langen Schatten. Ein
        Baum wirft gleichzeitig einen $8$\,m langen Schatten. Wie hoch ist der
        Baum?
  \part Ein Quadrat mit der Seitenlänge $a$ wird mit dem Faktor $k$ vergrössert.
        Geben Sie den Umfang und den Flächeninhalt des neuen Quadrats an.
\end{parts}
\begin{solution}
  a) $\overline{SB} = 6$\,cm, also $\overline{SB'} = 5 \cdot \frac{6}{4} =
  7.5$\,cm und $\overline{BB'} = 3 \cdot \frac{6}{4} = 4.5$\,cm \quad
  b) $12$\,m \quad
  c) Umfang $4ka$, Flächeninhalt $k^2a^2$
  \diagnose{Bei a) zeigen Ergebnisse wie $\overline{BB'} = 1.5$\,cm oder
  $\overline{SB'} = 2.5$\,cm, dass mit der Teilstrecke $\overline{AB}$ statt
  mit $\overline{SB}$ gerechnet wird. Additive Ansätze ($3 + 2 = 5$) zeigen
  fehlendes Verständnis der Ähnlichkeit. $ka^2$ bei c) zeigt, dass die Fläche
  wie die Länge skaliert wird.}
\end{solution}
\raum{4cm}

% ---------------------------------------------------------------------------
% 9 Sprache ↔ Mathematik
% ---------------------------------------------------------------------------
\aufgabe{9.6}
\begin{parts}
  \part Handyabo A kostet CHF 20 pro Monat plus CHF 0.10 pro Gesprächsminute.
        Abo B kostet pauschal CHF 35 pro Monat. Stellen Sie für beide Abos
        einen Term für die Monatskosten bei $x$ Minuten auf. Ab wie vielen
        Minuten ist Abo B günstiger?
  \part Für einen Schulausflug fahren $200$ Personen mit Bussen, die je
        $45$ Plätze haben. Wie viele Busse braucht es?
\end{parts}
\begin{solution}
  a) A: $20 + 0.1x$; B: $35$. Aus $20 + 0.1x = 35$ folgt $x = 150$. Abo B ist
  ab mehr als $150$ Minuten günstiger. \quad
  b) $200 \geteilt 45 \approx 4.4$, also braucht es $5$ Busse.
  \diagnose{Die Antwort «$150$» ohne Aussage, was darunter und darüber gilt,
  zeigt, dass die Lösung nicht im Kontext gedeutet wird. Bei b) sind $4.44$
  oder $4$ Busse die klassischen Zeichen dafür, dass ein Rechenergebnis nicht
  auf die Situation zurückbezogen wird.}
\end{solution}
\raum{4.5cm}


\end{questions}

% ===========================================================================
% Auswertungsraster (für die Lehrperson)
% ===========================================================================
\newpage
\section*{Auswertung \normalsize(wird von der Lehrperson ausgefüllt)}

Pro Aufgabe: \quad \textbf{2} = sicher gelöst \quad \textbf{1} = teilweise
gelöst \quad \textbf{0} = nicht gelöst. \\
Die \emph{Bruchstelle} ist das tiefste Niveau eines Strangs, auf dem die
Aufgaben nicht mehr sicher gelöst werden.

\begin{center}
\small
\renewcommand{\arraystretch}{2.2}
\newcommand{\feld}[1]{#1\,\framebox[0.6cm]{\rule{0pt}{0.9em}}}
\begin{tabular}{|>{\raggedright\arraybackslash}p{3.5cm}|c|c|c|c|p{1.9cm}|}
  \hline
  \textbf{Strang}
    & \textbf{A}\,\footnotesize Grundlagen & \textbf{B}\,\footnotesize Numerisch
    & \textbf{C}\,\footnotesize Symbol. & \textbf{D}\,\footnotesize Formal
    & \footnotesize\textbf{Bruchstelle} \\ \hline
  1 Zahlen und Stellenwert          & \feld{1.1}\ \feld{1.2} & \feld{1.3}\ \feld{1.4} & \feld{1.5} & \feld{1.6} & \\ \hline
  2 Operationen und Rechenreihenfolge & \feld{2.1}\ \feld{2.2} & \feld{2.3}\ \feld{2.4} & \feld{2.5} & \feld{2.6} & \\ \hline
  3 Brüche                          & \feld{3.1}\ \feld{3.2} & \feld{3.3}\ \feld{3.4} & \feld{3.5} & \feld{3.6} & \\ \hline
  4 Proportionalität und Prozent    & \feld{4.1}\ \feld{4.2} & \feld{4.3}\ \feld{4.4} & \feld{4.5} & \feld{4.6} & \\ \hline
  5 Terme                           & \feld{5.1}\ \feld{5.2} & \feld{5.3}\ \feld{5.4} & \feld{5.5} & \feld{5.6} & \\ \hline
  6 Gleichungen                     & \feld{6.1}\ \feld{6.2} & \feld{6.3}\ \feld{6.4} & \feld{6.5} & \feld{6.6} & \\ \hline
  7 Funktionen und Graphen          & \feld{7.1}\ \feld{7.2} & \feld{7.3}\ \feld{7.4} & \feld{7.5} & \feld{7.6} & \\ \hline
  8 Geometrie und Grössen           & \feld{8.1}\ \feld{8.2} & \feld{8.3}\ \feld{8.4} & \feld{8.5} & \feld{8.6} & \\ \hline
  9 Sprache $\leftrightarrow$ Mathematik & \feld{9.1}\ \feld{9.2} & \feld{9.3}\ \feld{9.4} & \feld{9.5} & \feld{9.6} & \\ \hline
\end{tabular}
\end{center}

\textbf{Bemerkungen:}

\vspace{4cm}

\end{document}
"
% Jeweils zweimal kompilieren (Seitenzahlen).
%
% Aufbau: 9 Stränge × 4 Niveaus
%   A Grundlagen  (Aufgaben x.1, x.2)
%   B Numerisch   (Aufgaben x.3, x.4)
%   C Symbolisch  (Aufgabe  x.5)
%   D Formal      (Aufgabe  x.6)
%%%%%%%%%%%%%%%%%%%%%%%%%%%%%%%%%%%%%%%%%%%%%%%%%%%%%%%%%%%%%%%%%%%%%%%%%%%%%%

\ifdefined\mitloesungen \printanswers \fi

\usepackage[T1]{fontenc}
\usepackage{lmodern}
\usepackage[nswissgerman]{babel}
\usepackage[hidelinks]{hyperref}
\usepackage{amssymb, amsmath, setspace, multicol, array, needspace}
\usepackage{tikz}
\usepackage[margin=2cm]{geometry}

\usepackage{color}
\colorgrids
\definecolor{GridColor}{gray}{0.7}
\setlength{\gridsize}{4mm}

\singlespacing
\parindent 0ex

% ---------------------------------------------------------------------------
% Kopfzeilen und Angaben
% ---------------------------------------------------------------------------
\newcommand{\titel}{Standortbestimmung Mathematik}
\newcommand{\schuljahr}{Schuljahr 2026/27}
\newcommand{\zeit}{90 Minuten (Teile A und B)}
\newcommand{\hilfsmittel}{keine (kein Taschenrechner)}

% ---------------------------------------------------------------------------
% Hilfsbefehle
% ---------------------------------------------------------------------------
% Aufgabe mit fester Nummer (Strang.Nummer) statt fortlaufender Zählung
\newcommand{\aufgabe}[1]{%
  \needspace{8\baselineskip}%
  \renewcommand{\questionlabel}{\textbf{#1}}%
  \question}

% Teil-Überschrift über die volle Breite
\newcommand{\teil}[2]{\uplevel{\vspace{1ex}\textbf{\large Teil #1: #2}\par\vspace{1ex}}}

% Divisionszeichen (Doppelpunkt) mit korrektem Abstand
\newcommand{\geteilt}{\mathbin{:}}
% Punkt mit Koordinaten (x | y)
\newcommand{\pkt}[2]{(#1\,|\,#2)}
% Leeres Kästchen für Lückenaufgaben
\newcommand{\kaestchen}{\fbox{\rule{0pt}{1.1em}\hspace{1.6em}}}
% Platz für die Antwort (Raster), nur in der Version ohne Lösungen
\newcommand{\raum}[1]{\ifprintanswers\else\fillwithgrid{#1}\fi}
% Ankreuzfelder: \wahl = falsche, \wahlr = richtige Option
\newcommand{\wahl}[1]{\mbox{$\Box$\,#1}}
\newcommand{\wahlr}[1]{\mbox{\ifprintanswers$\blacksquare$\,\textbf{#1}\else$\Box$\,#1\fi}}
% Mehr Zeilenabstand in Aufgabenteilen mit Brüchen
\newcommand{\bruchabstand}{\setlength{\itemsep}{2.5ex}}
% Interpretationshinweis in den Lösungen
\newcommand{\diagnose}[1]{\par\smallskip\textit{Interpretation:} #1}

% ---------------------------------------------------------------------------
% Deutsche Bezeichnungen für die exam-Klasse
% ---------------------------------------------------------------------------
\renewcommand{\solutiontitle}{\noindent\textbf{Lösung:}\enspace}
\checkboxchar{$\Box$}
\checkedchar{$\blacksquare$}
\CorrectChoiceEmphasis{\bfseries}
\pointsinmargin
\nopointsinmargin

\begin{document}

\pagestyle{headandfoot}
\firstpageheader{}{}{}
\runningheader{\titel\ifprintanswers\ – Lösungen\fi}{}{\schuljahr}
\runningheadrule
\footer{}{Seite \thepage\ von \numpages}{}

% ===========================================================================
% Titelseite
% ===========================================================================
\begin{center}
  {\LARGE\bfseries \titel\ifprintanswers\ – Lösungen\fi}\\[1ex]
  {\large Eintritt ins Gymnasium – \schuljahr}
\end{center}
\vspace{1em}

\begin{tabular}{@{}p{0.5\textwidth}l@{}}
  \textbf{Name:} \makebox[5.5cm]{\hrulefill} & \textbf{Klasse:} \makebox[3cm]{\hrulefill} \\[2ex]
  \textbf{Zeit:} \zeit                       & \textbf{Hilfsmittel:} \hilfsmittel
\end{tabular}

\vspace{1.5em}
\rule{\textwidth}{.4pt}

\subsection*{Hinweise}

\begin{itemize}
  \item Diese Standortbestimmung ist \textbf{keine Prüfung} und wird nicht benotet. Sie zeigt
        uns, wo Sie in Mathematik stehen, damit wir Sie gezielt unterstützen können.
  \item Die Aufgaben werden von Teil zu Teil schwieriger. Es ist normal, dass
        Sie nicht alle Aufgaben lösen können. Versuchen Sie trotzdem jede Aufgabe.
  \item Lösen Sie die Aufgaben \textbf{ohne Taschenrechner}.
  \item Schreiben Sie Ihre Überlegungen und Zwischenschritte auf. Ein Lösungsweg
        sagt uns oft mehr als das Ergebnis.
  \item Wenn Sie eine Aufgabe nicht lösen können, gehen Sie zur nächsten weiter.
\end{itemize}

\vfill
\newpage

% ===========================================================================
\begin{questions}
% ===========================================================================

\teil{A}{Grundlagen}

% ---------------------------------------------------------------------------
% 1 Zahlen und Stellenwert
% ---------------------------------------------------------------------------
\aufgabe{1.1}
Ordnen Sie die Zahlen der Grösse nach, beginnend mit der kleinsten:
\[
  0.5; \quad 0.45; \quad 0.405; \quad 0.54
\]
\begin{solution}
  $0.405 < 0.45 < 0.5 < 0.54$
  \diagnose{Wird $0.405$ oder $0.45$ grösser als $0.5$ eingeordnet, so gilt
  vermutlich die Vorstellung «mehr Ziffern bedeutet eine grössere Zahl».
  Der Stellenwert der Dezimalstellen ist nicht gesichert.}
\end{solution}
\raum{2.5cm}

\aufgabe{1.2}
\begin{parts}
  \part Berechnen Sie: \quad $3.7 \cdot 100$ \qquad $45 \geteilt 1000$
  \part Welche Zahl liegt genau in der Mitte zwischen $0.3$ und $0.4$?
\end{parts}
\begin{solution}
  a) $370$; \ $0.045$ \qquad b) $0.35$
  \diagnose{Ergebnisse wie $3.700$ oder $0.45$ deuten auf die auswendig gelernte
  Regel «Nullen anhängen bzw.\ wegstreichen» hin, ohne Verständnis des
  Stellenwerts. Bei b) zeigt die Antwort «gibt es nicht», dass die Dichte der
  Dezimalzahlen nicht verstanden ist.}
\end{solution}
\raum{3cm}

% ---------------------------------------------------------------------------
% 2 Operationen und Rechenreihenfolge
% ---------------------------------------------------------------------------
\aufgabe{2.1}
Ein Band von $3.6$\,m Länge wird in Stücke von $0.4$\,m Länge geschnitten.
Welche Rechnung passt zu der Frage, wie viele Stücke entstehen? Kreuzen Sie an.

\begin{oneparcheckboxes}
  \choice $3.6 \cdot 0.4$
  \CorrectChoice $3.6 \geteilt 0.4$
  \choice $0.4 \geteilt 3.6$
  \choice $3.6 - 0.4$
\end{oneparcheckboxes}
\begin{solution}
  $3.6 \geteilt 0.4 = 9$ Stücke.
  \diagnose{Wer die Division meidet, hat oft die Vorstellung, man dürfe nur durch
  ganze Zahlen teilen. Die Wahl von $0.4 \geteilt 3.6$ zeigt die Regel «immer die
  grössere durch die kleinere Zahl» in umgekehrter Form bzw.\ fehlendes Verständnis
  der Division als Aufteilen.}
\end{solution}

\aufgabe{2.2}
Entscheiden Sie \textbf{ohne zu rechnen}, ob das Ergebnis grösser oder kleiner
als $50$ ist.
\begin{parts}
  \part $50 \cdot 0.8$ \hfill \wahl{grösser als 50}\quad \wahlr{kleiner als 50}
  \part $50 \geteilt 0.8$ \hfill \wahlr{grösser als 50}\quad \wahl{kleiner als 50}
  \part $50 \geteilt 1.25$ \hfill \wahl{grösser als 50}\quad \wahlr{kleiner als 50}
\end{parts}
\begin{solution}
  a) kleiner ($40$) \quad b) grösser ($62.5$) \quad c) kleiner ($40$)
  \diagnose{Fehler bei a) und b) zeigen die Fehlvorstellung «Multiplizieren
  vergrössert, Dividieren verkleinert», die aus dem Rechnen mit natürlichen
  Zahlen stammt.}
\end{solution}

% ---------------------------------------------------------------------------
% 3 Brüche
% ---------------------------------------------------------------------------
\aufgabe{3.1}
Markieren Sie die Zahlen $\dfrac{3}{4}$ und $\dfrac{5}{4}$ auf dem Zahlenstrahl.

\begin{center}
\begin{tikzpicture}[x=6cm]
  \draw[->, thick] (-0.05,0) -- (2.15,0);
  \foreach \x in {0,1,2} {
    \draw[thick] (\x,0.15) -- (\x,-0.15) node[below] {$\x$};
  }
  \ifprintanswers
    \foreach \x/\l in {0.75/{\frac{3}{4}}, 1.25/{\frac{5}{4}}} {
      \draw[very thick] (\x,0.25) -- (\x,-0.25);
      \node[above] at (\x,0.25) {$\l$};
    }
  \fi
\end{tikzpicture}
\end{center}
\begin{solution}
  Markierungen bei $0.75$ und $1.25$ (siehe Zahlenstrahl).
  \diagnose{Wer $\frac{5}{4}$ nicht einträgt oder für unmöglich hält, versteht
  Brüche nur als «Teil eines Ganzen», nicht als Zahlen. Falsche Positionen von
  $\frac{3}{4}$ zeigen Schwierigkeiten beim Einteilen der Einheit.}
\end{solution}

\aufgabe{3.2}
\begin{parts}
  \part Welcher Bruch ist jeweils grösser? Begründen Sie kurz.
  \[
    \frac{2}{3} \text{ oder } \frac{3}{4} \qquad\qquad
    \frac{5}{8} \text{ oder } \frac{5}{6} \qquad\qquad
    \frac{7}{8} \text{ oder } \frac{8}{9}
  \]
  \part Ergänzen Sie: \quad $\dfrac{2}{3} = \dfrac{\kaestchen}{12}$
\end{parts}
\begin{solution}
  a) $\frac{3}{4}$; \ $\frac{5}{6}$; \ $\frac{8}{9}$ (z.\,B.\ weil $\frac{8}{9}$
  nur $\frac{1}{9}$ von $1$ entfernt ist, $\frac{7}{8}$ aber $\frac{1}{8}$)
  \qquad b) $8$
  \diagnose{Wird $\frac{5}{8}$ als grösser gewählt («8 ist grösser als 6»), so
  werden Zähler und Nenner getrennt als natürliche Zahlen betrachtet. Das letzte
  Paar unterscheidet auswendig gelernte Verfahren (gleichnamig machen) von
  Zahlverständnis (Abstand zu~1).}
\end{solution}
\raum{3.5cm}

% ---------------------------------------------------------------------------
% 4 Proportionalität und Prozent
% ---------------------------------------------------------------------------
\aufgabe{4.1}
3 Hefte kosten CHF 7.50. Wie viel kosten 12 Hefte? Wie viel kosten 5 Hefte?
\begin{solution}
  12 Hefte: CHF 30.00 \qquad 5 Hefte: CHF 12.50 (1 Heft kostet CHF 2.50)
  \diagnose{Gelingen die 12 Hefte, die 5 Hefte aber nicht, so wird nur
  vervielfacht, der Weg über den Einheitspreis (Dreisatz) fehlt.}
\end{solution}
\raum{3cm}

\aufgabe{4.2}
Für 4 Personen braucht ein Rezept 6 Eier. Wie viele Eier braucht man für
6 Personen?
\begin{solution}
  $9$ Eier
  \diagnose{Die Antwort $8$ («2 Personen mehr, also 2 Eier mehr») ist das
  klassische Zeichen für additives statt multiplikatives Denken.}
\end{solution}
\raum{3cm}

% ---------------------------------------------------------------------------
% 5 Terme
% ---------------------------------------------------------------------------
\aufgabe{5.1}
Aus Hölzchen werden Quadrate in einer Reihe gelegt.

\begin{center}
\begin{tikzpicture}[x=0.8cm, y=0.8cm, line width=1.6pt, line cap=round]
  \newcommand{\reihe}[2]{% #1 = x-Start, #2 = Anzahl Quadrate
    \foreach \i in {1,...,#2} {
      \draw (#1+\i-1,0) -- (#1+\i,0);
      \draw (#1+\i-1,1) -- (#1+\i,1);
      \draw (#1+\i,0) -- (#1+\i,1);
    }
    \draw (#1,0) -- (#1,1);
  }
  \reihe{0}{1}  \node at (0.5,-0.6) {\small 4 Hölzchen};
  \reihe{3}{2}  \node at (4,-0.6)   {\small 7 Hölzchen};
  \reihe{7}{3}  \node at (8.5,-0.6) {\small 10 Hölzchen};
\end{tikzpicture}
\end{center}
\begin{parts}
  \part Wie viele Hölzchen braucht man für 10 Quadrate?
  \part Wie viele Hölzchen braucht man für $n$ Quadrate?
\end{parts}
\begin{solution}
  a) $31$ \qquad b) $3n + 1$ (oder gleichwertig, z.\,B.\ $4 + 3 \cdot (n-1)$)
  \diagnose{Gelingt a), aber nicht b), so kann ein Muster an Beispielen
  fortgesetzt, aber nicht verallgemeinert werden. Antworten wie $4n$ zeigen, dass
  die gemeinsamen Hölzchen nicht beachtet werden.}
\end{solution}
\raum{3.5cm}

\aufgabe{5.2}
Welcher Ausdruck ist grösser: $n + 2$ oder $2n$? Begründen Sie.
\begin{solution}
  Das hängt von $n$ ab: Für $n = 2$ sind beide gleich, für $n < 2$ ist $n + 2$
  grösser, für $n > 2$ ist $2n$ grösser.
  \diagnose{Die Antwort «$2n$, weil verdoppeln mehr ist» oder die Prüfung mit
  nur einem Wert zeigen, dass die Variable als feste Zahl und nicht als
  Platzhalter für viele Zahlen verstanden wird.}
\end{solution}
\raum{3cm}

% ---------------------------------------------------------------------------
% 6 Gleichungen
% ---------------------------------------------------------------------------
\aufgabe{6.1}
Welche Zahl gehört ins Kästchen?
\begin{parts}
  \part $8 + 4 = \kaestchen + 5$
  \part $3 \cdot \kaestchen = 24 \geteilt 2$
\end{parts}
\begin{solution}
  a) $7$ \qquad b) $4$
  \diagnose{Die Antworten $12$ oder $17$ bei a) zeigen, dass das
  Gleichheitszeichen als Aufforderung «jetzt kommt das Ergebnis» gelesen wird
  und nicht als Beziehung zwischen zwei gleichwertigen Seiten.}
\end{solution}

\aufgabe{6.2}
Ist die Aussage wahr oder falsch? Begründen Sie, möglichst ohne die Ergebnisse
auszurechnen.
\begin{parts}
  \part $37 + 25 = 36 + 26$ \hfill \wahlr{wahr}\quad \wahl{falsch}
  \part $47 - 18 = 48 - 17$ \hfill \wahl{wahr}\quad \wahlr{falsch}
  \part $6 \cdot 15 = 3 \cdot 30$ \hfill \wahlr{wahr}\quad \wahl{falsch}
\end{parts}
\begin{solution}
  a) wahr \quad b) falsch ($29 \neq 31$) \quad c) wahr
  \diagnose{Begründungen wie «links eins weniger, rechts eins mehr» zeigen
  relationales Denken. Wer b) als wahr einstuft, überträgt die
  Ausgleichsregel der Addition unreflektiert auf die Subtraktion.}
\end{solution}
\raum{3cm}

% ---------------------------------------------------------------------------
% 7 Funktionen und Graphen
% ---------------------------------------------------------------------------
\aufgabe{7.1}
Geben Sie die Koordinaten der Punkte $A$, $B$ und $C$ an.

\begin{center}
\begin{tikzpicture}[scale=0.55]
  \draw[gray!40, thin] (-5,-5) grid (5,5);
  \draw[->, thick] (-5.4,0) -- (5.6,0) node[right] {$x$};
  \draw[->, thick] (0,-5.4) -- (0,5.6) node[above] {$y$};
  \foreach \t in {-4,-2,2,4} {
    \node[below, font=\scriptsize] at (\t,0) {$\t$};
    \node[left,  font=\scriptsize] at (0,\t) {$\t$};
  }
  \foreach \p/\x/\y/\pos in {A/3/-2/right, B/-2/3/left, C/0/4/right} {
    \fill (\x,\y) circle (3.5pt);
    \node[\pos, font=\small] at (\x,\y) {$\p$};
  }
\end{tikzpicture}
\end{center}
$A = $ \fillin[$\pkt{3}{-2}$][2.5cm] \qquad
$B = $ \fillin[$\pkt{-2}{3}$][2.5cm] \qquad
$C = $ \fillin[$\pkt{0}{4}$][2.5cm]
\begin{solution}
  $A\pkt{3}{-2}$, \ $B\pkt{-2}{3}$, \ $C\pkt{0}{4}$
  \diagnose{Vertauschte Koordinaten bei $A$ und $B$ zeigen, dass die
  Reihenfolge $(x\,|\,y)$ nicht gesichert ist. Fehler bei $C$ (z.\,B.\ $(4\,|\,0)$)
  deuten auf Unsicherheit bei Punkten auf den Achsen hin.}
\end{solution}

\aufgabe{7.2}
Die $y$-Werte in der Tabelle werden alle nach derselben Regel aus den
$x$-Werten berechnet. Ergänzen Sie die Tabelle sinngemäss.
\begin{center}
  \renewcommand{\arraystretch}{1.4}
  \begin{tabular}{|c|*{4}{>{\centering\arraybackslash}p{1.2cm}|}}
    \hline
    $x$ & $0$ & $1$ & $3$  & $6$ \\ \hline
    $y$ & $2$ & $5$ & $11$ & \ifprintanswers $\mathbf{20}$\fi \\ \hline
  \end{tabular}
\end{center}
\begin{solution}
  $y = 20$ (Zuordnung $y = 3x + 2$)
  \diagnose{Da die $x$-Werte nicht in gleichen Schritten wachsen, scheitert ein
  rein schrittweises Weiterzählen (z.\,B.\ «immer $+6$» ergibt $17$). Das zeigt,
  ob der Zusammenhang zwischen $x$ und $y$ erkannt wird.}
\end{solution}
\raum{2.5cm}

% ---------------------------------------------------------------------------
% 8 Geometrie und Grössen
% ---------------------------------------------------------------------------
\aufgabe{8.1}
Ergänzen Sie.
\begin{parts}
  \part $3.5\text{ m} = $ \fillin[$350$][2cm] cm
  \part $250\text{ g} = $ \fillin[$0.25$][2cm] kg
  \part $1\text{ m}^2 = $ \fillin[$10\,000$][2cm] cm$^2$
  \part $2\text{ h }15\text{ min} = $ \fillin[$2.25$][2cm] h
\end{parts}
\begin{solution}
  a) $350$ \quad b) $0.25$ \quad c) $10\,000$ \quad d) $2.25$
  \diagnose{Die Antwort $100$ bei c) zeigt, dass die Umrechnungszahl für
  Flächeneinheiten nicht verstanden ist. Die Antwort $2.15$ bei d) zeigt, dass
  Zeitangaben wie Dezimalzahlen behandelt werden.}
\end{solution}

\aufgabe{8.2}
Ein Rechteck ist $6$\,cm lang und $4$\,cm breit.
\begin{parts}
  \part Berechnen Sie den Umfang und den Flächeninhalt.
  \part Hat jedes Rechteck mit demselben Umfang auch denselben Flächeninhalt?
        Begründen Sie.
\end{parts}
\begin{solution}
  a) Umfang $20$\,cm, Flächeninhalt $24$\,cm$^2$
  \qquad b) Nein, z.\,B.\ hat ein Rechteck mit $7$\,cm und $3$\,cm ebenfalls den
  Umfang $20$\,cm, aber den Flächeninhalt $21$\,cm$^2$.
  \diagnose{Verwechslungen von Umfang und Fläche (auch bei den Einheiten) sowie
  die Antwort «ja» bei b) zeigen, dass die beiden Grössen nicht klar
  unterschieden werden.}
\end{solution}
\raum{3.5cm}

% ---------------------------------------------------------------------------
% 9 Sprache ↔ Mathematik
% ---------------------------------------------------------------------------
\aufgabe{9.1}
Schreiben Sie die Rechnung auf und berechnen Sie das Ergebnis.
\begin{parts}
  \part die Summe von $12$ und $8$
  \part die Differenz von $20$ und $5$
  \part das Produkt von $6$ und $7$
  \part der Quotient von $36$ und $4$
\end{parts}
\begin{solution}
  a) $12 + 8 = 20$ \quad b) $20 - 5 = 15$ \quad c) $6 \cdot 7 = 42$ \quad
  d) $36 \geteilt 4 = 9$
  \diagnose{Unbekannte Fachbegriffe zeigen sich hier direkt. Bei b) und d)
  auf die Reihenfolge in der aufgeschriebenen Rechnung achten ($5 - 20$,
  $4 \geteilt 36$).}
\end{solution}
\raum{3cm}

\aufgabe{9.2}
\begin{parts}
  \part Mia hat dreimal so viele Karten wie Leo. Leo hat 12 Karten. Wie viele
        Karten hat Mia?
  \part Nora hat 12 Karten. Das sind dreimal so viele wie Tim. Wie viele Karten
        hat Tim?
\end{parts}
\begin{solution}
  a) $36$ \qquad b) $4$
  \diagnose{Die Antwort $36$ bei b) zeigt eine Schlüsselwort-Strategie
  («mal» bedeutet multiplizieren), ohne die Beziehung zwischen den Grössen zu
  erfassen.}
\end{solution}
\raum{3cm}

% ===========================================================================
\newpage
\teil{B}{Rechnen mit Zahlen}
% ===========================================================================

% ---------------------------------------------------------------------------
% 1 Zahlen und Stellenwert
% ---------------------------------------------------------------------------
\aufgabe{1.3}
\begin{parts}
  \part Runden Sie.
  \begin{subparts}
    \subpart $3.4567$ auf zwei Dezimalstellen
    \subpart $0.0496$ auf zwei Dezimalstellen
    \subpart $28\,349$ auf Tausender
  \end{subparts}
  \part Welches Ergebnis ist plausibel? Kreuzen Sie an, ohne genau zu rechnen.
        \[
          0.48 \cdot 205 \approx \quad
          \text{\wahl{$10$}} \quad \text{\wahlr{$100$}} \quad
          \text{\wahl{$1000$}} \quad \text{\wahl{$10\,000$}}
        \]
\end{parts}
\begin{solution}
  a) $3.46$; \ $0.05$; \ $28\,000$ \qquad b) $100$ (genau: $98.4$)
  \diagnose{$3.45$ zeigt Abschneiden statt Runden. Antworten wie $0.5$ oder
  $0.04$ bei $0.0496$ zeigen, dass der Stellenwert beim Runden verloren geht.
  Fehler bei b) deuten auf fehlende Überschlagsrechnung hin.}
\end{solution}
\raum{3cm}

\aufgabe{1.4}
\begin{parts}
  \part Ordnen Sie der Grösse nach, beginnend mit der kleinsten Zahl:
        \quad $-3.5; \quad -3; \quad 0.5; \quad -4; \quad -0.5$
  \part Welche Zahl ist um $5$ grösser als $-3$?
\end{parts}
\begin{solution}
  a) $-4 < -3.5 < -3 < -0.5 < 0.5$ \qquad b) $2$
  \diagnose{Wird $-3.5$ grösser als $-3$ eingeordnet, so werden negative Zahlen
  nach ihrem Betrag geordnet. Antworten wie $-8$ oder $8$ bei b) zeigen
  Unsicherheit beim Bewegen auf dem Zahlenstrahl.}
\end{solution}
\raum{2.5cm}

% ---------------------------------------------------------------------------
% 2 Operationen und Rechenreihenfolge
% ---------------------------------------------------------------------------
\aufgabe{2.3}
Berechnen Sie.
\begin{multicols}{3}
\begin{parts}
  \part $-7 + 12$
  \part $-4 - 9$
  \part $(-6) \cdot (-3)$
  \part $(-15) \geteilt 3$
  \part $5 - (-8)$
\end{parts}
\end{multicols}
\begin{solution}
  a) $5$ \quad b) $-13$ \quad c) $18$ \quad d) $-5$ \quad e) $13$
  \diagnose{Typische Fehler sind $-5$ oder $13$ bei b) sowie $-3$ bei e). Sie
  zeigen, dass Vorzeichen und Rechenzeichen nicht sicher unterschieden werden.}
\end{solution}
\raum{2.5cm}

\aufgabe{2.4}
Berechnen Sie.
\begin{multicols}{3}
\begin{parts}
  \part $12 - 3 \cdot 4$
  \part $(12 - 3) \cdot 4$
  \part $20 \geteilt 4 \cdot 5$
  \part $2 \cdot 3^2$
  \part $18 - 6 \geteilt 2 + 1$
\end{parts}
\end{multicols}
\begin{solution}
  a) $0$ \quad b) $36$ \quad c) $25$ \quad d) $18$ \quad e) $16$
  \diagnose{Typische Fehler: c) $1$ (Multiplikation vor Division), d) $36$
  (zuerst multipliziert, dann quadriert), e) $7$ (von links nach rechts
  gerechnet) oder $14$ (Addition vor Subtraktion). Sie zeigen, welche Regel der
  Rechenreihenfolge fehlt.}
\end{solution}
\raum{3cm}

% ---------------------------------------------------------------------------
% 3 Brüche
% ---------------------------------------------------------------------------
\aufgabe{3.3}
Berechnen Sie und kürzen Sie das Ergebnis so weit wie möglich.
\begin{multicols}{3}
\begin{parts}\bruchabstand
  \part $\dfrac{2}{3} + \dfrac{1}{4}$
  \part $\dfrac{5}{6} - \dfrac{1}{2}$
  \part $\dfrac{3}{4} \cdot \dfrac{2}{9}$
  \part $\dfrac{2}{5} \geteilt \dfrac{4}{15}$
  \part $3 \cdot \dfrac{2}{5}$
\end{parts}
\end{multicols}
\begin{solution}
  a) $\frac{11}{12}$ \quad b) $\frac{1}{3}$ \quad c) $\frac{1}{6}$ \quad
  d) $\frac{3}{2}$ \quad e) $\frac{6}{5}$
  \diagnose{$\frac{3}{7}$ bei a) zeigt, dass Zähler und Nenner getrennt addiert
  werden. $\frac{6}{15}$ bei e) zeigt, dass Zähler und Nenner mit $3$
  multipliziert werden (Verwechslung mit dem Erweitern). Fehler nur bei d)
  deuten auf eine nicht verstandene Divisionsregel hin.}
\end{solution}
\raum{4cm}

\aufgabe{3.4}
Ergänzen Sie die Tabelle.
\begin{center}
  \newcommand{\lsg}[1]{\ifprintanswers$\mathbf{#1}$\fi}
  \newcommand{\zh}{\rule[-2.4ex]{0pt}{6.6ex}}% Zeilenhöhe für Brüche
  \begin{tabular}{|c|c|c|}
    \hline
    \zh\makebox[2.6cm]{\textbf{Bruch}} & \makebox[2.6cm]{\textbf{Dezimalzahl}}
      & \makebox[2.6cm]{\textbf{Prozent}} \\ \hline
    \zh$\dfrac{3}{4}$          & \lsg{0.75}           & \lsg{75\,\%} \\ \hline
    \zh\lsg{\dfrac{1}{20}}     & $0.05$               & \lsg{5\,\%} \\ \hline
    \zh\lsg{\dfrac{5}{4}}      & \lsg{1.25}           & $125\,\%$ \\ \hline
    \zh$\dfrac{1}{3}$          & \lsg{0.333\ldots}    & \lsg{33.3\,\%} \\ \hline
  \end{tabular}
\end{center}
\begin{solution}
  Siehe Tabelle. Für $\frac{1}{3}$ sind auch $0.\overline{3}$ und
  $33\frac{1}{3}\,\%$ korrekt.
  \diagnose{$0.05 = 50\,\%$ zeigt Unsicherheit beim Stellenwert. Wer bei
  $125\,\%$ keinen Bruch findet, hält Prozentangaben über $100\,\%$ oder Brüche
  grösser als $1$ für unmöglich.}
\end{solution}
\raum{2.5cm}

% ---------------------------------------------------------------------------
% 4 Proportionalität und Prozent
% ---------------------------------------------------------------------------
\aufgabe{4.3}
Ein Pullover kostet CHF 80.
\begin{parts}
  \part Berechnen Sie $15\,\%$ von CHF 80.
  \part Der Preis des Pullovers wird um $25\,\%$ reduziert. Wie viel kostet er
        jetzt?
\end{parts}
\begin{solution}
  a) CHF 12 \qquad b) CHF 60
  \diagnose{Wer bei b) nur den Rabatt (CHF 20) angibt, unterscheidet nicht
  zwischen Prozentwert und neuem Wert. Fehler bei a) zeigen, dass die
  Grundidee von «Prozent als Hundertstel» fehlt.}
\end{solution}
\raum{3cm}

\aufgabe{4.4}
Ein Preis steigt von CHF 40 auf CHF 50.
\begin{parts}
  \part Um wie viel Prozent ist der Preis gestiegen?
  \part Danach sinkt der Preis wieder von CHF 50 auf CHF 40. Um wie viel Prozent
        ist er gesunken?
\end{parts}
\begin{solution}
  a) $25\,\%$ \qquad b) $20\,\%$
  \diagnose{Die Antwort $25\,\%$ bei b) (oder $20\,\%$ bei a)) zeigt, dass
  nicht erkannt wird, dass sich der Prozentsatz auf den jeweiligen Grundwert
  bezieht. Dieser Fehler ist sehr häufig.}
\end{solution}
\raum{3cm}

% ---------------------------------------------------------------------------
% 5 Terme
% ---------------------------------------------------------------------------
\aufgabe{5.3}
Vereinfachen Sie so weit wie möglich.
\begin{multicols}{3}
\begin{parts}
  \part $3a + 5 + 2a - 1$
  \part $4x + 3y - x + 2y$
  \part $2a + 3b$
  \part $3x \cdot 4x$
  \part $a^2 + a + a^2$
\end{parts}
\end{multicols}
\begin{solution}
  a) $5a + 4$ \quad b) $3x + 5y$ \quad c) nicht vereinfachbar \quad
  d) $12x^2$ \quad e) $2a^2 + a$
  \diagnose{$5ab$ bei c) sowie $a^4$, $3a^2$ oder $2a^3$ bei e) zeigen, dass
  ungleichartige Terme zusammengefasst werden. Die Aussage «$2a + 3b$ ist keine
  Lösung» weist auf die Vorstellung hin, dass ein Term ein einzelnes Ergebnis
  haben muss.}
\end{solution}
\raum{3cm}

\aufgabe{5.4}
Ist die Aussage richtig oder falsch? Kreuzen Sie an.
\begin{parts}
  \part $2a + 3b = 5ab$     \hfill \wahl{richtig}\quad \wahlr{falsch}
  \part $x + x = x^2$       \hfill \wahl{richtig}\quad \wahlr{falsch}
  \part $3x \cdot 2 = 6x$   \hfill \wahlr{richtig}\quad \wahl{falsch}
  \part $a + a^2 = a^3$     \hfill \wahl{richtig}\quad \wahlr{falsch}
  \part $4y - y = 4$        \hfill \wahl{richtig}\quad \wahlr{falsch}
\end{parts}
\begin{solution}
  a) falsch \quad b) falsch \quad c) richtig \quad d) falsch \quad e) falsch
  \diagnose{Wiedererkennungsversion von 5.3. Wer hier richtig entscheidet, bei
  5.3 aber Fehler macht, kennt die Regeln, wendet sie aber nicht sicher an.
  Bei e) zeigt «richtig», dass $y$ als Einheit oder Abkürzung gelesen wird
  («4 Äpfel minus 1 Apfel»).}
\end{solution}

% ---------------------------------------------------------------------------
% 6 Gleichungen
% ---------------------------------------------------------------------------
\aufgabe{6.3}
Lösen Sie die Gleichungen.
\begin{multicols}{3}
\begin{parts}\bruchabstand
  \part $x + 7 = 3$
  \part $3x = 21$
  \part $4x - 5 = 11$
  \part $12 = 5x + 2$
  \part $\dfrac{x}{3} = 6$
\end{parts}
\end{multicols}
\begin{solution}
  a) $x = -4$ \quad b) $x = 7$ \quad c) $x = 4$ \quad d) $x = 2$ \quad
  e) $x = 18$
  \diagnose{Bei d) steht die Variable rechts, was manche Schülerinnen und
  Schüler verunsichert. $x = 2$ bei e) zeigt, dass Umkehroperationen
  verwechselt werden.}
\end{solution}
\raum{3.5cm}

\aufgabe{6.4}
Lösen Sie die Gleichungen.
\begin{multicols}{3}
\begin{parts}
  \part $5x + 3 = 2x + 15$
  \part $7 - 2x = 1$
  \part $3x = 0$
\end{parts}
\end{multicols}
\begin{solution}
  a) $x = 4$ \quad b) $x = 3$ \quad c) $x = 0$
  \diagnose{$x = -3$ bei b) zeigt Vorzeichenfehler beim Umformen. Bei c)
  antworten manche «keine Lösung» oder «$x$ ist nicht definiert», weil $0$
  nicht als Lösung akzeptiert wird.}
\end{solution}
\raum{3.5cm}

% ---------------------------------------------------------------------------
% 7 Funktionen und Graphen
% ---------------------------------------------------------------------------
\aufgabe{7.3}
Der Graph zeigt, wie weit eine Person von zu Hause entfernt ist.

\begin{center}
\begin{tikzpicture}[x=0.22cm, y=1.1cm]
  \draw[gray!35, thin, xstep=5, ystep=0.5] (0,0) grid (45,3.5);
  \draw[->, thick] (0,0) -- (47,0) node[right] {\small Zeit (min)};
  \draw[->, thick] (0,0) -- (0,3.9) node[above] {\small Entfernung (km)};
  \foreach \t in {0,10,20,30,40} \node[below, font=\scriptsize] at (\t,0) {$\t$};
  \foreach \s in {1,2,3} \node[left, font=\scriptsize] at (0,\s) {$\s$};
  \draw[very thick] (0,0) -- (20,3) -- (30,3) -- (40,0);
\end{tikzpicture}
\end{center}
\begin{parts}
  \part Wie weit ist die Person nach 20 Minuten von zu Hause entfernt?
  \part Was geschieht zwischen der 20.\ und der 30.\ Minute?
  \part In welchem Abschnitt bewegt sich die Person am schnellsten?
        Begründen Sie.
  \part Berechnen Sie die Geschwindigkeit im ersten Abschnitt in km/h.
\end{parts}
\begin{solution}
  a) $3$\,km \quad b) Die Person bleibt stehen (macht eine Pause). \quad
  c) Auf dem Rückweg (30.--40.\ Minute), weil der Graph dort am steilsten ist.
  \quad d) $3$\,km in $20$\,min, also $9$\,km/h.
  \diagnose{Wird bei c) der «höchste» Abschnitt gewählt, so wird die Höhe des
  Graphen statt seiner Steigung gelesen. Bei b) zeigt «sie geht geradeaus
  weiter», dass der Graph als Bild des Weges gedeutet wird.}
\end{solution}
\raum{3.5cm}

\aufgabe{7.4}
Eine Velofahrerin fährt über einen Hügel: zuerst bergauf, dann bergab.
Welcher Graph zeigt ihre Geschwindigkeit $v$ in Abhängigkeit von der Zeit $t$
am besten? Kreuzen Sie an.

\begin{center}
\newcommand{\achsen}{%
  \draw[->] (0,0) -- (2.4,0) node[below] {\scriptsize $t$};
  \draw[->] (0,0) -- (0,1.7) node[left]  {\scriptsize $v$};}
\begin{tikzpicture}[thick]
  \begin{scope}[xshift=0cm]
    \achsen \draw[very thick] (0,0.3) .. controls (0.8,1.6) and (1.4,1.6) .. (2.2,0.3);
    \node at (1.1,-0.5) {\wahl{Graph 1}};
  \end{scope}
  \begin{scope}[xshift=3.3cm]
    \achsen \draw[very thick] (0,1.0) .. controls (0.6,0.3) and (0.9,0.3) .. (1.1,0.35)
                                     .. controls (1.4,0.5) and (1.8,1.3) .. (2.2,1.5);
    \node at (1.1,-0.5) {\wahlr{Graph 2}};
  \end{scope}
  \begin{scope}[xshift=6.6cm]
    \achsen \draw[very thick] (0,1.4) -- (2.2,0.3);
    \node at (1.1,-0.5) {\wahl{Graph 3}};
  \end{scope}
  \begin{scope}[xshift=9.9cm]
    \achsen \draw[very thick] (0,0.3) -- (1.1,1.3) -- (2.2,1.3);
    \node at (1.1,-0.5) {\wahl{Graph 4}};
  \end{scope}
\end{tikzpicture}
\end{center}
\begin{solution}
  Graph 2: Bergauf nimmt die Geschwindigkeit ab, bergab nimmt sie zu.
  \diagnose{Die Wahl von Graph 1 zeigt die klassische Fehlvorstellung «Graph
  als Bild der Situation»: Der Graph sieht aus wie der Hügel, beschreibt aber
  nicht die Geschwindigkeit.}
\end{solution}

% ---------------------------------------------------------------------------
% 8 Geometrie und Grössen
% ---------------------------------------------------------------------------
\aufgabe{8.3}
Aus einem Rechteck wird an einer Ecke ein kleineres Rechteck herausgeschnitten.
Berechnen Sie den Flächeninhalt und den Umfang der Figur.

\begin{center}
\begin{tikzpicture}[scale=0.5]
  \draw[very thick, fill=gray!10] (0,0) -- (8,0) -- (8,4) -- (5,4) -- (5,6) -- (0,6) -- cycle;
  \draw[dashed, gray] (5,4) -- (5,6) (5,6) -- (8,6) -- (8,4);
  \node[below] at (4,0) {\small $8$\,cm};
  \node[left]  at (0,3) {\small $6$\,cm};
  \node[above] at (6.5,6) {\small $3$\,cm};
  \node[right] at (8,5) {\small $2$\,cm};
\end{tikzpicture}
\end{center}
\begin{solution}
  Flächeninhalt: $8 \cdot 6 - 3 \cdot 2 = 42$\,cm$^2$ \qquad
  Umfang: $28$\,cm (gleich wie beim ganzen Rechteck)
  \diagnose{Wer beim Umfang etwas abzieht (z.\,B.\ $28 - 10 = 18$\,cm),
  überträgt die Überlegung zur Fläche auf den Umfang. Das zeigt, dass Umfang
  und Fläche nicht als verschiedene Grössen verstanden werden.}
\end{solution}
\raum{3.5cm}

\aufgabe{8.4}
\begin{parts}
  \part In einem Dreieck sind zwei Winkel $47^\circ$ und $68^\circ$ gross. Wie
        gross ist der dritte Winkel?
  \part Ein Aquarium ist $50$\,cm lang, $40$\,cm breit und $30$\,cm hoch. Wie
        viele Liter Wasser fasst es?
\end{parts}
\begin{solution}
  a) $65^\circ$ \qquad b) $60\,000$\,cm$^3 = 60$\,dm$^3 = 60$\,Liter
  \diagnose{Ergebnisse wie $600$\,L oder $6$\,L bei b) zeigen, dass der
  Zusammenhang $1\,\text{dm}^3 = 1$\,L bzw.\ die Umrechnung von
  Volumeneinheiten nicht gesichert ist.}
\end{solution}
\raum{3cm}

% ---------------------------------------------------------------------------
% 9 Sprache ↔ Mathematik
% ---------------------------------------------------------------------------
\aufgabe{9.3}
Schreiben Sie als \textbf{einen} Rechenausdruck und berechnen Sie das Ergebnis.
\begin{parts}
  \part das Dreifache der Differenz von $12$ und $5$
  \part die Differenz aus dem Dreifachen von $12$ und $5$
  \part Subtrahieren Sie $4$ von $10$.
  \part die Hälfte der Summe von $9$ und $15$
\end{parts}
\begin{solution}
  a) $3 \cdot (12 - 5) = 21$ \quad b) $3 \cdot 12 - 5 = 31$ \quad
  c) $10 - 4 = 6$ \quad d) $(9 + 15) \geteilt 2 = 12$
  \diagnose{a) und b) enthalten dieselben Wörter in anderer Reihenfolge und
  ergeben verschiedene Strukturen. Fehlende Klammern bei a) oder d) zeigen,
  dass die Struktur des Satzes nicht in die Rechnung übertragen wird. $4 - 10$
  bei c) ist die typische Umkehrung der Reihenfolge.}
\end{solution}
\raum{3.5cm}

\aufgabe{9.4}
Eine Klasse mit 24 Schülerinnen und Schülern macht einen Ausflug. Die Bahnfahrt
kostet pro Person CHF 18.50, der Museumseintritt pro Person CHF 12. Die Klasse
fährt um 8:15 Uhr ab und ist um 16:40 Uhr zurück. Die Schule übernimmt
CHF 300 der Gesamtkosten. Wie viel muss jede Person selbst bezahlen?
\begin{solution}
  Kosten pro Person: $18.50 + 12 = 30.50$; Gesamtkosten:
  $24 \cdot 30.50 = 732$; nach Abzug: $732 - 300 = 432$; pro Person:
  $432 \geteilt 24 = 18$. Jede Person bezahlt CHF 18.
  \diagnose{Die Zeitangaben sind überflüssig. Wer sie verwendet, kann die
  relevanten Informationen nicht auswählen. Wird CHF 300 von den Kosten einer
  Person abgezogen, so fehlt die Unterscheidung zwischen Gesamt- und
  Einzelgrösse.}
\end{solution}
\raum{4.5cm}

% ===========================================================================
\newpage
\teil{C}{Rechnen mit Variablen}
% ===========================================================================

% ---------------------------------------------------------------------------
% 1 Zahlen und Stellenwert
% ---------------------------------------------------------------------------
\aufgabe{1.5}
\begin{parts}
  \part Schreiben Sie in wissenschaftlicher Schreibweise:
        \quad $45\,000\,000$ \qquad $0.00072$
  \part Schreiben Sie als Dezimalzahl:
        \quad $3.1 \cdot 10^5$ \qquad $6 \cdot 10^{-3}$
  \part Berechnen Sie die beiden Ausdrücke und geben Sie die Ergebnisse in
        wissenschaftlicher Schreibweise bzw.\ als Zehnerpotenz an.
  \begin{multicols}{2}
  \begin{subparts}
    \subpart $(2 \cdot 10^3) \cdot (4 \cdot 10^5)$
    \subpart $10^6 \geteilt 10^2$
  \end{subparts}
  \end{multicols}
\end{parts}
\begin{solution}
  a) $4.5 \cdot 10^7$; \ $7.2 \cdot 10^{-4}$ \quad
  b) $310\,000$; \ $0.006$ \quad
  c) $8 \cdot 10^8$; \ $10^4$
  \diagnose{Falsch gezählte Stellen zeigen fehlendes Stellenwertverständnis.
  Werden negative Exponenten als negative Zahlen gedeutet ($6 \cdot 10^{-3}
  = -6000$), fehlt die Bedeutung der Schreibweise. $10^{15}$ bzw.\ $10^3$ bei
  c) zeigen, dass Exponenten multipliziert bzw.\ dividiert statt addiert bzw.\
  subtrahiert werden.}
\end{solution}
\raum{3.5cm}

% ---------------------------------------------------------------------------
% 2 Operationen und Rechenreihenfolge
% ---------------------------------------------------------------------------
\aufgabe{2.5}
Berechnen Sie den Wert der Terme für $a = -2$ und $b = 3$.
\begin{multicols}{3}
\begin{parts}
  \part $3a + b$
  \part $a^2$
  \part $-a^2$
  \part $2b - a$
  \part $a \cdot b^2$
\end{parts}
\end{multicols}
\begin{solution}
  a) $-3$ \quad b) $4$ \quad c) $-4$ \quad d) $8$ \quad e) $-18$
  \diagnose{Eine der aussagekräftigsten Aufgaben. $-4$ bei b) zeigt, dass ohne
  Klammern eingesetzt wird ($-2^2$ statt $(-2)^2$). $4$ bei c) zeigt
  Unsicherheit mit dem Minuszeichen vor der Potenz. $36$ bei e) zeigt, dass
  $(ab)^2$ gerechnet wird.}
\end{solution}
\raum{3cm}

% ---------------------------------------------------------------------------
% 3 Brüche
% ---------------------------------------------------------------------------
\aufgabe{3.5}
Fassen Sie zusammen und vereinfachen Sie so weit wie möglich.
\begin{multicols}{2}
\begin{parts}\bruchabstand
  \part $\dfrac{a}{4} + \dfrac{b}{6}$
  \part $\dfrac{x}{3} - \dfrac{x}{5}$
  \part $\dfrac{3}{x} + \dfrac{2}{x}$
  \part $\dfrac{2}{a} \cdot \dfrac{a}{6}$
\end{parts}
\end{multicols}
\begin{solution}
  a) $\frac{3a + 2b}{12}$ \quad b) $\frac{2x}{15}$ \quad c) $\frac{5}{x}$ \quad
  d) $\frac{1}{3}$
  \diagnose{Diese Aufgabe entspricht 3.3 mit Variablen. Gelingt 3.3, aber nicht
  3.5, so liegt die Schwierigkeit beim Übergang zur Algebra, nicht bei den
  Bruchregeln. $\frac{5}{2x}$ bei c) oder $\frac{a+b}{10}$ bei a) zeigen
  getrenntes Addieren von Zählern und Nennern.}
\end{solution}
\raum{4cm}

% ---------------------------------------------------------------------------
% 4 Proportionalität und Prozent
% ---------------------------------------------------------------------------
\aufgabe{4.5}
\begin{parts}
  \part Entscheiden Sie bei jeder Tabelle, ob die Zuordnung proportional,
        umgekehrt proportional oder keines von beidem ist. Geben Sie bei
        proportionalen Zuordnungen die Gleichung $y = k \cdot x$ an.
  \begin{center}
    \renewcommand{\arraystretch}{1.3}
    \newcommand{\tab}[5]{%
      \begin{tabular}[t]{@{}c@{}}
        \small Tabelle #1 \\[0.5ex]
        \begin{tabular}{|c|c|c|c|c|}
          \hline $x$ & 1 & 2 & 4 & 5 \\ \hline
                 $y$ & #2 & #3 & #4 & #5 \\ \hline
        \end{tabular}
      \end{tabular}}
    \tab{1}{3}{6}{12}{15} \qquad \tab{2}{5}{7}{11}{13} \qquad \tab{3}{20}{10}{5}{4}
  \end{center}
  \part $4$ Personen brauchen für eine Arbeit $6$ Tage. Wie lange brauchen
        $3$ Personen bei gleichem Arbeitstempo?
\end{parts}
\begin{solution}
  a) Tabelle 1: proportional, $y = 3x$; \ Tabelle 2: keines von beidem
  ($y = 2x + 3$); \ Tabelle 3: umgekehrt proportional ($x \cdot y = 20$)
  \qquad b) $8$ Tage
  \diagnose{Tabelle 2 ist die zentrale Falle: «gleichmässig wachsend» wird mit
  proportional verwechselt. $4.5$ Tage bei b) zeigen, dass ohne Prüfung direkt
  proportional gerechnet wird.}
\end{solution}
\raum{3.5cm}

% ---------------------------------------------------------------------------
% 5 Terme
% ---------------------------------------------------------------------------
\aufgabe{5.5}
\begin{parts}
  \part Multiplizieren Sie aus und vereinfachen Sie.
  \begin{multicols}{4}
    $3 \cdot (x + 4)$ \columnbreak

    $(-2) \cdot (a - 5)$ \columnbreak

    $(x + 3)(x + 2)$ \columnbreak

    $(a + b)^2$
  \end{multicols}
  \part Klammern Sie so weit wie möglich aus.
        \qquad $6x + 9$ \qquad\qquad $x^2 - 5x$
\end{parts}
\begin{solution}
  a) $3x + 12$; \ $-2a + 10$; \ $x^2 + 5x + 6$; \ $a^2 + 2ab + b^2$ \quad
  b) $3 \cdot (2x + 3)$; \ $x(x - 5)$
  \diagnose{$3x + 4$ zeigt, dass nur der erste Summand multipliziert wird.
  $-2a - 10$ zeigt Vorzeichenfehler. $(a + b)^2 = a^2 + b^2$ ist der häufigste
  Fehler überhaupt und zeigt, dass Potenzieren über die Addition «verteilt»
  wird.}
\end{solution}
\raum{4cm}

% ---------------------------------------------------------------------------
% 6 Gleichungen
% ---------------------------------------------------------------------------
\aufgabe{6.5}
\begin{parts}
  \part Lösen Sie die Gleichungen.
  \begin{multicols}{3}
    $3 \cdot (x - 2) = 2x + 1$ \columnbreak

    $\dfrac{x}{2} + \dfrac{x}{3} = 10$ \columnbreak

    $\dfrac{x + 1}{4} = \dfrac{x - 2}{3}$
  \end{multicols}
  \part Prüfen Sie, \textbf{ohne die Gleichung zu lösen}, ob $x = 2$ eine Lösung
        der folgenden Gleichung ist.
        \[ 5x - 3 = 2 \cdot (x + 1) + 1 \]
\end{parts}
\begin{solution}
  a) $x = 7$; \ $x = 12$; \ $x = 11$ \qquad
  b) Ja: Für $x = 2$ ergeben beide Seiten $7$.
  \diagnose{Teil b) prüft, ob verstanden wird, was eine Lösung ist, unabhängig
  vom Lösungsverfahren. Wer b) nur durch Lösen der Gleichung beantworten kann,
  verwendet ein Verfahren ohne Begriffsverständnis.}
\end{solution}
\raum{4.5cm}

% ---------------------------------------------------------------------------
% 7 Funktionen und Graphen
% ---------------------------------------------------------------------------
\aufgabe{7.5}
\begin{parts}
  \part Geben Sie für die Gerade im linken Koordinatensystem die Steigung, den
        $y$-Achsenabschnitt und die Funktionsgleichung an.
  \part Zeichnen Sie den Graphen von $y = -\frac{1}{2}x + 3$ ins rechte
        Koordinatensystem.
\end{parts}

\begin{center}
\newcommand{\koord}[4]{% xmin, xmax, ymin, ymax
  \draw[gray!35, thin] (#1,#3) grid (#2,#4);
  \draw[->, thick] (#1-0.4,0) -- (#2+0.5,0) node[right] {\small $x$};
  \draw[->, thick] (0,#3-0.4) -- (0,#4+0.5) node[above] {\small $y$};
  \foreach \t in {#1,...,#2} {\ifnum\t=0\else
    \draw (\t,0.1) -- (\t,-0.1) node[below, font=\tiny] {$\t$};\fi}
  \foreach \t in {#3,...,#4} {\ifnum\t=0\else
    \draw (0.1,\t) -- (-0.1,\t) node[left, font=\tiny] {$\t$};\fi}}
\begin{tikzpicture}[scale=0.5]
  \koord{-3}{4}{-3}{7}
  \draw[very thick, domain=-2:3] plot (\x, {2*\x + 1});
  \fill (0,1) circle (4pt) (2,5) circle (4pt);
\end{tikzpicture}
\hspace{1.5cm}
\begin{tikzpicture}[scale=0.5]
  \koord{-2}{8}{-2}{6}
  \ifprintanswers
    \draw[very thick, domain=-2:8] plot (\x, {-0.5*\x + 3});
  \fi
\end{tikzpicture}
\end{center}
\begin{parts}
  \setcounter{partno}{2}
  \part Liegt der Punkt $\pkt{4}{7}$ auf der Geraden $y = 2x - 1$? Begründen Sie.
  \part Die Werte in der Tabelle gehören zu einer linearen Funktion.
        Bestimmen Sie die Funktionsgleichung.
  \quad
  \begin{tabular}{|c|c|c|c|}
    \hline $x$ & $0$ & $1$ & $2$  \\ \hline
           $y$ & $4$ & $1$ & $-2$ \\ \hline
  \end{tabular}
\end{parts}
\begin{solution}
  a) Steigung $2$, $y$-Achsenabschnitt $1$, $y = 2x + 1$ \quad
  b) Gerade durch $\pkt{0}{3}$ und $\pkt{6}{0}$ (siehe Grafik) \quad
  c) Ja, denn $2 \cdot 4 - 1 = 7$. \quad
  d) $y = -3x + 4$
  \diagnose{Die Aufgabe prüft alle Wechsel zwischen Graph, Gleichung und
  Tabelle. Typische Fehler: Steigung als $\frac{\Delta x}{\Delta y}$ berechnet
  ($\frac{1}{2}$), negative Steigung steigend gezeichnet, $y$-Achsenabschnitt
  und Nullstelle verwechselt.}
\end{solution}
\raum{3.5cm}

% ---------------------------------------------------------------------------
% 8 Geometrie und Grössen
% ---------------------------------------------------------------------------
\aufgabe{8.5}
\begin{parts}
  \part Ein rechtwinkliges Dreieck hat die Katheten $6$\,cm und $8$\,cm.
        Berechnen Sie die Länge der Hypotenuse.
  \part In einem rechtwinkligen Dreieck ist die Hypotenuse $13$\,cm und eine
        Kathete $5$\,cm lang. Berechnen Sie die Länge der anderen Kathete.
  \part Eine $10$\,m lange Leiter lehnt an einer senkrechten Wand. Ihr Fuss
        steht $6$\,m von der Wand entfernt. In welcher Höhe berührt die Leiter
        die Wand?
  \part Ist ein Dreieck mit den Seitenlängen $7$\,cm, $24$\,cm und $25$\,cm
        rechtwinklig? Begründen Sie.
\end{parts}
\begin{solution}
  a) $10$\,cm \quad b) $12$\,cm \quad c) $8$\,m \quad
  d) Ja, denn $7^2 + 24^2 = 49 + 576 = 625 = 25^2$.
  \diagnose{$\sqrt{194}$ bei b) zeigt, dass $a^2 + b^2 = c^2$ mechanisch
  angewendet wird, ohne die Hypotenuse zu identifizieren. c) verlangt, das
  rechtwinklige Dreieck in einer Situation zu erkennen. d) prüft die Umkehrung
  des Satzes von Pythagoras.}
\end{solution}
\raum{5cm}

% ---------------------------------------------------------------------------
% 9 Sprache ↔ Mathematik
% ---------------------------------------------------------------------------
\aufgabe{9.5}
\begin{parts}
  \part Schreiben Sie als Term: \quad «$5$ weniger als das Doppelte von $x$»
        \quad und \quad «um $3$ grösser als die Hälfte von $y$»
  \part An einer Schule gibt es sechsmal so viele Schülerinnen und Schüler wie
        Lehrpersonen. Schreiben Sie diese Aussage als Gleichung. Verwenden Sie
        $S$ für die Anzahl Schülerinnen und Schüler und $L$ für die Anzahl
        Lehrpersonen.
  \part Beschreiben Sie die Terme in Worten: \quad $3 \cdot (x + 4)$ \quad und \quad
        $3x + 4$
  \part Verdoppelt man eine Zahl und addiert $7$, so erhält man $31$. Stellen
        Sie eine Gleichung auf und bestimmen Sie die Zahl.
\end{parts}
\begin{solution}
  a) $2x - 5$; \ $\frac{y}{2} + 3$ \quad
  b) $S = 6L$ \quad
  c) «das Dreifache der Summe von $x$ und $4$»; «die Summe aus dem Dreifachen
  von $x$ und $4$» \quad
  d) $2x + 7 = 31$, $x = 12$
  \diagnose{b) ist die klassische «Studierende-und-Professoren»-Aufgabe: $6S = L$
  entsteht durch Wort-für-Wort-Übersetzung, ohne die Beziehung zu erfassen.
  $5 - 2x$ bei a) zeigt die Umkehrung bei «weniger als». c) ist die
  Gegenrichtung zu 9.3.}
\end{solution}
\raum{4.5cm}

% ===========================================================================
\newpage
\teil{D}{Vertiefung}
% ===========================================================================

% ---------------------------------------------------------------------------
% 1 Zahlen und Stellenwert
% ---------------------------------------------------------------------------
\aufgabe{1.6}
\begin{parts}
  \part Geben Sie für jede Zahl alle Zahlenmengen ($\mathbb{N}$, $\mathbb{Z}$,
        $\mathbb{Q}$, $\mathbb{R}$) an, zu denen sie gehört.
        \[
          -3 \qquad 0.75 \qquad \sqrt{2} \qquad \frac{12}{4} \qquad 0.333\ldots
        \]
  \part Kann das Ergebnis stimmen? Entscheiden Sie \textbf{ohne nachzurechnen}
        und begründen Sie.
        \begin{itemize}
          \item Lena rechnet: $0.6 \geteilt 0.03 = 2$.
          \item Ein Schüler berechnet das Volumen eines Schulzimmers und erhält
                $18$\,m$^3$.
        \end{itemize}
\end{parts}
\begin{solution}
  a) $-3$: $\mathbb{Z}, \mathbb{Q}, \mathbb{R}$; \
  $0.75$: $\mathbb{Q}, \mathbb{R}$; \
  $\sqrt{2}$: $\mathbb{R}$; \
  $\frac{12}{4}$: $\mathbb{N}, \mathbb{Z}, \mathbb{Q}, \mathbb{R}$; \
  $0.333\ldots$: $\mathbb{Q}, \mathbb{R}$ \\
  b) Beides ist nicht plausibel: Division durch eine Zahl kleiner als $1$
  vergrössert (richtig: $20$). Ein Schulzimmer hat etwa $150$ bis
  $200$\,m$^3$.
  \diagnose{Wird $\frac{12}{4}$ nur $\mathbb{Q}$ zugeordnet, so wird die Zahl
  nach ihrer Schreibweise statt nach ihrem Wert beurteilt. Die Zahl $0$ wurde
  bewusst vermieden, da Lehrmittel uneinheitlich sind, ob $0 \in \mathbb{N}$.}
\end{solution}
\raum{4cm}

% ---------------------------------------------------------------------------
% 2 Operationen und Rechenreihenfolge
% ---------------------------------------------------------------------------
\aufgabe{2.6}
\begin{parts}
  \part Welche Terme sind für alle Werte von $x$ gleichwertig mit $2 \cdot (x - 3)$?
        Kreuzen Sie an.

        \wahl{$2x - 3$} \quad \wahlr{$2x - 6$} \quad \wahlr{$(-2) \cdot (3 - x)$}
        \quad \wahlr{$(x - 3) \cdot 2$} \quad \wahl{$2 \cdot (3 - x)$}
  \part Gilt $(a - b)^2 = (b - a)^2$ für alle Zahlen $a$ und $b$? Begründen Sie.
\end{parts}
\begin{solution}
  a) $2x - 6$, \ $(-2) \cdot (3 - x)$, \ $(x - 3) \cdot 2$ \qquad
  b) Ja, denn $b - a = -(a - b)$, und beim Quadrieren fällt das Vorzeichen weg.
  \diagnose{Die Aufgabe prüft, ob Gleichwertigkeit an der Struktur erkannt wird.
  Wer nur einzelne Werte einsetzt, kann b) nicht allgemein begründen.}
\end{solution}
\raum{3cm}

% ---------------------------------------------------------------------------
% 3 Brüche
% ---------------------------------------------------------------------------
\aufgabe{3.6}
\begin{parts}
  \part Fassen Sie zusammen: \quad $\dfrac{1}{x} + \dfrac{1}{x + 1}$
  \part Kürzen Sie, falls möglich:
        \quad $\dfrac{3x + 6}{x + 2}$ \qquad $\dfrac{x^2 - 9}{x - 3}$
        \qquad $\dfrac{x + 5}{5}$
\end{parts}
\begin{solution}
  a) $\frac{2x + 1}{x(x + 1)}$ \quad
  b) $3$; \ $x + 3$; \ nicht kürzbar
  \diagnose{$\frac{x + 5}{5} = x + 1$ oder $= x$ zeigt «Kürzen aus Summen». Wer
  bei den ersten beiden Brüchen nicht faktorisiert, erkennt die Struktur nicht.
  Die Angabe der Definitionslücken ($x \neq -2$, $x \neq 3$) wird auf diesem
  Niveau nicht verlangt.}
\end{solution}
\raum{3.5cm}

% ---------------------------------------------------------------------------
% 4 Proportionalität und Prozent
% ---------------------------------------------------------------------------
\aufgabe{4.6}
\begin{parts}
  \part Ein Preis wird zuerst um $20\,\%$ erhöht und danach um $20\,\%$
        gesenkt. Ist er danach gleich hoch wie zu Beginn? Begründen Sie.
  \part Nach einer Erhöhung um $25\,\%$ kostet ein Artikel CHF 150. Wie viel
        kostete er vorher?
  \part Eine Bevölkerung von $P$ Personen wächst jährlich um $3\,\%$. Geben Sie
        einen Term für die Bevölkerung nach $n$ Jahren an.
\end{parts}
\begin{solution}
  a) Nein: $1.2 \cdot 0.8 = 0.96$, der Preis ist $4\,\%$ tiefer. \quad
  b) $150 \geteilt 1.25 = 120$, also CHF 120 \quad
  c) $P \cdot 1.03^n$
  \diagnose{CHF 112.50 bei b) zeigt, dass $25\,\%$ vom neuen statt vom alten
  Wert abgezogen werden. $P + 0.03n$ oder $P \cdot 0.03^n$ bei c) zeigen, dass
  der Wachstumsfaktor nicht verstanden ist.}
\end{solution}
\raum{4cm}

% ---------------------------------------------------------------------------
% 5 Terme
% ---------------------------------------------------------------------------
\aufgabe{5.6}
\begin{parts}
  \part Faktorisieren Sie so weit wie möglich.
  \begin{multicols}{4}
    $x^2 - 9$ \columnbreak

    $x^2 + 6x + 9$ \columnbreak

    $4a^2b - 6ab^2$ \columnbreak

    $ax + ay - bx - by$
  \end{multicols}
  \part Ergänzen Sie: \quad $x^2 + \kaestchen\,x + 25 = (x + \kaestchen)^2$
\end{parts}
\begin{solution}
  a) $(x - 3)(x + 3)$; \ $(x + 3)^2$; \ $2ab(2a - 3b)$; \ $(a - b)(x + y)$
  \quad b) $10$ und $5$
  \diagnose{Faktorisieren verlangt, die binomischen Formeln «rückwärts» zu
  lesen. Wer 5.5 lösen kann, hier aber scheitert, kennt die Formeln nur in eine
  Richtung.}
\end{solution}
\raum{3.5cm}

% ---------------------------------------------------------------------------
% 6 Gleichungen
% ---------------------------------------------------------------------------
\aufgabe{6.6}
\begin{parts}
  \part Lösen Sie nach der angegebenen Variable auf.
  \begin{multicols}{3}
    $A = \frac{1}{2} \cdot g \cdot h$ \ nach $h$ \columnbreak

    $v = \dfrac{s}{t}$ \ nach $t$ \columnbreak

    $F = \frac{9}{5} \cdot C + 32$ \ nach $C$
  \end{multicols}
  \part Lösen Sie nach $x$ auf ($a \neq 3$): \quad $ax = 3x + 5$
  \part Bestimmen Sie die Lösungsmenge.
        \quad $2 \cdot (x + 1) = 2x + 2$ \qquad $2 \cdot (x + 1) = 2x + 3$
\end{parts}
\begin{solution}
  a) $h = \frac{2A}{g}$; \ $t = \frac{s}{v}$; \ $C = \frac{5}{9} \cdot (F - 32)$
  \quad b) $x = \frac{5}{a - 3}$ \quad
  c) $\mathbb{L} = \mathbb{R}$; \ $\mathbb{L} = \{\,\}$
  \diagnose{Bei b) muss $x$ ausgeklammert werden ($x(a - 3) = 5$). Wer nur
  Gleichungen mit Zahlen lösen kann, kommt hier nicht weiter. Bei c) zeigt
  die Reaktion «$0 = 1$, also habe ich mich verrechnet», dass Lösungsmengen
  ohne genau eine Lösung unbekannt sind.}
\end{solution}
\raum{4.5cm}

% ---------------------------------------------------------------------------
% 7 Funktionen und Graphen
% ---------------------------------------------------------------------------
\aufgabe{7.6}
\begin{parts}
  \part Was haben alle Geraden mit der Gleichung $y = mx + 2$ gemeinsam? Was
        haben alle Geraden mit der Gleichung $y = 3x + q$ gemeinsam?
  \part Berechnen Sie den Schnittpunkt der Geraden $y = 2x + 1$ und
        $y = -x + 7$.
  \part Für welchen Wert von $m$ verläuft die Gerade $y = mx + 2$ durch den
        Punkt $\pkt{4}{0}$?
\end{parts}
\begin{solution}
  a) Alle verlaufen durch $\pkt{0}{2}$; alle sind parallel (Steigung $3$).
  \quad b) $2x + 1 = -x + 7 \Rightarrow x = 2$, also $S\pkt{2}{5}$
  \quad c) $0 = 4m + 2 \Rightarrow m = -\frac{1}{2}$
  \diagnose{Die Aufgabe prüft, ob ein Buchstabe als Parameter (fest, aber
  beliebig) und nicht als Variable verstanden wird, und ob ein Schnittpunkt als
  Lösung einer Gleichung erkannt wird.}
\end{solution}
\raum{4cm}

% ---------------------------------------------------------------------------
% 8 Geometrie und Grössen
% ---------------------------------------------------------------------------
\aufgabe{8.6}
\begin{parts}
  \part Die Geraden $AA'$ und $BB'$ sind parallel. Berechnen Sie die Längen
        $\overline{SB'}$ und $\overline{BB'}$. (Skizze nicht massstäblich)

  \begin{center}
  \begin{tikzpicture}[scale=0.8]
    \coordinate (S) at (0,0);
    \coordinate (A) at (4,0);   \coordinate (B) at (6,0);
    \coordinate (A2) at (3.2,2.4); \coordinate (B2) at (4.8,3.6);
    \draw[thick] (S) -- (7,0);
    \draw[thick] (S) -- (5.6,4.2);
    \draw[thick] (A) -- (A2);
    \draw[thick] (B) -- (B2);
    \foreach \p/\pos in {S/left, A/below, B/below, A2/above left, B2/above left}
      \fill (\p) circle (2pt);
    \node[left]        at (S)  {$S$};
    \node[below]       at (A)  {$A$};
    \node[below]       at (B)  {$B$};
    \node[above left]  at (A2) {$A'$};
    \node[above left]  at (B2) {$B'$};
    \node[below, font=\small] at (2,0)       {$4$\,cm};
    \node[below, font=\small] at (5,0)       {$2$\,cm};
    \node[above left, font=\small] at (1.6,1.2) {$5$\,cm};
    \node[right, font=\small] at (3.6,1.2)   {$3$\,cm};
  \end{tikzpicture}
  \end{center}
  \part Eine $1.80$\,m grosse Person wirft einen $1.20$\,m langen Schatten. Ein
        Baum wirft gleichzeitig einen $8$\,m langen Schatten. Wie hoch ist der
        Baum?
  \part Ein Quadrat mit der Seitenlänge $a$ wird mit dem Faktor $k$ vergrössert.
        Geben Sie den Umfang und den Flächeninhalt des neuen Quadrats an.
\end{parts}
\begin{solution}
  a) $\overline{SB} = 6$\,cm, also $\overline{SB'} = 5 \cdot \frac{6}{4} =
  7.5$\,cm und $\overline{BB'} = 3 \cdot \frac{6}{4} = 4.5$\,cm \quad
  b) $12$\,m \quad
  c) Umfang $4ka$, Flächeninhalt $k^2a^2$
  \diagnose{Bei a) zeigen Ergebnisse wie $\overline{BB'} = 1.5$\,cm oder
  $\overline{SB'} = 2.5$\,cm, dass mit der Teilstrecke $\overline{AB}$ statt
  mit $\overline{SB}$ gerechnet wird. Additive Ansätze ($3 + 2 = 5$) zeigen
  fehlendes Verständnis der Ähnlichkeit. $ka^2$ bei c) zeigt, dass die Fläche
  wie die Länge skaliert wird.}
\end{solution}
\raum{4cm}

% ---------------------------------------------------------------------------
% 9 Sprache ↔ Mathematik
% ---------------------------------------------------------------------------
\aufgabe{9.6}
\begin{parts}
  \part Handyabo A kostet CHF 20 pro Monat plus CHF 0.10 pro Gesprächsminute.
        Abo B kostet pauschal CHF 35 pro Monat. Stellen Sie für beide Abos
        einen Term für die Monatskosten bei $x$ Minuten auf. Ab wie vielen
        Minuten ist Abo B günstiger?
  \part Für einen Schulausflug fahren $200$ Personen mit Bussen, die je
        $45$ Plätze haben. Wie viele Busse braucht es?
\end{parts}
\begin{solution}
  a) A: $20 + 0.1x$; B: $35$. Aus $20 + 0.1x = 35$ folgt $x = 150$. Abo B ist
  ab mehr als $150$ Minuten günstiger. \quad
  b) $200 \geteilt 45 \approx 4.4$, also braucht es $5$ Busse.
  \diagnose{Die Antwort «$150$» ohne Aussage, was darunter und darüber gilt,
  zeigt, dass die Lösung nicht im Kontext gedeutet wird. Bei b) sind $4.44$
  oder $4$ Busse die klassischen Zeichen dafür, dass ein Rechenergebnis nicht
  auf die Situation zurückbezogen wird.}
\end{solution}
\raum{4.5cm}


\end{questions}

% ===========================================================================
% Auswertungsraster (für die Lehrperson)
% ===========================================================================
\newpage
\section*{Auswertung \normalsize(wird von der Lehrperson ausgefüllt)}

Pro Aufgabe: \quad \textbf{2} = sicher gelöst \quad \textbf{1} = teilweise
gelöst \quad \textbf{0} = nicht gelöst. \\
Die \emph{Bruchstelle} ist das tiefste Niveau eines Strangs, auf dem die
Aufgaben nicht mehr sicher gelöst werden.

\begin{center}
\small
\renewcommand{\arraystretch}{2.2}
\newcommand{\feld}[1]{#1\,\framebox[0.6cm]{\rule{0pt}{0.9em}}}
\begin{tabular}{|>{\raggedright\arraybackslash}p{3.5cm}|c|c|c|c|p{1.9cm}|}
  \hline
  \textbf{Strang}
    & \textbf{A}\,\footnotesize Grundlagen & \textbf{B}\,\footnotesize Numerisch
    & \textbf{C}\,\footnotesize Symbol. & \textbf{D}\,\footnotesize Formal
    & \footnotesize\textbf{Bruchstelle} \\ \hline
  1 Zahlen und Stellenwert          & \feld{1.1}\ \feld{1.2} & \feld{1.3}\ \feld{1.4} & \feld{1.5} & \feld{1.6} & \\ \hline
  2 Operationen und Rechenreihenfolge & \feld{2.1}\ \feld{2.2} & \feld{2.3}\ \feld{2.4} & \feld{2.5} & \feld{2.6} & \\ \hline
  3 Brüche                          & \feld{3.1}\ \feld{3.2} & \feld{3.3}\ \feld{3.4} & \feld{3.5} & \feld{3.6} & \\ \hline
  4 Proportionalität und Prozent    & \feld{4.1}\ \feld{4.2} & \feld{4.3}\ \feld{4.4} & \feld{4.5} & \feld{4.6} & \\ \hline
  5 Terme                           & \feld{5.1}\ \feld{5.2} & \feld{5.3}\ \feld{5.4} & \feld{5.5} & \feld{5.6} & \\ \hline
  6 Gleichungen                     & \feld{6.1}\ \feld{6.2} & \feld{6.3}\ \feld{6.4} & \feld{6.5} & \feld{6.6} & \\ \hline
  7 Funktionen und Graphen          & \feld{7.1}\ \feld{7.2} & \feld{7.3}\ \feld{7.4} & \feld{7.5} & \feld{7.6} & \\ \hline
  8 Geometrie und Grössen           & \feld{8.1}\ \feld{8.2} & \feld{8.3}\ \feld{8.4} & \feld{8.5} & \feld{8.6} & \\ \hline
  9 Sprache $\leftrightarrow$ Mathematik & \feld{9.1}\ \feld{9.2} & \feld{9.3}\ \feld{9.4} & \feld{9.5} & \feld{9.6} & \\ \hline
\end{tabular}
\end{center}

\textbf{Bemerkungen:}

\vspace{4cm}

\end{document}
"
% Jeweils zweimal kompilieren (Seitenzahlen).
%
% Aufbau: 9 Stränge × 4 Niveaus
%   A Grundlagen  (Aufgaben x.1, x.2)
%   B Numerisch   (Aufgaben x.3, x.4)
%   C Symbolisch  (Aufgabe  x.5)
%   D Formal      (Aufgabe  x.6)
%%%%%%%%%%%%%%%%%%%%%%%%%%%%%%%%%%%%%%%%%%%%%%%%%%%%%%%%%%%%%%%%%%%%%%%%%%%%%%

\ifdefined\mitloesungen \printanswers \fi

\usepackage[T1]{fontenc}
\usepackage{lmodern}
\usepackage[nswissgerman]{babel}
\usepackage[hidelinks]{hyperref}
\usepackage{amssymb, amsmath, setspace, multicol, array, needspace}
\usepackage{tikz}
\usepackage[margin=2cm]{geometry}

\usepackage{color}
\colorgrids
\definecolor{GridColor}{gray}{0.7}
\setlength{\gridsize}{4mm}

\singlespacing
\parindent 0ex

% ---------------------------------------------------------------------------
% Kopfzeilen und Angaben
% ---------------------------------------------------------------------------
\newcommand{\titel}{Standortbestimmung Mathematik}
\newcommand{\schuljahr}{Schuljahr 2026/27}
\newcommand{\zeit}{90 Minuten (Teile A und B)}
\newcommand{\hilfsmittel}{keine (kein Taschenrechner)}

% ---------------------------------------------------------------------------
% Hilfsbefehle
% ---------------------------------------------------------------------------
% Aufgabe mit fester Nummer (Strang.Nummer) statt fortlaufender Zählung
\newcommand{\aufgabe}[1]{%
  \needspace{8\baselineskip}%
  \renewcommand{\questionlabel}{\textbf{#1}}%
  \question}

% Teil-Überschrift über die volle Breite
\newcommand{\teil}[2]{\uplevel{\vspace{1ex}\textbf{\large Teil #1: #2}\par\vspace{1ex}}}

% Divisionszeichen (Doppelpunkt) mit korrektem Abstand
\newcommand{\geteilt}{\mathbin{:}}
% Punkt mit Koordinaten (x | y)
\newcommand{\pkt}[2]{(#1\,|\,#2)}
% Leeres Kästchen für Lückenaufgaben
\newcommand{\kaestchen}{\fbox{\rule{0pt}{1.1em}\hspace{1.6em}}}
% Platz für die Antwort (Raster), nur in der Version ohne Lösungen
\newcommand{\raum}[1]{\ifprintanswers\else\fillwithgrid{#1}\fi}
% Ankreuzfelder: \wahl = falsche, \wahlr = richtige Option
\newcommand{\wahl}[1]{\mbox{$\Box$\,#1}}
\newcommand{\wahlr}[1]{\mbox{\ifprintanswers$\blacksquare$\,\textbf{#1}\else$\Box$\,#1\fi}}
% Mehr Zeilenabstand in Aufgabenteilen mit Brüchen
\newcommand{\bruchabstand}{\setlength{\itemsep}{2.5ex}}
% Interpretationshinweis in den Lösungen
\newcommand{\diagnose}[1]{\par\smallskip\textit{Interpretation:} #1}

% ---------------------------------------------------------------------------
% Deutsche Bezeichnungen für die exam-Klasse
% ---------------------------------------------------------------------------
\renewcommand{\solutiontitle}{\noindent\textbf{Lösung:}\enspace}
\checkboxchar{$\Box$}
\checkedchar{$\blacksquare$}
\CorrectChoiceEmphasis{\bfseries}
\pointsinmargin
\nopointsinmargin

\begin{document}

\pagestyle{headandfoot}
\firstpageheader{}{}{}
\runningheader{\titel\ifprintanswers\ – Lösungen\fi}{}{\schuljahr}
\runningheadrule
\footer{}{Seite \thepage\ von \numpages}{}

% ===========================================================================
% Titelseite
% ===========================================================================
\begin{center}
  {\LARGE\bfseries \titel\ifprintanswers\ – Lösungen\fi}\\[1ex]
  {\large Eintritt ins Gymnasium – \schuljahr}
\end{center}
\vspace{1em}

\begin{tabular}{@{}p{0.5\textwidth}l@{}}
  \textbf{Name:} \makebox[5.5cm]{\hrulefill} & \textbf{Klasse:} \makebox[3cm]{\hrulefill} \\[2ex]
  \textbf{Zeit:} \zeit                       & \textbf{Hilfsmittel:} \hilfsmittel
\end{tabular}

\vspace{1.5em}
\rule{\textwidth}{.4pt}

\subsection*{Hinweise}

\begin{itemize}
  \item Diese Standortbestimmung ist \textbf{keine Prüfung} und wird nicht benotet. Sie zeigt
        uns, wo Sie in Mathematik stehen, damit wir Sie gezielt unterstützen können.
  \item Die Aufgaben werden von Teil zu Teil schwieriger. Es ist normal, dass
        Sie nicht alle Aufgaben lösen können. Versuchen Sie trotzdem jede Aufgabe.
  \item Lösen Sie die Aufgaben \textbf{ohne Taschenrechner}.
  \item Schreiben Sie Ihre Überlegungen und Zwischenschritte auf. Ein Lösungsweg
        sagt uns oft mehr als das Ergebnis.
  \item Wenn Sie eine Aufgabe nicht lösen können, gehen Sie zur nächsten weiter.
\end{itemize}

\vfill
\newpage

% ===========================================================================
\begin{questions}
% ===========================================================================

\teil{A}{Grundlagen}

% ---------------------------------------------------------------------------
% 1 Zahlen und Stellenwert
% ---------------------------------------------------------------------------
\aufgabe{1.1}
Ordnen Sie die Zahlen der Grösse nach, beginnend mit der kleinsten:
\[
  0.5; \quad 0.45; \quad 0.405; \quad 0.54
\]
\begin{solution}
  $0.405 < 0.45 < 0.5 < 0.54$
  \diagnose{Wird $0.405$ oder $0.45$ grösser als $0.5$ eingeordnet, so gilt
  vermutlich die Vorstellung «mehr Ziffern bedeutet eine grössere Zahl».
  Der Stellenwert der Dezimalstellen ist nicht gesichert.}
\end{solution}
\raum{2.5cm}

\aufgabe{1.2}
\begin{parts}
  \part Berechnen Sie: \quad $3.7 \cdot 100$ \qquad $45 \geteilt 1000$
  \part Welche Zahl liegt genau in der Mitte zwischen $0.3$ und $0.4$?
\end{parts}
\begin{solution}
  a) $370$; \ $0.045$ \qquad b) $0.35$
  \diagnose{Ergebnisse wie $3.700$ oder $0.45$ deuten auf die auswendig gelernte
  Regel «Nullen anhängen bzw.\ wegstreichen» hin, ohne Verständnis des
  Stellenwerts. Bei b) zeigt die Antwort «gibt es nicht», dass die Dichte der
  Dezimalzahlen nicht verstanden ist.}
\end{solution}
\raum{3cm}

% ---------------------------------------------------------------------------
% 2 Operationen und Rechenreihenfolge
% ---------------------------------------------------------------------------
\aufgabe{2.1}
Ein Band von $3.6$\,m Länge wird in Stücke von $0.4$\,m Länge geschnitten.
Welche Rechnung passt zu der Frage, wie viele Stücke entstehen? Kreuzen Sie an.

\begin{oneparcheckboxes}
  \choice $3.6 \cdot 0.4$
  \CorrectChoice $3.6 \geteilt 0.4$
  \choice $0.4 \geteilt 3.6$
  \choice $3.6 - 0.4$
\end{oneparcheckboxes}
\begin{solution}
  $3.6 \geteilt 0.4 = 9$ Stücke.
  \diagnose{Wer die Division meidet, hat oft die Vorstellung, man dürfe nur durch
  ganze Zahlen teilen. Die Wahl von $0.4 \geteilt 3.6$ zeigt die Regel «immer die
  grössere durch die kleinere Zahl» in umgekehrter Form bzw.\ fehlendes Verständnis
  der Division als Aufteilen.}
\end{solution}

\aufgabe{2.2}
Entscheiden Sie \textbf{ohne zu rechnen}, ob das Ergebnis grösser oder kleiner
als $50$ ist.
\begin{parts}
  \part $50 \cdot 0.8$ \hfill \wahl{grösser als 50}\quad \wahlr{kleiner als 50}
  \part $50 \geteilt 0.8$ \hfill \wahlr{grösser als 50}\quad \wahl{kleiner als 50}
  \part $50 \geteilt 1.25$ \hfill \wahl{grösser als 50}\quad \wahlr{kleiner als 50}
\end{parts}
\begin{solution}
  a) kleiner ($40$) \quad b) grösser ($62.5$) \quad c) kleiner ($40$)
  \diagnose{Fehler bei a) und b) zeigen die Fehlvorstellung «Multiplizieren
  vergrössert, Dividieren verkleinert», die aus dem Rechnen mit natürlichen
  Zahlen stammt.}
\end{solution}

% ---------------------------------------------------------------------------
% 3 Brüche
% ---------------------------------------------------------------------------
\aufgabe{3.1}
Markieren Sie die Zahlen $\dfrac{3}{4}$ und $\dfrac{5}{4}$ auf dem Zahlenstrahl.

\begin{center}
\begin{tikzpicture}[x=6cm]
  \draw[->, thick] (-0.05,0) -- (2.15,0);
  \foreach \x in {0,1,2} {
    \draw[thick] (\x,0.15) -- (\x,-0.15) node[below] {$\x$};
  }
  \ifprintanswers
    \foreach \x/\l in {0.75/{\frac{3}{4}}, 1.25/{\frac{5}{4}}} {
      \draw[very thick] (\x,0.25) -- (\x,-0.25);
      \node[above] at (\x,0.25) {$\l$};
    }
  \fi
\end{tikzpicture}
\end{center}
\begin{solution}
  Markierungen bei $0.75$ und $1.25$ (siehe Zahlenstrahl).
  \diagnose{Wer $\frac{5}{4}$ nicht einträgt oder für unmöglich hält, versteht
  Brüche nur als «Teil eines Ganzen», nicht als Zahlen. Falsche Positionen von
  $\frac{3}{4}$ zeigen Schwierigkeiten beim Einteilen der Einheit.}
\end{solution}

\aufgabe{3.2}
\begin{parts}
  \part Welcher Bruch ist jeweils grösser? Begründen Sie kurz.
  \[
    \frac{2}{3} \text{ oder } \frac{3}{4} \qquad\qquad
    \frac{5}{8} \text{ oder } \frac{5}{6} \qquad\qquad
    \frac{7}{8} \text{ oder } \frac{8}{9}
  \]
  \part Ergänzen Sie: \quad $\dfrac{2}{3} = \dfrac{\kaestchen}{12}$
\end{parts}
\begin{solution}
  a) $\frac{3}{4}$; \ $\frac{5}{6}$; \ $\frac{8}{9}$ (z.\,B.\ weil $\frac{8}{9}$
  nur $\frac{1}{9}$ von $1$ entfernt ist, $\frac{7}{8}$ aber $\frac{1}{8}$)
  \qquad b) $8$
  \diagnose{Wird $\frac{5}{8}$ als grösser gewählt («8 ist grösser als 6»), so
  werden Zähler und Nenner getrennt als natürliche Zahlen betrachtet. Das letzte
  Paar unterscheidet auswendig gelernte Verfahren (gleichnamig machen) von
  Zahlverständnis (Abstand zu~1).}
\end{solution}
\raum{3.5cm}

% ---------------------------------------------------------------------------
% 4 Proportionalität und Prozent
% ---------------------------------------------------------------------------
\aufgabe{4.1}
3 Hefte kosten CHF 7.50. Wie viel kosten 12 Hefte? Wie viel kosten 5 Hefte?
\begin{solution}
  12 Hefte: CHF 30.00 \qquad 5 Hefte: CHF 12.50 (1 Heft kostet CHF 2.50)
  \diagnose{Gelingen die 12 Hefte, die 5 Hefte aber nicht, so wird nur
  vervielfacht, der Weg über den Einheitspreis (Dreisatz) fehlt.}
\end{solution}
\raum{3cm}

\aufgabe{4.2}
Für 4 Personen braucht ein Rezept 6 Eier. Wie viele Eier braucht man für
6 Personen?
\begin{solution}
  $9$ Eier
  \diagnose{Die Antwort $8$ («2 Personen mehr, also 2 Eier mehr») ist das
  klassische Zeichen für additives statt multiplikatives Denken.}
\end{solution}
\raum{3cm}

% ---------------------------------------------------------------------------
% 5 Terme
% ---------------------------------------------------------------------------
\aufgabe{5.1}
Aus Hölzchen werden Quadrate in einer Reihe gelegt.

\begin{center}
\begin{tikzpicture}[x=0.8cm, y=0.8cm, line width=1.6pt, line cap=round]
  \newcommand{\reihe}[2]{% #1 = x-Start, #2 = Anzahl Quadrate
    \foreach \i in {1,...,#2} {
      \draw (#1+\i-1,0) -- (#1+\i,0);
      \draw (#1+\i-1,1) -- (#1+\i,1);
      \draw (#1+\i,0) -- (#1+\i,1);
    }
    \draw (#1,0) -- (#1,1);
  }
  \reihe{0}{1}  \node at (0.5,-0.6) {\small 4 Hölzchen};
  \reihe{3}{2}  \node at (4,-0.6)   {\small 7 Hölzchen};
  \reihe{7}{3}  \node at (8.5,-0.6) {\small 10 Hölzchen};
\end{tikzpicture}
\end{center}
\begin{parts}
  \part Wie viele Hölzchen braucht man für 10 Quadrate?
  \part Wie viele Hölzchen braucht man für $n$ Quadrate?
\end{parts}
\begin{solution}
  a) $31$ \qquad b) $3n + 1$ (oder gleichwertig, z.\,B.\ $4 + 3 \cdot (n-1)$)
  \diagnose{Gelingt a), aber nicht b), so kann ein Muster an Beispielen
  fortgesetzt, aber nicht verallgemeinert werden. Antworten wie $4n$ zeigen, dass
  die gemeinsamen Hölzchen nicht beachtet werden.}
\end{solution}
\raum{3.5cm}

\aufgabe{5.2}
Welcher Ausdruck ist grösser: $n + 2$ oder $2n$? Begründen Sie.
\begin{solution}
  Das hängt von $n$ ab: Für $n = 2$ sind beide gleich, für $n < 2$ ist $n + 2$
  grösser, für $n > 2$ ist $2n$ grösser.
  \diagnose{Die Antwort «$2n$, weil verdoppeln mehr ist» oder die Prüfung mit
  nur einem Wert zeigen, dass die Variable als feste Zahl und nicht als
  Platzhalter für viele Zahlen verstanden wird.}
\end{solution}
\raum{3cm}

% ---------------------------------------------------------------------------
% 6 Gleichungen
% ---------------------------------------------------------------------------
\aufgabe{6.1}
Welche Zahl gehört ins Kästchen?
\begin{parts}
  \part $8 + 4 = \kaestchen + 5$
  \part $3 \cdot \kaestchen = 24 \geteilt 2$
\end{parts}
\begin{solution}
  a) $7$ \qquad b) $4$
  \diagnose{Die Antworten $12$ oder $17$ bei a) zeigen, dass das
  Gleichheitszeichen als Aufforderung «jetzt kommt das Ergebnis» gelesen wird
  und nicht als Beziehung zwischen zwei gleichwertigen Seiten.}
\end{solution}

\aufgabe{6.2}
Ist die Aussage wahr oder falsch? Begründen Sie, möglichst ohne die Ergebnisse
auszurechnen.
\begin{parts}
  \part $37 + 25 = 36 + 26$ \hfill \wahlr{wahr}\quad \wahl{falsch}
  \part $47 - 18 = 48 - 17$ \hfill \wahl{wahr}\quad \wahlr{falsch}
  \part $6 \cdot 15 = 3 \cdot 30$ \hfill \wahlr{wahr}\quad \wahl{falsch}
\end{parts}
\begin{solution}
  a) wahr \quad b) falsch ($29 \neq 31$) \quad c) wahr
  \diagnose{Begründungen wie «links eins weniger, rechts eins mehr» zeigen
  relationales Denken. Wer b) als wahr einstuft, überträgt die
  Ausgleichsregel der Addition unreflektiert auf die Subtraktion.}
\end{solution}
\raum{3cm}

% ---------------------------------------------------------------------------
% 7 Funktionen und Graphen
% ---------------------------------------------------------------------------
\aufgabe{7.1}
Geben Sie die Koordinaten der Punkte $A$, $B$ und $C$ an.

\begin{center}
\begin{tikzpicture}[scale=0.55]
  \draw[gray!40, thin] (-5,-5) grid (5,5);
  \draw[->, thick] (-5.4,0) -- (5.6,0) node[right] {$x$};
  \draw[->, thick] (0,-5.4) -- (0,5.6) node[above] {$y$};
  \foreach \t in {-4,-2,2,4} {
    \node[below, font=\scriptsize] at (\t,0) {$\t$};
    \node[left,  font=\scriptsize] at (0,\t) {$\t$};
  }
  \foreach \p/\x/\y/\pos in {A/3/-2/right, B/-2/3/left, C/0/4/right} {
    \fill (\x,\y) circle (3.5pt);
    \node[\pos, font=\small] at (\x,\y) {$\p$};
  }
\end{tikzpicture}
\end{center}
$A = $ \fillin[$\pkt{3}{-2}$][2.5cm] \qquad
$B = $ \fillin[$\pkt{-2}{3}$][2.5cm] \qquad
$C = $ \fillin[$\pkt{0}{4}$][2.5cm]
\begin{solution}
  $A\pkt{3}{-2}$, \ $B\pkt{-2}{3}$, \ $C\pkt{0}{4}$
  \diagnose{Vertauschte Koordinaten bei $A$ und $B$ zeigen, dass die
  Reihenfolge $(x\,|\,y)$ nicht gesichert ist. Fehler bei $C$ (z.\,B.\ $(4\,|\,0)$)
  deuten auf Unsicherheit bei Punkten auf den Achsen hin.}
\end{solution}

\aufgabe{7.2}
Die $y$-Werte in der Tabelle werden alle nach derselben Regel aus den
$x$-Werten berechnet. Ergänzen Sie die Tabelle sinngemäss.
\begin{center}
  \renewcommand{\arraystretch}{1.4}
  \begin{tabular}{|c|*{4}{>{\centering\arraybackslash}p{1.2cm}|}}
    \hline
    $x$ & $0$ & $1$ & $3$  & $6$ \\ \hline
    $y$ & $2$ & $5$ & $11$ & \ifprintanswers $\mathbf{20}$\fi \\ \hline
  \end{tabular}
\end{center}
\begin{solution}
  $y = 20$ (Zuordnung $y = 3x + 2$)
  \diagnose{Da die $x$-Werte nicht in gleichen Schritten wachsen, scheitert ein
  rein schrittweises Weiterzählen (z.\,B.\ «immer $+6$» ergibt $17$). Das zeigt,
  ob der Zusammenhang zwischen $x$ und $y$ erkannt wird.}
\end{solution}
\raum{2.5cm}

% ---------------------------------------------------------------------------
% 8 Geometrie und Grössen
% ---------------------------------------------------------------------------
\aufgabe{8.1}
Ergänzen Sie.
\begin{parts}
  \part $3.5\text{ m} = $ \fillin[$350$][2cm] cm
  \part $250\text{ g} = $ \fillin[$0.25$][2cm] kg
  \part $1\text{ m}^2 = $ \fillin[$10\,000$][2cm] cm$^2$
  \part $2\text{ h }15\text{ min} = $ \fillin[$2.25$][2cm] h
\end{parts}
\begin{solution}
  a) $350$ \quad b) $0.25$ \quad c) $10\,000$ \quad d) $2.25$
  \diagnose{Die Antwort $100$ bei c) zeigt, dass die Umrechnungszahl für
  Flächeneinheiten nicht verstanden ist. Die Antwort $2.15$ bei d) zeigt, dass
  Zeitangaben wie Dezimalzahlen behandelt werden.}
\end{solution}

\aufgabe{8.2}
Ein Rechteck ist $6$\,cm lang und $4$\,cm breit.
\begin{parts}
  \part Berechnen Sie den Umfang und den Flächeninhalt.
  \part Hat jedes Rechteck mit demselben Umfang auch denselben Flächeninhalt?
        Begründen Sie.
\end{parts}
\begin{solution}
  a) Umfang $20$\,cm, Flächeninhalt $24$\,cm$^2$
  \qquad b) Nein, z.\,B.\ hat ein Rechteck mit $7$\,cm und $3$\,cm ebenfalls den
  Umfang $20$\,cm, aber den Flächeninhalt $21$\,cm$^2$.
  \diagnose{Verwechslungen von Umfang und Fläche (auch bei den Einheiten) sowie
  die Antwort «ja» bei b) zeigen, dass die beiden Grössen nicht klar
  unterschieden werden.}
\end{solution}
\raum{3.5cm}

% ---------------------------------------------------------------------------
% 9 Sprache ↔ Mathematik
% ---------------------------------------------------------------------------
\aufgabe{9.1}
Schreiben Sie die Rechnung auf und berechnen Sie das Ergebnis.
\begin{parts}
  \part die Summe von $12$ und $8$
  \part die Differenz von $20$ und $5$
  \part das Produkt von $6$ und $7$
  \part der Quotient von $36$ und $4$
\end{parts}
\begin{solution}
  a) $12 + 8 = 20$ \quad b) $20 - 5 = 15$ \quad c) $6 \cdot 7 = 42$ \quad
  d) $36 \geteilt 4 = 9$
  \diagnose{Unbekannte Fachbegriffe zeigen sich hier direkt. Bei b) und d)
  auf die Reihenfolge in der aufgeschriebenen Rechnung achten ($5 - 20$,
  $4 \geteilt 36$).}
\end{solution}
\raum{3cm}

\aufgabe{9.2}
\begin{parts}
  \part Mia hat dreimal so viele Karten wie Leo. Leo hat 12 Karten. Wie viele
        Karten hat Mia?
  \part Nora hat 12 Karten. Das sind dreimal so viele wie Tim. Wie viele Karten
        hat Tim?
\end{parts}
\begin{solution}
  a) $36$ \qquad b) $4$
  \diagnose{Die Antwort $36$ bei b) zeigt eine Schlüsselwort-Strategie
  («mal» bedeutet multiplizieren), ohne die Beziehung zwischen den Grössen zu
  erfassen.}
\end{solution}
\raum{3cm}

% ===========================================================================
\newpage
\teil{B}{Rechnen mit Zahlen}
% ===========================================================================

% ---------------------------------------------------------------------------
% 1 Zahlen und Stellenwert
% ---------------------------------------------------------------------------
\aufgabe{1.3}
\begin{parts}
  \part Runden Sie.
  \begin{subparts}
    \subpart $3.4567$ auf zwei Dezimalstellen
    \subpart $0.0496$ auf zwei Dezimalstellen
    \subpart $28\,349$ auf Tausender
  \end{subparts}
  \part Welches Ergebnis ist plausibel? Kreuzen Sie an, ohne genau zu rechnen.
        \[
          0.48 \cdot 205 \approx \quad
          \text{\wahl{$10$}} \quad \text{\wahlr{$100$}} \quad
          \text{\wahl{$1000$}} \quad \text{\wahl{$10\,000$}}
        \]
\end{parts}
\begin{solution}
  a) $3.46$; \ $0.05$; \ $28\,000$ \qquad b) $100$ (genau: $98.4$)
  \diagnose{$3.45$ zeigt Abschneiden statt Runden. Antworten wie $0.5$ oder
  $0.04$ bei $0.0496$ zeigen, dass der Stellenwert beim Runden verloren geht.
  Fehler bei b) deuten auf fehlende Überschlagsrechnung hin.}
\end{solution}
\raum{3cm}

\aufgabe{1.4}
\begin{parts}
  \part Ordnen Sie der Grösse nach, beginnend mit der kleinsten Zahl:
        \quad $-3.5; \quad -3; \quad 0.5; \quad -4; \quad -0.5$
  \part Welche Zahl ist um $5$ grösser als $-3$?
\end{parts}
\begin{solution}
  a) $-4 < -3.5 < -3 < -0.5 < 0.5$ \qquad b) $2$
  \diagnose{Wird $-3.5$ grösser als $-3$ eingeordnet, so werden negative Zahlen
  nach ihrem Betrag geordnet. Antworten wie $-8$ oder $8$ bei b) zeigen
  Unsicherheit beim Bewegen auf dem Zahlenstrahl.}
\end{solution}
\raum{2.5cm}

% ---------------------------------------------------------------------------
% 2 Operationen und Rechenreihenfolge
% ---------------------------------------------------------------------------
\aufgabe{2.3}
Berechnen Sie.
\begin{multicols}{3}
\begin{parts}
  \part $-7 + 12$
  \part $-4 - 9$
  \part $(-6) \cdot (-3)$
  \part $(-15) \geteilt 3$
  \part $5 - (-8)$
\end{parts}
\end{multicols}
\begin{solution}
  a) $5$ \quad b) $-13$ \quad c) $18$ \quad d) $-5$ \quad e) $13$
  \diagnose{Typische Fehler sind $-5$ oder $13$ bei b) sowie $-3$ bei e). Sie
  zeigen, dass Vorzeichen und Rechenzeichen nicht sicher unterschieden werden.}
\end{solution}
\raum{2.5cm}

\aufgabe{2.4}
Berechnen Sie.
\begin{multicols}{3}
\begin{parts}
  \part $12 - 3 \cdot 4$
  \part $(12 - 3) \cdot 4$
  \part $20 \geteilt 4 \cdot 5$
  \part $2 \cdot 3^2$
  \part $18 - 6 \geteilt 2 + 1$
\end{parts}
\end{multicols}
\begin{solution}
  a) $0$ \quad b) $36$ \quad c) $25$ \quad d) $18$ \quad e) $16$
  \diagnose{Typische Fehler: c) $1$ (Multiplikation vor Division), d) $36$
  (zuerst multipliziert, dann quadriert), e) $7$ (von links nach rechts
  gerechnet) oder $14$ (Addition vor Subtraktion). Sie zeigen, welche Regel der
  Rechenreihenfolge fehlt.}
\end{solution}
\raum{3cm}

% ---------------------------------------------------------------------------
% 3 Brüche
% ---------------------------------------------------------------------------
\aufgabe{3.3}
Berechnen Sie und kürzen Sie das Ergebnis so weit wie möglich.
\begin{multicols}{3}
\begin{parts}\bruchabstand
  \part $\dfrac{2}{3} + \dfrac{1}{4}$
  \part $\dfrac{5}{6} - \dfrac{1}{2}$
  \part $\dfrac{3}{4} \cdot \dfrac{2}{9}$
  \part $\dfrac{2}{5} \geteilt \dfrac{4}{15}$
  \part $3 \cdot \dfrac{2}{5}$
\end{parts}
\end{multicols}
\begin{solution}
  a) $\frac{11}{12}$ \quad b) $\frac{1}{3}$ \quad c) $\frac{1}{6}$ \quad
  d) $\frac{3}{2}$ \quad e) $\frac{6}{5}$
  \diagnose{$\frac{3}{7}$ bei a) zeigt, dass Zähler und Nenner getrennt addiert
  werden. $\frac{6}{15}$ bei e) zeigt, dass Zähler und Nenner mit $3$
  multipliziert werden (Verwechslung mit dem Erweitern). Fehler nur bei d)
  deuten auf eine nicht verstandene Divisionsregel hin.}
\end{solution}
\raum{4cm}

\aufgabe{3.4}
Ergänzen Sie die Tabelle.
\begin{center}
  \newcommand{\lsg}[1]{\ifprintanswers$\mathbf{#1}$\fi}
  \newcommand{\zh}{\rule[-2.4ex]{0pt}{6.6ex}}% Zeilenhöhe für Brüche
  \begin{tabular}{|c|c|c|}
    \hline
    \zh\makebox[2.6cm]{\textbf{Bruch}} & \makebox[2.6cm]{\textbf{Dezimalzahl}}
      & \makebox[2.6cm]{\textbf{Prozent}} \\ \hline
    \zh$\dfrac{3}{4}$          & \lsg{0.75}           & \lsg{75\,\%} \\ \hline
    \zh\lsg{\dfrac{1}{20}}     & $0.05$               & \lsg{5\,\%} \\ \hline
    \zh\lsg{\dfrac{5}{4}}      & \lsg{1.25}           & $125\,\%$ \\ \hline
    \zh$\dfrac{1}{3}$          & \lsg{0.333\ldots}    & \lsg{33.3\,\%} \\ \hline
  \end{tabular}
\end{center}
\begin{solution}
  Siehe Tabelle. Für $\frac{1}{3}$ sind auch $0.\overline{3}$ und
  $33\frac{1}{3}\,\%$ korrekt.
  \diagnose{$0.05 = 50\,\%$ zeigt Unsicherheit beim Stellenwert. Wer bei
  $125\,\%$ keinen Bruch findet, hält Prozentangaben über $100\,\%$ oder Brüche
  grösser als $1$ für unmöglich.}
\end{solution}
\raum{2.5cm}

% ---------------------------------------------------------------------------
% 4 Proportionalität und Prozent
% ---------------------------------------------------------------------------
\aufgabe{4.3}
Ein Pullover kostet CHF 80.
\begin{parts}
  \part Berechnen Sie $15\,\%$ von CHF 80.
  \part Der Preis des Pullovers wird um $25\,\%$ reduziert. Wie viel kostet er
        jetzt?
\end{parts}
\begin{solution}
  a) CHF 12 \qquad b) CHF 60
  \diagnose{Wer bei b) nur den Rabatt (CHF 20) angibt, unterscheidet nicht
  zwischen Prozentwert und neuem Wert. Fehler bei a) zeigen, dass die
  Grundidee von «Prozent als Hundertstel» fehlt.}
\end{solution}
\raum{3cm}

\aufgabe{4.4}
Ein Preis steigt von CHF 40 auf CHF 50.
\begin{parts}
  \part Um wie viel Prozent ist der Preis gestiegen?
  \part Danach sinkt der Preis wieder von CHF 50 auf CHF 40. Um wie viel Prozent
        ist er gesunken?
\end{parts}
\begin{solution}
  a) $25\,\%$ \qquad b) $20\,\%$
  \diagnose{Die Antwort $25\,\%$ bei b) (oder $20\,\%$ bei a)) zeigt, dass
  nicht erkannt wird, dass sich der Prozentsatz auf den jeweiligen Grundwert
  bezieht. Dieser Fehler ist sehr häufig.}
\end{solution}
\raum{3cm}

% ---------------------------------------------------------------------------
% 5 Terme
% ---------------------------------------------------------------------------
\aufgabe{5.3}
Vereinfachen Sie so weit wie möglich.
\begin{multicols}{3}
\begin{parts}
  \part $3a + 5 + 2a - 1$
  \part $4x + 3y - x + 2y$
  \part $2a + 3b$
  \part $3x \cdot 4x$
  \part $a^2 + a + a^2$
\end{parts}
\end{multicols}
\begin{solution}
  a) $5a + 4$ \quad b) $3x + 5y$ \quad c) nicht vereinfachbar \quad
  d) $12x^2$ \quad e) $2a^2 + a$
  \diagnose{$5ab$ bei c) sowie $a^4$, $3a^2$ oder $2a^3$ bei e) zeigen, dass
  ungleichartige Terme zusammengefasst werden. Die Aussage «$2a + 3b$ ist keine
  Lösung» weist auf die Vorstellung hin, dass ein Term ein einzelnes Ergebnis
  haben muss.}
\end{solution}
\raum{3cm}

\aufgabe{5.4}
Ist die Aussage richtig oder falsch? Kreuzen Sie an.
\begin{parts}
  \part $2a + 3b = 5ab$     \hfill \wahl{richtig}\quad \wahlr{falsch}
  \part $x + x = x^2$       \hfill \wahl{richtig}\quad \wahlr{falsch}
  \part $3x \cdot 2 = 6x$   \hfill \wahlr{richtig}\quad \wahl{falsch}
  \part $a + a^2 = a^3$     \hfill \wahl{richtig}\quad \wahlr{falsch}
  \part $4y - y = 4$        \hfill \wahl{richtig}\quad \wahlr{falsch}
\end{parts}
\begin{solution}
  a) falsch \quad b) falsch \quad c) richtig \quad d) falsch \quad e) falsch
  \diagnose{Wiedererkennungsversion von 5.3. Wer hier richtig entscheidet, bei
  5.3 aber Fehler macht, kennt die Regeln, wendet sie aber nicht sicher an.
  Bei e) zeigt «richtig», dass $y$ als Einheit oder Abkürzung gelesen wird
  («4 Äpfel minus 1 Apfel»).}
\end{solution}

% ---------------------------------------------------------------------------
% 6 Gleichungen
% ---------------------------------------------------------------------------
\aufgabe{6.3}
Lösen Sie die Gleichungen.
\begin{multicols}{3}
\begin{parts}\bruchabstand
  \part $x + 7 = 3$
  \part $3x = 21$
  \part $4x - 5 = 11$
  \part $12 = 5x + 2$
  \part $\dfrac{x}{3} = 6$
\end{parts}
\end{multicols}
\begin{solution}
  a) $x = -4$ \quad b) $x = 7$ \quad c) $x = 4$ \quad d) $x = 2$ \quad
  e) $x = 18$
  \diagnose{Bei d) steht die Variable rechts, was manche Schülerinnen und
  Schüler verunsichert. $x = 2$ bei e) zeigt, dass Umkehroperationen
  verwechselt werden.}
\end{solution}
\raum{3.5cm}

\aufgabe{6.4}
Lösen Sie die Gleichungen.
\begin{multicols}{3}
\begin{parts}
  \part $5x + 3 = 2x + 15$
  \part $7 - 2x = 1$
  \part $3x = 0$
\end{parts}
\end{multicols}
\begin{solution}
  a) $x = 4$ \quad b) $x = 3$ \quad c) $x = 0$
  \diagnose{$x = -3$ bei b) zeigt Vorzeichenfehler beim Umformen. Bei c)
  antworten manche «keine Lösung» oder «$x$ ist nicht definiert», weil $0$
  nicht als Lösung akzeptiert wird.}
\end{solution}
\raum{3.5cm}

% ---------------------------------------------------------------------------
% 7 Funktionen und Graphen
% ---------------------------------------------------------------------------
\aufgabe{7.3}
Der Graph zeigt, wie weit eine Person von zu Hause entfernt ist.

\begin{center}
\begin{tikzpicture}[x=0.22cm, y=1.1cm]
  \draw[gray!35, thin, xstep=5, ystep=0.5] (0,0) grid (45,3.5);
  \draw[->, thick] (0,0) -- (47,0) node[right] {\small Zeit (min)};
  \draw[->, thick] (0,0) -- (0,3.9) node[above] {\small Entfernung (km)};
  \foreach \t in {0,10,20,30,40} \node[below, font=\scriptsize] at (\t,0) {$\t$};
  \foreach \s in {1,2,3} \node[left, font=\scriptsize] at (0,\s) {$\s$};
  \draw[very thick] (0,0) -- (20,3) -- (30,3) -- (40,0);
\end{tikzpicture}
\end{center}
\begin{parts}
  \part Wie weit ist die Person nach 20 Minuten von zu Hause entfernt?
  \part Was geschieht zwischen der 20.\ und der 30.\ Minute?
  \part In welchem Abschnitt bewegt sich die Person am schnellsten?
        Begründen Sie.
  \part Berechnen Sie die Geschwindigkeit im ersten Abschnitt in km/h.
\end{parts}
\begin{solution}
  a) $3$\,km \quad b) Die Person bleibt stehen (macht eine Pause). \quad
  c) Auf dem Rückweg (30.--40.\ Minute), weil der Graph dort am steilsten ist.
  \quad d) $3$\,km in $20$\,min, also $9$\,km/h.
  \diagnose{Wird bei c) der «höchste» Abschnitt gewählt, so wird die Höhe des
  Graphen statt seiner Steigung gelesen. Bei b) zeigt «sie geht geradeaus
  weiter», dass der Graph als Bild des Weges gedeutet wird.}
\end{solution}
\raum{3.5cm}

\aufgabe{7.4}
Eine Velofahrerin fährt über einen Hügel: zuerst bergauf, dann bergab.
Welcher Graph zeigt ihre Geschwindigkeit $v$ in Abhängigkeit von der Zeit $t$
am besten? Kreuzen Sie an.

\begin{center}
\newcommand{\achsen}{%
  \draw[->] (0,0) -- (2.4,0) node[below] {\scriptsize $t$};
  \draw[->] (0,0) -- (0,1.7) node[left]  {\scriptsize $v$};}
\begin{tikzpicture}[thick]
  \begin{scope}[xshift=0cm]
    \achsen \draw[very thick] (0,0.3) .. controls (0.8,1.6) and (1.4,1.6) .. (2.2,0.3);
    \node at (1.1,-0.5) {\wahl{Graph 1}};
  \end{scope}
  \begin{scope}[xshift=3.3cm]
    \achsen \draw[very thick] (0,1.0) .. controls (0.6,0.3) and (0.9,0.3) .. (1.1,0.35)
                                     .. controls (1.4,0.5) and (1.8,1.3) .. (2.2,1.5);
    \node at (1.1,-0.5) {\wahlr{Graph 2}};
  \end{scope}
  \begin{scope}[xshift=6.6cm]
    \achsen \draw[very thick] (0,1.4) -- (2.2,0.3);
    \node at (1.1,-0.5) {\wahl{Graph 3}};
  \end{scope}
  \begin{scope}[xshift=9.9cm]
    \achsen \draw[very thick] (0,0.3) -- (1.1,1.3) -- (2.2,1.3);
    \node at (1.1,-0.5) {\wahl{Graph 4}};
  \end{scope}
\end{tikzpicture}
\end{center}
\begin{solution}
  Graph 2: Bergauf nimmt die Geschwindigkeit ab, bergab nimmt sie zu.
  \diagnose{Die Wahl von Graph 1 zeigt die klassische Fehlvorstellung «Graph
  als Bild der Situation»: Der Graph sieht aus wie der Hügel, beschreibt aber
  nicht die Geschwindigkeit.}
\end{solution}

% ---------------------------------------------------------------------------
% 8 Geometrie und Grössen
% ---------------------------------------------------------------------------
\aufgabe{8.3}
Aus einem Rechteck wird an einer Ecke ein kleineres Rechteck herausgeschnitten.
Berechnen Sie den Flächeninhalt und den Umfang der Figur.

\begin{center}
\begin{tikzpicture}[scale=0.5]
  \draw[very thick, fill=gray!10] (0,0) -- (8,0) -- (8,4) -- (5,4) -- (5,6) -- (0,6) -- cycle;
  \draw[dashed, gray] (5,4) -- (5,6) (5,6) -- (8,6) -- (8,4);
  \node[below] at (4,0) {\small $8$\,cm};
  \node[left]  at (0,3) {\small $6$\,cm};
  \node[above] at (6.5,6) {\small $3$\,cm};
  \node[right] at (8,5) {\small $2$\,cm};
\end{tikzpicture}
\end{center}
\begin{solution}
  Flächeninhalt: $8 \cdot 6 - 3 \cdot 2 = 42$\,cm$^2$ \qquad
  Umfang: $28$\,cm (gleich wie beim ganzen Rechteck)
  \diagnose{Wer beim Umfang etwas abzieht (z.\,B.\ $28 - 10 = 18$\,cm),
  überträgt die Überlegung zur Fläche auf den Umfang. Das zeigt, dass Umfang
  und Fläche nicht als verschiedene Grössen verstanden werden.}
\end{solution}
\raum{3.5cm}

\aufgabe{8.4}
\begin{parts}
  \part In einem Dreieck sind zwei Winkel $47^\circ$ und $68^\circ$ gross. Wie
        gross ist der dritte Winkel?
  \part Ein Aquarium ist $50$\,cm lang, $40$\,cm breit und $30$\,cm hoch. Wie
        viele Liter Wasser fasst es?
\end{parts}
\begin{solution}
  a) $65^\circ$ \qquad b) $60\,000$\,cm$^3 = 60$\,dm$^3 = 60$\,Liter
  \diagnose{Ergebnisse wie $600$\,L oder $6$\,L bei b) zeigen, dass der
  Zusammenhang $1\,\text{dm}^3 = 1$\,L bzw.\ die Umrechnung von
  Volumeneinheiten nicht gesichert ist.}
\end{solution}
\raum{3cm}

% ---------------------------------------------------------------------------
% 9 Sprache ↔ Mathematik
% ---------------------------------------------------------------------------
\aufgabe{9.3}
Schreiben Sie als \textbf{einen} Rechenausdruck und berechnen Sie das Ergebnis.
\begin{parts}
  \part das Dreifache der Differenz von $12$ und $5$
  \part die Differenz aus dem Dreifachen von $12$ und $5$
  \part Subtrahieren Sie $4$ von $10$.
  \part die Hälfte der Summe von $9$ und $15$
\end{parts}
\begin{solution}
  a) $3 \cdot (12 - 5) = 21$ \quad b) $3 \cdot 12 - 5 = 31$ \quad
  c) $10 - 4 = 6$ \quad d) $(9 + 15) \geteilt 2 = 12$
  \diagnose{a) und b) enthalten dieselben Wörter in anderer Reihenfolge und
  ergeben verschiedene Strukturen. Fehlende Klammern bei a) oder d) zeigen,
  dass die Struktur des Satzes nicht in die Rechnung übertragen wird. $4 - 10$
  bei c) ist die typische Umkehrung der Reihenfolge.}
\end{solution}
\raum{3.5cm}

\aufgabe{9.4}
Eine Klasse mit 24 Schülerinnen und Schülern macht einen Ausflug. Die Bahnfahrt
kostet pro Person CHF 18.50, der Museumseintritt pro Person CHF 12. Die Klasse
fährt um 8:15 Uhr ab und ist um 16:40 Uhr zurück. Die Schule übernimmt
CHF 300 der Gesamtkosten. Wie viel muss jede Person selbst bezahlen?
\begin{solution}
  Kosten pro Person: $18.50 + 12 = 30.50$; Gesamtkosten:
  $24 \cdot 30.50 = 732$; nach Abzug: $732 - 300 = 432$; pro Person:
  $432 \geteilt 24 = 18$. Jede Person bezahlt CHF 18.
  \diagnose{Die Zeitangaben sind überflüssig. Wer sie verwendet, kann die
  relevanten Informationen nicht auswählen. Wird CHF 300 von den Kosten einer
  Person abgezogen, so fehlt die Unterscheidung zwischen Gesamt- und
  Einzelgrösse.}
\end{solution}
\raum{4.5cm}

% ===========================================================================
\newpage
\teil{C}{Rechnen mit Variablen}
% ===========================================================================

% ---------------------------------------------------------------------------
% 1 Zahlen und Stellenwert
% ---------------------------------------------------------------------------
\aufgabe{1.5}
\begin{parts}
  \part Schreiben Sie in wissenschaftlicher Schreibweise:
        \quad $45\,000\,000$ \qquad $0.00072$
  \part Schreiben Sie als Dezimalzahl:
        \quad $3.1 \cdot 10^5$ \qquad $6 \cdot 10^{-3}$
  \part Berechnen Sie die beiden Ausdrücke und geben Sie die Ergebnisse in
        wissenschaftlicher Schreibweise bzw.\ als Zehnerpotenz an.
  \begin{multicols}{2}
  \begin{subparts}
    \subpart $(2 \cdot 10^3) \cdot (4 \cdot 10^5)$
    \subpart $10^6 \geteilt 10^2$
  \end{subparts}
  \end{multicols}
\end{parts}
\begin{solution}
  a) $4.5 \cdot 10^7$; \ $7.2 \cdot 10^{-4}$ \quad
  b) $310\,000$; \ $0.006$ \quad
  c) $8 \cdot 10^8$; \ $10^4$
  \diagnose{Falsch gezählte Stellen zeigen fehlendes Stellenwertverständnis.
  Werden negative Exponenten als negative Zahlen gedeutet ($6 \cdot 10^{-3}
  = -6000$), fehlt die Bedeutung der Schreibweise. $10^{15}$ bzw.\ $10^3$ bei
  c) zeigen, dass Exponenten multipliziert bzw.\ dividiert statt addiert bzw.\
  subtrahiert werden.}
\end{solution}
\raum{3.5cm}

% ---------------------------------------------------------------------------
% 2 Operationen und Rechenreihenfolge
% ---------------------------------------------------------------------------
\aufgabe{2.5}
Berechnen Sie den Wert der Terme für $a = -2$ und $b = 3$.
\begin{multicols}{3}
\begin{parts}
  \part $3a + b$
  \part $a^2$
  \part $-a^2$
  \part $2b - a$
  \part $a \cdot b^2$
\end{parts}
\end{multicols}
\begin{solution}
  a) $-3$ \quad b) $4$ \quad c) $-4$ \quad d) $8$ \quad e) $-18$
  \diagnose{Eine der aussagekräftigsten Aufgaben. $-4$ bei b) zeigt, dass ohne
  Klammern eingesetzt wird ($-2^2$ statt $(-2)^2$). $4$ bei c) zeigt
  Unsicherheit mit dem Minuszeichen vor der Potenz. $36$ bei e) zeigt, dass
  $(ab)^2$ gerechnet wird.}
\end{solution}
\raum{3cm}

% ---------------------------------------------------------------------------
% 3 Brüche
% ---------------------------------------------------------------------------
\aufgabe{3.5}
Fassen Sie zusammen und vereinfachen Sie so weit wie möglich.
\begin{multicols}{2}
\begin{parts}\bruchabstand
  \part $\dfrac{a}{4} + \dfrac{b}{6}$
  \part $\dfrac{x}{3} - \dfrac{x}{5}$
  \part $\dfrac{3}{x} + \dfrac{2}{x}$
  \part $\dfrac{2}{a} \cdot \dfrac{a}{6}$
\end{parts}
\end{multicols}
\begin{solution}
  a) $\frac{3a + 2b}{12}$ \quad b) $\frac{2x}{15}$ \quad c) $\frac{5}{x}$ \quad
  d) $\frac{1}{3}$
  \diagnose{Diese Aufgabe entspricht 3.3 mit Variablen. Gelingt 3.3, aber nicht
  3.5, so liegt die Schwierigkeit beim Übergang zur Algebra, nicht bei den
  Bruchregeln. $\frac{5}{2x}$ bei c) oder $\frac{a+b}{10}$ bei a) zeigen
  getrenntes Addieren von Zählern und Nennern.}
\end{solution}
\raum{4cm}

% ---------------------------------------------------------------------------
% 4 Proportionalität und Prozent
% ---------------------------------------------------------------------------
\aufgabe{4.5}
\begin{parts}
  \part Entscheiden Sie bei jeder Tabelle, ob die Zuordnung proportional,
        umgekehrt proportional oder keines von beidem ist. Geben Sie bei
        proportionalen Zuordnungen die Gleichung $y = k \cdot x$ an.
  \begin{center}
    \renewcommand{\arraystretch}{1.3}
    \newcommand{\tab}[5]{%
      \begin{tabular}[t]{@{}c@{}}
        \small Tabelle #1 \\[0.5ex]
        \begin{tabular}{|c|c|c|c|c|}
          \hline $x$ & 1 & 2 & 4 & 5 \\ \hline
                 $y$ & #2 & #3 & #4 & #5 \\ \hline
        \end{tabular}
      \end{tabular}}
    \tab{1}{3}{6}{12}{15} \qquad \tab{2}{5}{7}{11}{13} \qquad \tab{3}{20}{10}{5}{4}
  \end{center}
  \part $4$ Personen brauchen für eine Arbeit $6$ Tage. Wie lange brauchen
        $3$ Personen bei gleichem Arbeitstempo?
\end{parts}
\begin{solution}
  a) Tabelle 1: proportional, $y = 3x$; \ Tabelle 2: keines von beidem
  ($y = 2x + 3$); \ Tabelle 3: umgekehrt proportional ($x \cdot y = 20$)
  \qquad b) $8$ Tage
  \diagnose{Tabelle 2 ist die zentrale Falle: «gleichmässig wachsend» wird mit
  proportional verwechselt. $4.5$ Tage bei b) zeigen, dass ohne Prüfung direkt
  proportional gerechnet wird.}
\end{solution}
\raum{3.5cm}

% ---------------------------------------------------------------------------
% 5 Terme
% ---------------------------------------------------------------------------
\aufgabe{5.5}
\begin{parts}
  \part Multiplizieren Sie aus und vereinfachen Sie.
  \begin{multicols}{4}
    $3 \cdot (x + 4)$ \columnbreak

    $(-2) \cdot (a - 5)$ \columnbreak

    $(x + 3)(x + 2)$ \columnbreak

    $(a + b)^2$
  \end{multicols}
  \part Klammern Sie so weit wie möglich aus.
        \qquad $6x + 9$ \qquad\qquad $x^2 - 5x$
\end{parts}
\begin{solution}
  a) $3x + 12$; \ $-2a + 10$; \ $x^2 + 5x + 6$; \ $a^2 + 2ab + b^2$ \quad
  b) $3 \cdot (2x + 3)$; \ $x(x - 5)$
  \diagnose{$3x + 4$ zeigt, dass nur der erste Summand multipliziert wird.
  $-2a - 10$ zeigt Vorzeichenfehler. $(a + b)^2 = a^2 + b^2$ ist der häufigste
  Fehler überhaupt und zeigt, dass Potenzieren über die Addition «verteilt»
  wird.}
\end{solution}
\raum{4cm}

% ---------------------------------------------------------------------------
% 6 Gleichungen
% ---------------------------------------------------------------------------
\aufgabe{6.5}
\begin{parts}
  \part Lösen Sie die Gleichungen.
  \begin{multicols}{3}
    $3 \cdot (x - 2) = 2x + 1$ \columnbreak

    $\dfrac{x}{2} + \dfrac{x}{3} = 10$ \columnbreak

    $\dfrac{x + 1}{4} = \dfrac{x - 2}{3}$
  \end{multicols}
  \part Prüfen Sie, \textbf{ohne die Gleichung zu lösen}, ob $x = 2$ eine Lösung
        der folgenden Gleichung ist.
        \[ 5x - 3 = 2 \cdot (x + 1) + 1 \]
\end{parts}
\begin{solution}
  a) $x = 7$; \ $x = 12$; \ $x = 11$ \qquad
  b) Ja: Für $x = 2$ ergeben beide Seiten $7$.
  \diagnose{Teil b) prüft, ob verstanden wird, was eine Lösung ist, unabhängig
  vom Lösungsverfahren. Wer b) nur durch Lösen der Gleichung beantworten kann,
  verwendet ein Verfahren ohne Begriffsverständnis.}
\end{solution}
\raum{4.5cm}

% ---------------------------------------------------------------------------
% 7 Funktionen und Graphen
% ---------------------------------------------------------------------------
\aufgabe{7.5}
\begin{parts}
  \part Geben Sie für die Gerade im linken Koordinatensystem die Steigung, den
        $y$-Achsenabschnitt und die Funktionsgleichung an.
  \part Zeichnen Sie den Graphen von $y = -\frac{1}{2}x + 3$ ins rechte
        Koordinatensystem.
\end{parts}

\begin{center}
\newcommand{\koord}[4]{% xmin, xmax, ymin, ymax
  \draw[gray!35, thin] (#1,#3) grid (#2,#4);
  \draw[->, thick] (#1-0.4,0) -- (#2+0.5,0) node[right] {\small $x$};
  \draw[->, thick] (0,#3-0.4) -- (0,#4+0.5) node[above] {\small $y$};
  \foreach \t in {#1,...,#2} {\ifnum\t=0\else
    \draw (\t,0.1) -- (\t,-0.1) node[below, font=\tiny] {$\t$};\fi}
  \foreach \t in {#3,...,#4} {\ifnum\t=0\else
    \draw (0.1,\t) -- (-0.1,\t) node[left, font=\tiny] {$\t$};\fi}}
\begin{tikzpicture}[scale=0.5]
  \koord{-3}{4}{-3}{7}
  \draw[very thick, domain=-2:3] plot (\x, {2*\x + 1});
  \fill (0,1) circle (4pt) (2,5) circle (4pt);
\end{tikzpicture}
\hspace{1.5cm}
\begin{tikzpicture}[scale=0.5]
  \koord{-2}{8}{-2}{6}
  \ifprintanswers
    \draw[very thick, domain=-2:8] plot (\x, {-0.5*\x + 3});
  \fi
\end{tikzpicture}
\end{center}
\begin{parts}
  \setcounter{partno}{2}
  \part Liegt der Punkt $\pkt{4}{7}$ auf der Geraden $y = 2x - 1$? Begründen Sie.
  \part Die Werte in der Tabelle gehören zu einer linearen Funktion.
        Bestimmen Sie die Funktionsgleichung.
  \quad
  \begin{tabular}{|c|c|c|c|}
    \hline $x$ & $0$ & $1$ & $2$  \\ \hline
           $y$ & $4$ & $1$ & $-2$ \\ \hline
  \end{tabular}
\end{parts}
\begin{solution}
  a) Steigung $2$, $y$-Achsenabschnitt $1$, $y = 2x + 1$ \quad
  b) Gerade durch $\pkt{0}{3}$ und $\pkt{6}{0}$ (siehe Grafik) \quad
  c) Ja, denn $2 \cdot 4 - 1 = 7$. \quad
  d) $y = -3x + 4$
  \diagnose{Die Aufgabe prüft alle Wechsel zwischen Graph, Gleichung und
  Tabelle. Typische Fehler: Steigung als $\frac{\Delta x}{\Delta y}$ berechnet
  ($\frac{1}{2}$), negative Steigung steigend gezeichnet, $y$-Achsenabschnitt
  und Nullstelle verwechselt.}
\end{solution}
\raum{3.5cm}

% ---------------------------------------------------------------------------
% 8 Geometrie und Grössen
% ---------------------------------------------------------------------------
\aufgabe{8.5}
\begin{parts}
  \part Ein rechtwinkliges Dreieck hat die Katheten $6$\,cm und $8$\,cm.
        Berechnen Sie die Länge der Hypotenuse.
  \part In einem rechtwinkligen Dreieck ist die Hypotenuse $13$\,cm und eine
        Kathete $5$\,cm lang. Berechnen Sie die Länge der anderen Kathete.
  \part Eine $10$\,m lange Leiter lehnt an einer senkrechten Wand. Ihr Fuss
        steht $6$\,m von der Wand entfernt. In welcher Höhe berührt die Leiter
        die Wand?
  \part Ist ein Dreieck mit den Seitenlängen $7$\,cm, $24$\,cm und $25$\,cm
        rechtwinklig? Begründen Sie.
\end{parts}
\begin{solution}
  a) $10$\,cm \quad b) $12$\,cm \quad c) $8$\,m \quad
  d) Ja, denn $7^2 + 24^2 = 49 + 576 = 625 = 25^2$.
  \diagnose{$\sqrt{194}$ bei b) zeigt, dass $a^2 + b^2 = c^2$ mechanisch
  angewendet wird, ohne die Hypotenuse zu identifizieren. c) verlangt, das
  rechtwinklige Dreieck in einer Situation zu erkennen. d) prüft die Umkehrung
  des Satzes von Pythagoras.}
\end{solution}
\raum{5cm}

% ---------------------------------------------------------------------------
% 9 Sprache ↔ Mathematik
% ---------------------------------------------------------------------------
\aufgabe{9.5}
\begin{parts}
  \part Schreiben Sie als Term: \quad «$5$ weniger als das Doppelte von $x$»
        \quad und \quad «um $3$ grösser als die Hälfte von $y$»
  \part An einer Schule gibt es sechsmal so viele Schülerinnen und Schüler wie
        Lehrpersonen. Schreiben Sie diese Aussage als Gleichung. Verwenden Sie
        $S$ für die Anzahl Schülerinnen und Schüler und $L$ für die Anzahl
        Lehrpersonen.
  \part Beschreiben Sie die Terme in Worten: \quad $3 \cdot (x + 4)$ \quad und \quad
        $3x + 4$
  \part Verdoppelt man eine Zahl und addiert $7$, so erhält man $31$. Stellen
        Sie eine Gleichung auf und bestimmen Sie die Zahl.
\end{parts}
\begin{solution}
  a) $2x - 5$; \ $\frac{y}{2} + 3$ \quad
  b) $S = 6L$ \quad
  c) «das Dreifache der Summe von $x$ und $4$»; «die Summe aus dem Dreifachen
  von $x$ und $4$» \quad
  d) $2x + 7 = 31$, $x = 12$
  \diagnose{b) ist die klassische «Studierende-und-Professoren»-Aufgabe: $6S = L$
  entsteht durch Wort-für-Wort-Übersetzung, ohne die Beziehung zu erfassen.
  $5 - 2x$ bei a) zeigt die Umkehrung bei «weniger als». c) ist die
  Gegenrichtung zu 9.3.}
\end{solution}
\raum{4.5cm}

% ===========================================================================
\newpage
\teil{D}{Vertiefung}
% ===========================================================================

% ---------------------------------------------------------------------------
% 1 Zahlen und Stellenwert
% ---------------------------------------------------------------------------
\aufgabe{1.6}
\begin{parts}
  \part Geben Sie für jede Zahl alle Zahlenmengen ($\mathbb{N}$, $\mathbb{Z}$,
        $\mathbb{Q}$, $\mathbb{R}$) an, zu denen sie gehört.
        \[
          -3 \qquad 0.75 \qquad \sqrt{2} \qquad \frac{12}{4} \qquad 0.333\ldots
        \]
  \part Kann das Ergebnis stimmen? Entscheiden Sie \textbf{ohne nachzurechnen}
        und begründen Sie.
        \begin{itemize}
          \item Lena rechnet: $0.6 \geteilt 0.03 = 2$.
          \item Ein Schüler berechnet das Volumen eines Schulzimmers und erhält
                $18$\,m$^3$.
        \end{itemize}
\end{parts}
\begin{solution}
  a) $-3$: $\mathbb{Z}, \mathbb{Q}, \mathbb{R}$; \
  $0.75$: $\mathbb{Q}, \mathbb{R}$; \
  $\sqrt{2}$: $\mathbb{R}$; \
  $\frac{12}{4}$: $\mathbb{N}, \mathbb{Z}, \mathbb{Q}, \mathbb{R}$; \
  $0.333\ldots$: $\mathbb{Q}, \mathbb{R}$ \\
  b) Beides ist nicht plausibel: Division durch eine Zahl kleiner als $1$
  vergrössert (richtig: $20$). Ein Schulzimmer hat etwa $150$ bis
  $200$\,m$^3$.
  \diagnose{Wird $\frac{12}{4}$ nur $\mathbb{Q}$ zugeordnet, so wird die Zahl
  nach ihrer Schreibweise statt nach ihrem Wert beurteilt. Die Zahl $0$ wurde
  bewusst vermieden, da Lehrmittel uneinheitlich sind, ob $0 \in \mathbb{N}$.}
\end{solution}
\raum{4cm}

% ---------------------------------------------------------------------------
% 2 Operationen und Rechenreihenfolge
% ---------------------------------------------------------------------------
\aufgabe{2.6}
\begin{parts}
  \part Welche Terme sind für alle Werte von $x$ gleichwertig mit $2 \cdot (x - 3)$?
        Kreuzen Sie an.

        \wahl{$2x - 3$} \quad \wahlr{$2x - 6$} \quad \wahlr{$(-2) \cdot (3 - x)$}
        \quad \wahlr{$(x - 3) \cdot 2$} \quad \wahl{$2 \cdot (3 - x)$}
  \part Gilt $(a - b)^2 = (b - a)^2$ für alle Zahlen $a$ und $b$? Begründen Sie.
\end{parts}
\begin{solution}
  a) $2x - 6$, \ $(-2) \cdot (3 - x)$, \ $(x - 3) \cdot 2$ \qquad
  b) Ja, denn $b - a = -(a - b)$, und beim Quadrieren fällt das Vorzeichen weg.
  \diagnose{Die Aufgabe prüft, ob Gleichwertigkeit an der Struktur erkannt wird.
  Wer nur einzelne Werte einsetzt, kann b) nicht allgemein begründen.}
\end{solution}
\raum{3cm}

% ---------------------------------------------------------------------------
% 3 Brüche
% ---------------------------------------------------------------------------
\aufgabe{3.6}
\begin{parts}
  \part Fassen Sie zusammen: \quad $\dfrac{1}{x} + \dfrac{1}{x + 1}$
  \part Kürzen Sie, falls möglich:
        \quad $\dfrac{3x + 6}{x + 2}$ \qquad $\dfrac{x^2 - 9}{x - 3}$
        \qquad $\dfrac{x + 5}{5}$
\end{parts}
\begin{solution}
  a) $\frac{2x + 1}{x(x + 1)}$ \quad
  b) $3$; \ $x + 3$; \ nicht kürzbar
  \diagnose{$\frac{x + 5}{5} = x + 1$ oder $= x$ zeigt «Kürzen aus Summen». Wer
  bei den ersten beiden Brüchen nicht faktorisiert, erkennt die Struktur nicht.
  Die Angabe der Definitionslücken ($x \neq -2$, $x \neq 3$) wird auf diesem
  Niveau nicht verlangt.}
\end{solution}
\raum{3.5cm}

% ---------------------------------------------------------------------------
% 4 Proportionalität und Prozent
% ---------------------------------------------------------------------------
\aufgabe{4.6}
\begin{parts}
  \part Ein Preis wird zuerst um $20\,\%$ erhöht und danach um $20\,\%$
        gesenkt. Ist er danach gleich hoch wie zu Beginn? Begründen Sie.
  \part Nach einer Erhöhung um $25\,\%$ kostet ein Artikel CHF 150. Wie viel
        kostete er vorher?
  \part Eine Bevölkerung von $P$ Personen wächst jährlich um $3\,\%$. Geben Sie
        einen Term für die Bevölkerung nach $n$ Jahren an.
\end{parts}
\begin{solution}
  a) Nein: $1.2 \cdot 0.8 = 0.96$, der Preis ist $4\,\%$ tiefer. \quad
  b) $150 \geteilt 1.25 = 120$, also CHF 120 \quad
  c) $P \cdot 1.03^n$
  \diagnose{CHF 112.50 bei b) zeigt, dass $25\,\%$ vom neuen statt vom alten
  Wert abgezogen werden. $P + 0.03n$ oder $P \cdot 0.03^n$ bei c) zeigen, dass
  der Wachstumsfaktor nicht verstanden ist.}
\end{solution}
\raum{4cm}

% ---------------------------------------------------------------------------
% 5 Terme
% ---------------------------------------------------------------------------
\aufgabe{5.6}
\begin{parts}
  \part Faktorisieren Sie so weit wie möglich.
  \begin{multicols}{4}
    $x^2 - 9$ \columnbreak

    $x^2 + 6x + 9$ \columnbreak

    $4a^2b - 6ab^2$ \columnbreak

    $ax + ay - bx - by$
  \end{multicols}
  \part Ergänzen Sie: \quad $x^2 + \kaestchen\,x + 25 = (x + \kaestchen)^2$
\end{parts}
\begin{solution}
  a) $(x - 3)(x + 3)$; \ $(x + 3)^2$; \ $2ab(2a - 3b)$; \ $(a - b)(x + y)$
  \quad b) $10$ und $5$
  \diagnose{Faktorisieren verlangt, die binomischen Formeln «rückwärts» zu
  lesen. Wer 5.5 lösen kann, hier aber scheitert, kennt die Formeln nur in eine
  Richtung.}
\end{solution}
\raum{3.5cm}

% ---------------------------------------------------------------------------
% 6 Gleichungen
% ---------------------------------------------------------------------------
\aufgabe{6.6}
\begin{parts}
  \part Lösen Sie nach der angegebenen Variable auf.
  \begin{multicols}{3}
    $A = \frac{1}{2} \cdot g \cdot h$ \ nach $h$ \columnbreak

    $v = \dfrac{s}{t}$ \ nach $t$ \columnbreak

    $F = \frac{9}{5} \cdot C + 32$ \ nach $C$
  \end{multicols}
  \part Lösen Sie nach $x$ auf ($a \neq 3$): \quad $ax = 3x + 5$
  \part Bestimmen Sie die Lösungsmenge.
        \quad $2 \cdot (x + 1) = 2x + 2$ \qquad $2 \cdot (x + 1) = 2x + 3$
\end{parts}
\begin{solution}
  a) $h = \frac{2A}{g}$; \ $t = \frac{s}{v}$; \ $C = \frac{5}{9} \cdot (F - 32)$
  \quad b) $x = \frac{5}{a - 3}$ \quad
  c) $\mathbb{L} = \mathbb{R}$; \ $\mathbb{L} = \{\,\}$
  \diagnose{Bei b) muss $x$ ausgeklammert werden ($x(a - 3) = 5$). Wer nur
  Gleichungen mit Zahlen lösen kann, kommt hier nicht weiter. Bei c) zeigt
  die Reaktion «$0 = 1$, also habe ich mich verrechnet», dass Lösungsmengen
  ohne genau eine Lösung unbekannt sind.}
\end{solution}
\raum{4.5cm}

% ---------------------------------------------------------------------------
% 7 Funktionen und Graphen
% ---------------------------------------------------------------------------
\aufgabe{7.6}
\begin{parts}
  \part Was haben alle Geraden mit der Gleichung $y = mx + 2$ gemeinsam? Was
        haben alle Geraden mit der Gleichung $y = 3x + q$ gemeinsam?
  \part Berechnen Sie den Schnittpunkt der Geraden $y = 2x + 1$ und
        $y = -x + 7$.
  \part Für welchen Wert von $m$ verläuft die Gerade $y = mx + 2$ durch den
        Punkt $\pkt{4}{0}$?
\end{parts}
\begin{solution}
  a) Alle verlaufen durch $\pkt{0}{2}$; alle sind parallel (Steigung $3$).
  \quad b) $2x + 1 = -x + 7 \Rightarrow x = 2$, also $S\pkt{2}{5}$
  \quad c) $0 = 4m + 2 \Rightarrow m = -\frac{1}{2}$
  \diagnose{Die Aufgabe prüft, ob ein Buchstabe als Parameter (fest, aber
  beliebig) und nicht als Variable verstanden wird, und ob ein Schnittpunkt als
  Lösung einer Gleichung erkannt wird.}
\end{solution}
\raum{4cm}

% ---------------------------------------------------------------------------
% 8 Geometrie und Grössen
% ---------------------------------------------------------------------------
\aufgabe{8.6}
\begin{parts}
  \part Die Geraden $AA'$ und $BB'$ sind parallel. Berechnen Sie die Längen
        $\overline{SB'}$ und $\overline{BB'}$. (Skizze nicht massstäblich)

  \begin{center}
  \begin{tikzpicture}[scale=0.8]
    \coordinate (S) at (0,0);
    \coordinate (A) at (4,0);   \coordinate (B) at (6,0);
    \coordinate (A2) at (3.2,2.4); \coordinate (B2) at (4.8,3.6);
    \draw[thick] (S) -- (7,0);
    \draw[thick] (S) -- (5.6,4.2);
    \draw[thick] (A) -- (A2);
    \draw[thick] (B) -- (B2);
    \foreach \p/\pos in {S/left, A/below, B/below, A2/above left, B2/above left}
      \fill (\p) circle (2pt);
    \node[left]        at (S)  {$S$};
    \node[below]       at (A)  {$A$};
    \node[below]       at (B)  {$B$};
    \node[above left]  at (A2) {$A'$};
    \node[above left]  at (B2) {$B'$};
    \node[below, font=\small] at (2,0)       {$4$\,cm};
    \node[below, font=\small] at (5,0)       {$2$\,cm};
    \node[above left, font=\small] at (1.6,1.2) {$5$\,cm};
    \node[right, font=\small] at (3.6,1.2)   {$3$\,cm};
  \end{tikzpicture}
  \end{center}
  \part Eine $1.80$\,m grosse Person wirft einen $1.20$\,m langen Schatten. Ein
        Baum wirft gleichzeitig einen $8$\,m langen Schatten. Wie hoch ist der
        Baum?
  \part Ein Quadrat mit der Seitenlänge $a$ wird mit dem Faktor $k$ vergrössert.
        Geben Sie den Umfang und den Flächeninhalt des neuen Quadrats an.
\end{parts}
\begin{solution}
  a) $\overline{SB} = 6$\,cm, also $\overline{SB'} = 5 \cdot \frac{6}{4} =
  7.5$\,cm und $\overline{BB'} = 3 \cdot \frac{6}{4} = 4.5$\,cm \quad
  b) $12$\,m \quad
  c) Umfang $4ka$, Flächeninhalt $k^2a^2$
  \diagnose{Bei a) zeigen Ergebnisse wie $\overline{BB'} = 1.5$\,cm oder
  $\overline{SB'} = 2.5$\,cm, dass mit der Teilstrecke $\overline{AB}$ statt
  mit $\overline{SB}$ gerechnet wird. Additive Ansätze ($3 + 2 = 5$) zeigen
  fehlendes Verständnis der Ähnlichkeit. $ka^2$ bei c) zeigt, dass die Fläche
  wie die Länge skaliert wird.}
\end{solution}
\raum{4cm}

% ---------------------------------------------------------------------------
% 9 Sprache ↔ Mathematik
% ---------------------------------------------------------------------------
\aufgabe{9.6}
\begin{parts}
  \part Handyabo A kostet CHF 20 pro Monat plus CHF 0.10 pro Gesprächsminute.
        Abo B kostet pauschal CHF 35 pro Monat. Stellen Sie für beide Abos
        einen Term für die Monatskosten bei $x$ Minuten auf. Ab wie vielen
        Minuten ist Abo B günstiger?
  \part Für einen Schulausflug fahren $200$ Personen mit Bussen, die je
        $45$ Plätze haben. Wie viele Busse braucht es?
\end{parts}
\begin{solution}
  a) A: $20 + 0.1x$; B: $35$. Aus $20 + 0.1x = 35$ folgt $x = 150$. Abo B ist
  ab mehr als $150$ Minuten günstiger. \quad
  b) $200 \geteilt 45 \approx 4.4$, also braucht es $5$ Busse.
  \diagnose{Die Antwort «$150$» ohne Aussage, was darunter und darüber gilt,
  zeigt, dass die Lösung nicht im Kontext gedeutet wird. Bei b) sind $4.44$
  oder $4$ Busse die klassischen Zeichen dafür, dass ein Rechenergebnis nicht
  auf die Situation zurückbezogen wird.}
\end{solution}
\raum{4.5cm}


\end{questions}

% ===========================================================================
% Auswertungsraster (für die Lehrperson)
% ===========================================================================
\newpage
\section*{Auswertung \normalsize(wird von der Lehrperson ausgefüllt)}

Pro Aufgabe: \quad \textbf{2} = sicher gelöst \quad \textbf{1} = teilweise
gelöst \quad \textbf{0} = nicht gelöst. \\
Die \emph{Bruchstelle} ist das tiefste Niveau eines Strangs, auf dem die
Aufgaben nicht mehr sicher gelöst werden.

\begin{center}
\small
\renewcommand{\arraystretch}{2.2}
\newcommand{\feld}[1]{#1\,\framebox[0.6cm]{\rule{0pt}{0.9em}}}
\begin{tabular}{|>{\raggedright\arraybackslash}p{3.5cm}|c|c|c|c|p{1.9cm}|}
  \hline
  \textbf{Strang}
    & \textbf{A}\,\footnotesize Grundlagen & \textbf{B}\,\footnotesize Numerisch
    & \textbf{C}\,\footnotesize Symbol. & \textbf{D}\,\footnotesize Formal
    & \footnotesize\textbf{Bruchstelle} \\ \hline
  1 Zahlen und Stellenwert          & \feld{1.1}\ \feld{1.2} & \feld{1.3}\ \feld{1.4} & \feld{1.5} & \feld{1.6} & \\ \hline
  2 Operationen und Rechenreihenfolge & \feld{2.1}\ \feld{2.2} & \feld{2.3}\ \feld{2.4} & \feld{2.5} & \feld{2.6} & \\ \hline
  3 Brüche                          & \feld{3.1}\ \feld{3.2} & \feld{3.3}\ \feld{3.4} & \feld{3.5} & \feld{3.6} & \\ \hline
  4 Proportionalität und Prozent    & \feld{4.1}\ \feld{4.2} & \feld{4.3}\ \feld{4.4} & \feld{4.5} & \feld{4.6} & \\ \hline
  5 Terme                           & \feld{5.1}\ \feld{5.2} & \feld{5.3}\ \feld{5.4} & \feld{5.5} & \feld{5.6} & \\ \hline
  6 Gleichungen                     & \feld{6.1}\ \feld{6.2} & \feld{6.3}\ \feld{6.4} & \feld{6.5} & \feld{6.6} & \\ \hline
  7 Funktionen und Graphen          & \feld{7.1}\ \feld{7.2} & \feld{7.3}\ \feld{7.4} & \feld{7.5} & \feld{7.6} & \\ \hline
  8 Geometrie und Grössen           & \feld{8.1}\ \feld{8.2} & \feld{8.3}\ \feld{8.4} & \feld{8.5} & \feld{8.6} & \\ \hline
  9 Sprache $\leftrightarrow$ Mathematik & \feld{9.1}\ \feld{9.2} & \feld{9.3}\ \feld{9.4} & \feld{9.5} & \feld{9.6} & \\ \hline
\end{tabular}
\end{center}

\textbf{Bemerkungen:}

\vspace{4cm}

\end{document}
